\documentclass[11pt,a4paper]{article}
\usepackage[utf8]{inputenc}
\usepackage[T1]{fontenc}
\usepackage[spanish,es-nodecimaldot]{babel}
\usepackage[margin=2.2cm]{geometry}
\usepackage{amsmath,amssymb,array,longtable,booktabs}
\usepackage{xcolor}
\usepackage{hyperref}
\hypersetup{colorlinks=true,linkcolor=blue!45!black,urlcolor=blue!55!black}
\setlength{\parindent}{0pt}
\setlength{\parskip}{6pt}
\setlength{\emergencystretch}{3em}
\newcommand{\code}[1]{\texttt{\detokenize{#1}}}
\newcommand{\sat}{\operatorname{sat}}
\newcommand{\trunc}{\operatorname{trunc}}
\title{EXPLICACIÓN EFECTOS PROYECTO FFB YOKE}
\author{Análisis del código MiYoke\_PD\_V2.4}
\date{8 de octubre de 2026}
\begin{document}
\maketitle
\begin{abstract}
Este documento explica el centrado dinámico, las turbulencias, las vibraciones en suelo, los buffets de tren de aterrizaje, pérdida y spoiler, el damper y la fricción, para roll y pitch. Cada efecto tiene su sección: datos de entrada, fórmulas, procesamiento y salida PWM. Se analizan también la adquisición de posición por SimConnect, potenciómetros y encoders, y la relación entre las resistencias al movimiento y los demás efectos. Se distingue el código implementado de los fundamentos físicos, los antecedentes y las mejoras propuestas. Los valores numéricos son los predeterminados del código inspeccionado; una configuración guardada en la placa puede reemplazarlos. No se ha verificado qué firmware está cargado ni se han realizado mediciones de fuerza en el mecanismo.
\end{abstract}
\tableofcontents
\newpage

\section{Arquitectura vigente SC8: cálculo en Raspberry Pi Pico 2}
El programa de Windows solicita las 22 variables originales de SimConnect, selecciona la muestra más reciente y la transmite. No calcula presión normalizada, ráfagas, envolventes de buffet ni fuerzas de suelo. La Pico 2 ejecuta esas operaciones en \code{PicoForceModels.h}, los modelos \code{GroundModel.h}, \code{GustModel.h} y \code{BuffetModel.h}, y después aplica \code{SimEffects.h}, ganancias, sensores locales y límites de PWM. Los archivos del mismo nombre dentro de \code{SimConnectBridge} son inclusiones de compatibilidad para pruebas; no son modelos ejecutados por el puente.

El simulador continúa calculando la física del avión en la PC. La placa calcula toda la respuesta háptica del dispositivo a partir de esa telemetría. La selección del COM, la conexión con SimConnect, el permiso de armado y los registros siguen en Windows. Los parámetros de referencia y opciones del lanzador se transmiten como configuración; sus fórmulas se evalúan en la Pico.

\subsection{Contrato de transporte}
SC8 transmite 113 bytes, con enteros y flotantes en orden little endian. Los datos del simulador se solicitan como \code{FLOAT64} y se convierten a \code{float32}; se conservan sus unidades originales. Ya no se transmiten fuerzas ni envolventes cuantizadas con escala 10000.
\begin{longtable}{p{0.18\linewidth}p{0.74\linewidth}}
\toprule Bytes & Contenido\\\midrule\endhead
0--1 & Cabecera A5 5A.\\
2--3 & Versión 8 y banderas: bit 0 armado, bit 1 pausa; mutuamente excluyentes.\\
4--5 & Secuencia de 16 bits.\\
6--9 & Marca de captura de 32 bits, en milisegundos del reloj de Windows.\\
10--97 & 22 valores float32, en el orden definido en \code{RawTelemetry.h}.\\
98--109 & Referencia en kt, ganancia previa de ráfagas y mezcla longitudinal: tres float32.\\
110 & Opciones: bit 0 ráfagas, bit 1 buffets.\\
111--112 & CRC16 de los primeros 111 bytes.\\\bottomrule
\end{longtable}
El orden de los 22 valores es: presión en psf, IAS, vientos X/Y/Z, estado de suelo, slew, AoA, ángulo de pérdida, flaps izquierdo/derecho, VS0/VS1, extensión de tren, tren retráctil, spoilers izquierdo/derecho, velocidad de suelo, yugo roll/pitch y aceleraciones X/Y. Para un campo de índice $i$ (desde cero), el desplazamiento es $10+4i$.

El receptor conserva compatibilidad SC1--SC7, pero el lanzador y puente actuales exigen \code{YOKE17SC8}. Se deben actualizar juntos firmware y ejecutable. La identidad USB usada para detectar la Pico y los ajustes persistidos del menú se conservan.

\subsection{Relojes, lotes y corrección de suelo}
El registro del vuelo contenía 5670 filas en 1152 instantes de recepción. En SC7, 4518 de los 5669 intervalos eran cero y reiniciaban el filtro. Sólo ocho filas tenían señal de suelo distinta de cero. No era un problema que se solucionara aumentando la ganancia.

Ahora el puente guarda la última muestra de cada despacho y transmite como máximo cada 20 ms. La marca enviada corresponde a la recepción de esa muestra en Windows, no a la escritura USB ni a un reloj físico del simulador. La Pico calcula los modelos una vez por nueva muestra aceptada, con la diferencia entre marcas. El bucle de motores continúa nominalmente cada 2 ms, manteniendo las salidas de los modelos entre muestras; la textura de pista y las portadoras avanzan con el reloj local. Una marca repetida o que retrocede no renueva el watchdog; un desarme o pausa puede llegar con la misma marca. El modelo de suelo conserva su memoria ante $\Delta t=0$ y no inventa un intervalo positivo.

Esta selección evita procesar como muestras independientes los datos entregados juntos, pero puede omitir picos dentro del lote. No recupera la cronología original del simulador. La prueba con CSV reproduce los lotes por igualdad de marca en milisegundos; no reconstruye exactamente la planificación de Windows ni mide los motores.

\subsection{Notación y registros actuales}
En las fórmulas siguientes, los subíndices $\mathrm{rx}$ identifican la entrada disponible para la etapa local; en SC8 las envolventes, presión normalizada y fuerzas se calculan en la Pico y no se decodifican de enteros. Los ejemplos tabulados de generaciones anteriores se conservan redondeados para explicar las leyes; el cálculo actual no aplica su antiguo escalón USB de 0.0001. Las ganancias del menú sí siguen almacenadas en diezmilésimas.

\code{telemetry.csv} contiene 34 columnas de datos crudos, configuración y estado. \code{sent=1} identifica un envío USB completado; las filas de recepción tienen \code{sent=0} y \code{armed=0}, lo cual no significa por sí solo que los motores se hayan desarmado. Ya no incluye fuerzas calculadas ni indicadores de efectos. El monitor sin placa sólo registra entradas.

\code{pico.log} guarda el diagnóstico recibido desde la placa. La línea \code{SC8 t=... gr=... gp=... R=... P=... us=...} muestra marca de muestra, señales de suelo multiplicadas por 10000 antes de ganancias, PWM total aplicado a cada eje y máximo tiempo observado de \code{Effects::update} en microsegundos. Se intenta registrar cada 100 ms sin bloquear; puede omitir picos. R/P incluyen todos los efectos activos. El tiempo medido no incluye el resto del bucle ni demuestra por sí solo que se cumpla siempre el período de 2 ms.

\section{CENTRADO DINÁMICO (RESORTE)}

\subsection{Qué hace y qué significa que sea dinámico}
El efecto intenta devolver el yugo a su centro. Cuanto más se aparta el mando del centro, mayor es el comando restaurador. Además, cuando aumenta la presión dinámica del aire, el resorte se endurece. En este proyecto el centro permanece fijo: lo que cambia con la telemetría es la \textbf{rigidez}, no la posición de equilibrio.

La idea se expresa como un resorte lineal:
\begin{equation}
u_a=-G\,g_a\,k_f\,x_a,\qquad a\in\{R,P\}.
\label{eq:resorte}
\end{equation}
Aquí $R$ es roll y $P$ es pitch. $x_a$ es cuánto se desplazó el mando, $k_f$ es la rigidez variable suavizada, $g_a$ ajusta la intensidad de ese eje y $G$ es la ganancia general. El signo menos hace que el motor tienda a volver al centro. $u_a$ es un \textbf{comando adimensional}: todavía no es PWM, newtons ni torque medido.

\subsubsection{Tres cantidades diferentes que no hay que confundir}
\begin{itemize}
\item \textbf{Posición normalizada} $x$: $0$ es el centro; $-1$ y $+1$ son los extremos calibrados. Es una entrada al cálculo.
\item \textbf{Presión normalizada} $Q$: relaciona la presión del simulador con una presión de referencia. En el protocolo vigente puede llegar a $4$.
\item \textbf{Comando normalizado} $u$: indica sentido e intensidad de la salida. La interfaz del motor admite $[-1,1]$, pero este firmware limita antes la suma a aproximadamente $\pm0.862745$.
\end{itemize}
Por lo tanto, ``generar un valor de $\pm1$'' significa trabajar en una escala normalizada. No significa que siempre se produzcan exactamente $-1$ o $+1$, ni que la presión esté limitada a uno.

\subsection{Recorrido completo en la versión actual}
\begin{center}
\fbox{\begin{minipage}{0.91\linewidth}
\textbf{Rama de rigidez:} Prepar3D/SimConnect $\rightarrow$ presión en psf $\rightarrow$ USB SC8 $\rightarrow$ conversión local a Pa $\rightarrow$ división por $q_{\rm ref}$ $\rightarrow$ límite $0\ldots4$ $\rightarrow$ rigidez de reposo y filtro de 100 ms.

\textbf{Rama de posición actual:} potenciómetro roll/pitch $\rightarrow$ ADC de 12 bits $\rightarrow$ promedio e histéresis $\rightarrow$ calibración de extremos y centro $\rightarrow$ posición $-1\ldots1$.

\textbf{Unión:} $-\text{posición}\times\text{rigidez}\times\text{ganancia del eje}\times\text{ganancia total}$ $\rightarrow$ límite de comando $\rightarrow$ polaridad $\rightarrow$ multiplicación por 255 y redondeo $\rightarrow$ límite PWM 220 $\rightarrow$ puente H.
\end{minipage}}
\end{center}

\textbf{Hallazgo importante:} \code{mis_defines.h} define \code{YOKE_USE_SIMULATION=2}. Por defecto, la fuerza depende de la presión de SimConnect y de la posición de los \textbf{potenciómetros locales}. No depende de la posición de SimConnect, aunque el puente siga solicitándola, validándola y transmitiéndola. Las opciones son:
\begin{center}
\begin{tabular}{cll}
\toprule
Valor & Posición utilizada & Procedencia\\
\midrule
0 & Encoders & Cuentas relativas al centro al arrancar\\
1 & SimConnect & \code{YOKE X POSITION} / \code{YOKE Y POSITION}\\
2 & Potenciómetros & ADC0/ADC1; opción predeterminada\\
\bottomrule
\end{tabular}
\end{center}
Una definición de compilación puede cambiar esa opción. \code{Build-Pico2.ps1 -Encoders} fuerza el modo 0 y \code{-Potentiometers} fuerza el modo 2. Analizar el archivo fuente no permite certificar el modo del binario ya cargado. El firmware actual no transmite ejes HID al simulador: el código prevé que otra placa envíe los mandos.

\subsection{Variables de SimConnect que intervienen}
\subsubsection{Datos que entran directamente en la fórmula}
\begin{longtable}{p{0.31\linewidth}p{0.18\linewidth}p{0.43\linewidth}}
\toprule Variable & Unidad solicitada & Uso concreto\\\midrule
\endhead
\code{DYNAMIC PRESSURE} & \code{psf} & Presión dinámica común a roll y pitch. Determina la rigidez en los tres modos.\\
\code{YOKE X POSITION} & \code{Position} & Posición de roll, entre $-1$ y $+1$. Se usa en la fuerza solamente con modo 1.\\
\code{YOKE Y POSITION} & \code{Position} & Posición de pitch, entre $-1$ y $+1$. Se usa en la fuerza solamente con modo 1.\\
\bottomrule
\end{longtable}
El SDK describe X como desplazamiento izquierda/derecha e Y como mando hacia dentro/fuera. Son posiciones de mandos, no los ángulos de alabeo o cabeceo del avión. La unidad psf es libra-fuerza por pie cuadrado. Fuente: \href{https://www.prepar3d.com/SDKv3/LearningCenter/utilities/variables/simulation_variables.html}{SDK de Prepar3D, Simulation Variables}.

\subsubsection{Datos que habilitan o invalidan la salida}
\code{IS SLEW ACTIVE} debe valer cero. Los eventos de sistema \code{Sim} y \code{Pause} indican simulación en marcha y pausa; son eventos, no factores multiplicativos del resorte. También se necesitan \code{--enable}, un puerto identificado y muestras recientes.

El paquete de entrada es compartido con otros efectos. Su validación global exige además valores finitos para \code{AIRSPEED INDICATED}, \code{AIRCRAFT WIND X}, \code{AIRCRAFT WIND Y}, \code{AIRCRAFT WIND Z}, \code{SIM ON GROUND} y slew; el estado de suelo debe ser 0 o 1. Un valor inválido de estas entradas puede desarmar también el centrado, aunque no participe en su ecuación. La presión debe estar entre 0 y 5000 psf y ambas posiciones del simulador entre $-1$ y $1$, incluso en modo de potenciómetros.

\subsubsection{Datos que no calculan la rigidez}
La velocidad indicada no se eleva al cuadrado para calcular este resorte. El puente recibe directamente \code{DYNAMIC PRESSURE}. No intervienen como factores del centrado los flaps, tren, spoilers, ángulo de ataque, pérdida, velocidad sobre suelo, trim, RPM ni aceleraciones. Estar en tierra no anula el resorte. Si el simulador informa presión dinámica con el avión quieto, el modelo responde a esa presión; ``en reposo'' no es una condición equivalente a $q=0$.

\subsection{Paso 1: presión dinámica y referencia}
\subsubsection{Fundamento físico}
La presión dinámica se representa mediante
\begin{equation}
q=\frac12\rho V^2.
\end{equation}
$\rho$ es la densidad del aire en kg/m$^3$ y $V$ la velocidad relativa al aire en m/s; el resultado se expresa en Pa. A igual densidad, duplicar $V$ cuadruplica $q$. Esta relación física está explicada por \href{https://www1.grc.nasa.gov/beginners-guide-to-aeronautics/dynamic-pressure-2/}{NASA Glenn, Dynamic Pressure}. Justifica que el resorte se endurezca con la presión; no determina por sí sola las ganancias del yugo ni reproduce el momento de bisagra de un avión concreto.

\subsubsection{Conversión que hace la Pico}
En \code{PicoForceModels.h}, dentro de \code{update}:
\begin{align}
q_{\rm Pa}&=47.88025898\,q_{\rm psf},\\
V_{\rm ref}&=0.514444444\,V_{\rm ref,kt},\\
q_{\rm ref}&=\frac12(1.225)V_{\rm ref}^2,\\
Q&=\sat_{[0,4]}\left(\frac{q_{\rm Pa}}{q_{\rm ref}}\right).
\label{eq:q}
\end{align}
Definimos $\sat_{[a,b]}(z)=\min(b,\max(a,z))$: deja pasar el valor dentro del rango y lo recorta si se sale.

El factor 47.88025898 convierte psf a Pa; 0.514444444 convierte nudos a m/s, a partir de $1852/3600$. El valor 1.225 es una densidad estándar de referencia, fijada en el código; no es una lectura de la densidad actual. La presión real de entrada ya proviene del simulador.

\code{referenceKnots} vale 120 por defecto y puede modificarse con \code{--reference-knots}. El puente admite 40 a 250 kt. Con 120 kt:
\begin{equation}
V_{\rm ref}\simeq61.7333\ {\rm m/s},\qquad
q_{\rm ref}\simeq2334.24\ {\rm Pa}\simeq48.752\ {\rm psf}.
\end{equation}
$Q=1$ significa que la presión iguala esa referencia; $Q=0.5$, que es la mitad; $Q=4$, cuatro veces la referencia. \textbf{Los 120 kt son un ajuste del sistema, no la Vne del avión.} El tope de 4, introducido en SC6 y conservado en SC8, es una decisión de implementación, no una ley aerodinámica.

\subsection{Paso 2: transporte del dato a la placa}
\subsubsection{Codificación y recuperación}
La trama SC8 lleva presión original en psf y posición del yugo como float32. \code{Receiver::feed} comprueba CRC, secuencia, rangos y reloj; \code{PicoForceModels::update} convierte la presión y calcula $Q$. La presión ocupa bytes 10--13 y las posiciones roll/pitch 82--85 y 86--89. No se truncan a diezmilésimas. Véase el contrato SC8 al comienzo.

\subsubsection{Compatibilidad y tiempos de corte}
Las versiones 6 y 7 permiten $Q\le4$; las versiones 1--5 admitidas por el receptor limitan $Q\le1$. El receptor habilita disponibilidad de centrado desde v3 y posiciones de simulador desde v4. En modo 1 no hay reemplazo automático por sensores locales si falta la posición simulada.

El puente exige que la muestra seleccionada tenga menos de 250 ms para enviarla armada. La placa requiere trama válida y armada de menos de 300 ms y \code{Serial.dtr()} activo. Al arrancar el efecto, al perder habilitación o si el intervalo de actualización supera 50 ms, se reinicia la rigidez filtrada a cero. La desconexión o pausa retiran el resorte de reposo también: no existe un centrado autónomo permanente sin telemetría válida. El menú activo y el guardado de ajustes también fuerzan PWM cero.

\subsection{Paso 3: obtener la posición normalizada}
\subsubsection{Modo actual: potenciómetros}
Roll se lee en GP26/ADC0 y pitch en GP27/ADC1. \code{AxisPotentiometers.h} descarta la primera conversión al cambiar canal y promedia cuatro lecturas de 12 bits:
\begin{equation}
A=\left\lfloor\frac{a_1+a_2+a_3+a_4+2}{4}\right\rfloor.
\end{equation}
$a_i$ son las lecturas individuales de 0 a 4095; $A$ es la lectura promedio redondeada. Antes de normalizar, \code{PotentiometerPosition.h} utiliza una lectura retenida $H$ con banda móvil $h=4$ cuentas:
\begin{equation}
H_n=\begin{cases}
A_n-h,&A_n>H_{n-1}+h,\\
A_n+h,&A_n<H_{n-1}-h,\\
H_{n-1},&\text{en otro caso}.
\end{cases}
\end{equation}
Al inicializar se toma $H=A$. Esta histéresis reduce pequeñas fluctuaciones. No es una zona muerta fija del resorte: sigue desplazamientos lentos acumulados y puede dejar un pequeño error residual al volver al centro.

Para un sensor cuyo ADC aumenta hacia el lado positivo, la normalización es:
\begin{equation}
x=\sat_{[-1,1]}\begin{cases}
\displaystyle\frac{H-C}{P-C},&H\ge C,\\[4pt]
\displaystyle\frac{H-C}{C-N},&H<C.
\end{cases}
\label{eq:pot}
\end{equation}
$N$ es el ADC del extremo mecánico negativo; $C$, el centro; y $P$, el extremo positivo. Así, $H=N$ da $-1$, $H=C$ da 0 y $H=P$ da $+1$. El código admite también sensores invertidos, con $N>C>P$: selecciona el lado según el signo de $P-C$ y usa $(H-C)/(P-C)$ en el lado positivo y $-(H-C)/(N-C)$ en el negativo.

Los valores iniciales de ambos ejes son $N=0$, $C=2048$, $P=4095$: deben sustituirse por la calibración real. \code{YokeMotion} conserva la posición limitada y calcula también velocidad, pero el resorte no utiliza esa velocidad ni el filtro de velocidad.

\textbf{Detalle importante del menú:} la calibración actual captura cuatro extremos, primero pitch negativo/positivo y después roll negativo/positivo. Calcula el centro como
\begin{equation}
C=\left\lfloor\frac{N+P}{2}\right\rfloor.
\end{equation}
No captura un tercer punto central. Rechaza recorridos ADC menores que 64 cuentas. Si el centro mecánico no coincide con el promedio eléctrico de los extremos, conviene implementar una calibración de tres puntos. El normalizador ya admite ese centro independiente, pero la rutina del menú no lo mide.

\subsubsection{Modo 0: encoders}
Las cuentas se refieren al centro establecido al arrancar y se dividen por el recorrido de cada lado:
\begin{equation}
x=\sat_{[-1,1]}\begin{cases}
c/c_{\max},&c\ge0,\\
c/(-c_{\min}),&c<0.
\end{cases}
\end{equation}
$c$ son las cuentas con su dirección configurada. Los límites actuales son $\pm1000$ para roll y $\pm2546$ para pitch; los comentarios consideran el de pitch una estimación. Corresponde medirlos. Esta opción no está activa por defecto.

\subsubsection{Modo 1: posición de SimConnect}
Se usa directamente $x_{{\rm rx},R}$ o $x_{{\rm rx},P}$, limitada a $[-1,1]$. Una muestra cero anula inmediatamente el aporte de centrado al procesarse, aunque la rigidez tarde en cambiar. En este modo la respuesta de posición depende del simulador y de la comunicación; ejecutar el lazo de placa cada 2 ms no produce nuevas muestras de SimConnect cada 2 ms.

\subsection{Paso 4: rigidez mínima y suavizado}
\subsubsection{Rigidez objetivo}
En \code{SimEffects.h::Effects::update}:
\begin{equation}
k^*=r+(1-r)Q_{\rm rx},\qquad r=0.15\text{ por defecto}.
\label{eq:rigidez}
\end{equation}
$r$ es \code{CenterRest}, mostrado como ``Centro reposo''. $r$ agrega resistencia al desplazamiento incluso con presión cero. $(1-r)Q_{\rm rx}$ es el aporte que cambia con la presión. Ambos son adimensionales.
\begin{center}
\begin{tabular}{rrrrr}
\toprule
$Q_{\rm rx}$ & 0 & 0.25 & 1 & 4\\
$k^*$ con $r=0.15$ & 0.15 & 0.3625 & 1 & 3.55\\
\bottomrule
\end{tabular}
\end{center}
La descripción antigua ``15\% del máximo'' resulta imprecisa para SC8: es el 15\% de la rigidez a $Q=1$, no el 15\% del máximo $3.55$. $k^*$ no está limitado a uno. Con $r=0$ queda $k^*=Q$; con $r=1$ resulta un resorte de rigidez constante $k^*=1$.

Una consecuencia poco evidente:
\begin{equation}
\frac{\partial k^*}{\partial r}=1-Q.
\end{equation}
Subir ``Centro reposo'' endurece por debajo de $Q=1$, no cambia nada en $Q=1$ y \textbf{ablanda por encima de $Q=1$}. Esto se debe a que se reutiliza una interpolación fuera del intervalo $0\ldots1$; no es una calibración específica de un avión.

\subsubsection{Filtro de primer orden}
El código suaviza la rigidez, guardándola en una variable llamada \code{centerPressure}:
\begin{align}
\alpha_n&=\frac{\Delta t_{\rm ms}}{100+\Delta t_{\rm ms}},\\
k_{f,n}&=k_{f,n-1}+\alpha_n(k^*_n-k_{f,n-1}).
\label{eq:filtro}
\end{align}
Aunque el nombre diga ``pressure'', esta variable ya representa \textbf{rigidez normalizada filtrada}, no presión en Pa. $\Delta t$ es el tiempo real entre actualizaciones en milisegundos. Con 2 ms, $\alpha=2/102\simeq0.019608$: en cada ciclo se corrige aproximadamente el 1.96\% de la diferencia pendiente.

Esta forma corresponde a una discretización de $\tau\dot{k}_f+k_f=k^*$ con $\tau=0.1$ s. A entrada constante, partiendo de cero y con ciclos de 2 ms, $k_{f,n}=k^*[1-(100/102)^n]$. A los 100 ms alcanza aproximadamente 62.8\% del objetivo y a los 300 ms, 94.9\%. No es una espera fija de 100 ms: empieza a responder desde el primer ciclo posterior al reinicio.

El filtro evita cambios bruscos al variar presión, pero no retrasa deliberadamente la posición. La posición de potenciómetros sí tiene promedio e histéresis propios. El valor 100 ms es un ajuste local en \code{SimEffects.h}, no una constante física tomada de NASA ni una exigencia de TelemFFB.

\subsection{Paso 5: cálculo separado para roll y pitch}
Con las ganancias iniciales, después del filtro:
\begin{align}
u_R&=-x_R\,k_f\,(0.40)(0.90)=-0.36\,k_f\,x_R,\\
u_P&=-x_P\,k_f\,(0.80)(0.90)=-0.72\,k_f\,x_P.
\end{align}
\code{CenterRoll} y \code{CenterPitch} son independientes en \code{GainSettings.h}. El valor inicial de roll se define como la mitad del de pitch en \code{mis_defines.h}, pero al ajustar el menú la relación no se mantiene automáticamente. El cociente 2:1 solamente describe los valores iniciales, a igual posición normalizada, antes de saturar.

\subsubsection{Signo y centro}
Si $x>0$, $u<0$; si $x<0$, $u>0$; si $x=0$, $u=0$. No hay un término de trim ni un centro móvil $x_0$: equivale a usar $x_0=0$ en $u=-Ggk_f(x-x_0)$. Al mantener el mando desplazado, el centrado persiste aunque la velocidad del mando sea cero.

\subsubsection{Qué significa la pendiente}
Antes de saturar, la pendiente es $\partial u_a/\partial x_a=-Gg_ak_f$. Duplicar la posición duplica el comando; duplicar una ganancia también lo duplica. Esta es una ley proporcional de posición. No tiene integrador, derivada ni medición de torque en esta cadena. El damper es otro efecto y no forma parte de la ecuación del resorte.

\subsection{Paso 6: del comando al PWM}
\subsubsection{Límite de la salida de efectos}
\code{Effects::command} limita la suma de efectos. Para analizar solamente centrado, suponiendo los demás aportes nulos:
\begin{equation}
L=\frac{220}{255}\simeq0.862745,\qquad
u_{\rm out,a}=\sat_{[-L,L]}(u_a).
\end{equation}
Con otros efectos activos se limita la \emph{suma}, no el resorte por separado. Por eso el PWM observado en funcionamiento normal puede no coincidir con una tabla de centrado aislado. Aquí no se desarrollan las fórmulas de los otros efectos.

\subsubsection{Polaridad, escala y redondeo}
En el sketch se aplica $s_R=\code{ROLL_MOTOR_DIRECTION}$ y $s_P=\code{PITCH_MOTOR_DIRECTION}$, ambos inicialmente $+1$. \code{MotorDriver.h::setCommand} hace:
\begin{align}
z_a&=\sat_{[-1,1]}(s_a u_{\rm out,a}),\\
M&=\sat_{[0,255]}(\code{PWM_MAX})=255,\\
p_a^*&=\operatorname{round}_{\rm sim}(Mz_a),\\
p_a&=\sat_{[-220,220]}(p_a^*).
\label{eq:pwm}
\end{align}
El redondeo simétrico suma 0.5 si el valor es positivo o resta 0.5 si es negativo y convierte a entero truncando hacia cero. Por ejemplo, $45.9\to46$ y $-45.9\to-46$. Cero produce exactamente cero. $M$ es el factor de escala; 220 es el techo duro, aplicado después. \textbf{No se multiplica por 220: por defecto se multiplica por 255 y se limita a 220.}

Al reducir \code{PWM_MAX}, se reduce además la escala disponible. Por ejemplo, con $M=200$, el límite previo $L$ deja como máximo aproximadamente $\operatorname{round}(200L)=173$ cuentas, no 200. Ambos límites se acumulan.

\subsubsection{Pines del puente H y duty cycle}
Si $p\ge0$, la salida LEFT queda en 0 y RIGHT recibe $p$. Si $p<0$, RIGHT queda en 0 y LEFT recibe $|p|$. En Pico: roll GP6/GP7 y enable GP10; pitch GP8/GP9 y enable GP11. La dirección mecánica debe verificarse porque depende del cableado y de la transmisión.

El setup configura \code{analogWriteRange(255)} y \code{analogWriteFreq(490)}. El duty cycle comandado en la salida activa es
\begin{equation}
D=\frac{|p|}{255};\qquad |p|=220\Longrightarrow D\simeq86.27\%.
\end{equation}
Los 490 Hz son la frecuencia PWM configurada; los 500 Hz son la frecuencia nominal del lazo de cálculo y los 50 Hz el máximo aproximado de envío de telemetría. Son tres tiempos diferentes.

PWM indica cuánto tiempo se activa una salida. No certifica corriente ni fuerza: motor, alimentación, velocidad, resistencia, pérdidas y transmisión intervienen en la fuerza real. El límite 220 no es un limitador de corriente y la conversión mostrada no constituye un lazo de control de corriente.

\subsection{Ejemplos numéricos completos}
\subsubsection{Media carrera a presión de referencia}
Supongamos $q_{\rm Pa}=q_{\rm ref}$, telemetría válida, rigidez estabilizada y $x_R=x_P=+0.5$. Entonces $Q=1$, $k_f=0.15+0.85=1$ y:
\begin{align*}
u_R&=-(0.5)(1)(0.40)(0.90)=-0.18, &p_R&=\operatorname{round}(-45.9)=-46,\\
u_P&=-(0.5)(1)(0.80)(0.90)=-0.36, &p_P&=\operatorname{round}(-91.8)=-92.
\end{align*}
Con polaridades $+1$, roll activa LEFT con 46 y pitch LEFT con 92; las otras salidas quedan en cero. Los duty cycles son aproximadamente 18.04\% y 36.08\%. Para posición $-0.5$ se invierten los signos y las salidas activas.

\subsubsection{Tabla a media carrera}
Los valores siguientes usan ganancias predeterminadas, $x=+0.5$, $M=255$, polaridad $+1$ y filtro estabilizado. Se omite el pequeño error de cuantización de la presión.
\begin{center}
\begin{tabular}{rrrrrr}
\toprule
$Q$ & $k_f$ & $u_R$ bruto & PWM roll & $u_P$ bruto & PWM pitch\\
\midrule
0 & 0.1500 & $-0.02700$ & $-7$ & $-0.05400$ & $-14$\\
0.25 & 0.3625 & $-0.06525$ & $-17$ & $-0.13050$ & $-33$\\
1 & 1.0000 & $-0.18000$ & $-46$ & $-0.36000$ & $-92$\\
2 & 1.8500 & $-0.33300$ & $-85$ & $-0.66600$ & $-170$\\
4 & 3.5500 & $-0.63900$ & $-163$ & $-1.27800$ & $-220$\\
\bottomrule
\end{tabular}
\end{center}
La última fila muestra por qué hay que distinguir el cálculo bruto de la salida: pitch pide $-1.278$, pero el efecto entrega $-L$ y el driver aplica $-220$.

\subsubsection{Ejemplo de potenciómetro y arranque del filtro}
Con calibración inicial $N=0$, $C=2048$, $P=4095$, una lectura retenida $H=3072$ representa $x=1024/2047\simeq0.50024$. Si $Q=1$, una vez estabilizado resulta aproximadamente $-46$ PWM roll y $-92$ PWM pitch. $H$ es la lectura \emph{después} de la histéresis, no necesariamente el ADC bruto que se imprime.

Al habilitar desde $k_f=0$, el primer update de arranque conserva cero. En el siguiente update de 2 ms, con $k^*=1$, $k_f=0.019608$. Para $x=0.5$ se obtienen aproximadamente $u_R=-0.003529$ y $u_P=-0.007059$, que redondean a $-1$ y $-2$ PWM. La salida crece luego hacia el valor de régimen.

\subsubsection{Dónde empieza la saturación}
El límite previo aparece cuando $Gg_ak_f|x_a|\ge L$. Para $Gg_ak_f>0$:
\begin{equation}
|x_a|_{\rm sat}=\frac{L}{Gg_ak_f}.
\end{equation}
Si el resultado supera uno, el recorrido calibrado no alcanza la saturación. Con $Q=4$ y ganancias iniciales, roll satura cerca de $|x_R|=0.675$ y pitch de $|x_P|=0.338$. Desde allí, alejar más el mando ya no aumenta el PWM del resorte aislado. El redondeo puede producir la última cuenta algo antes de ese umbral ideal.

En el extremo $|x|=1$, pitch empieza a saturar aproximadamente en $Q=1.233$ y roll en $Q=2.643$. Con la densidad de referencia, equivalen a velocidades de aproximadamente 133 y 195 kt mediante $V=120\sqrt Q$. Son equivalencias del modelo, no límites de vuelo ni una conversión exacta de IAS real.

\subsection{De dónde procede cada parte}
\subsubsection{Fundamento físico y adaptación de TelemFFB}
La ley restauradora lineal es el principio general de un resorte. La dependencia de presión se apoya en $q=\rho V^2/2$. El comentario de \code{SimEffects.h} identifica explícitamente el principio de resorte aerodinámico de TelemFFB. La documentación local histórica también lo describe como adaptación.

Se inspeccionó la copia local de
\begin{quote}\small
\code{VPforce-TelemFFB-refactor/telemffb/sim/msfs_xp/}\par
\code{MsfsXpFlightControlsMixIn.py}.
\end{quote}
En \code{_calculate_dynamic_pressures}, TelemFFB multiplica la presión por una escala interna $s=0.005$. En \code{_calculate_vne_and_gains} calcula $Q_{\rm gain}=1/(q_{\rm Vne}s)$ y lo asigna a las ganancias de controles. En \code{_calculate_control_coefficients}, el coeficiente de alerón incluye $(q s)Q_{\rm gain}$, además de la corrección de resbalamiento. Si se omite esa corrección:
\begin{equation}
(q s)\frac{1}{q_{\rm Vne}s}=\frac{q}{q_{\rm Vne}}.
\end{equation}
Esa es la correspondencia conceptual con $q/q_{\rm ref}$ de MiYoke. La escala 0.005 se cancela en esa simplificación y \textbf{no es una ganancia que deba agregarse al código actual}.

TelemFFB puede usar Vne configurada o calcular una referencia a partir de datos de diseño; también mezcla presión del flujo de hélice en elevador, permite curvas, ganancias base y desplazamientos de centro. \code{_calculate_joystick_spring_coefficients} ajusta y limita coeficientes, y \code{_update_joystick_controls} los entrega al dispositivo FFB. MiYoke hace una adaptación mucho más simple: una presión común, centro fijo y cálculo local de fuerza/PWM. No reproduce toda la cadena TelemFFB, ni sus constantes calibran automáticamente este yugo.

Referencia pública para localizar el antecedente: \href{https://github.com/walmis/VPforce-TelemFFB/blob/refactor/telemffb/sim/msfs_xp/MsfsXpFlightControlsMixIn.py}{TelemFFB, archivo de controles MSFS/X-Plane}. Las afirmaciones de implementación anteriores se basan en la copia local inspeccionada; la rama pública puede cambiar.

\subsubsection{Qué se conserva del proyecto alemán original}
Se inspeccionó la carpeta vecina \code{../../PROEYCTO ORIGINAL/Arduino_FFB_Yoke}. En \code{src/Joystick.cpp}, el resorte HID usa \code{ConditionForceCalculator}, con posición normalizada, coeficientes, centro, banda muerta y saturaciones recibidas en el efecto. También existe una rama de resorte automático predeterminado:
\begin{equation}
F_{\rm original}=\operatorname{int16}\left[-10000\,
\sat_{[-1,1]}\left(\frac{\text{posición}}{\text{máxima posición}}\right)
\frac{\text{defaultSpringGain}}{100}\right].
\end{equation}
Es una rama condicional del original, no prueba de que estuviera siempre activada. Comparte la idea de oponer fuerza al desplazamiento, pero esta expresión no utiliza la presión de SimConnect. No corresponde atribuirle el cociente $q/q_{\rm ref}$ del firmware actual.

\code{AxisCalibration.cpp::MoveMotor} del original activa uno u otro PWM según el signo. Ese esquema coincide con el empleado por \code{MotorDriver.h}, cuyo comentario declara ese antecedente. Hay además otra función \code{MoveMotor} en \code{Arduino.ino}: utiliza \code{map}, PWM mínimo/máximo y un umbral de fuerza, con una convención de pines diferente. Por ello, ``proviene del original'' describe la arquitectura de dos salidas y la idea de conversión, no identidad exacta de todos los signos, umbrales y ecuaciones.

\subsubsection{Qué es implementación o ajuste propio de MiYoke}
\begin{longtable}{p{0.43\linewidth}p{0.49\linewidth}}
\toprule Parte & Procedencia comprobable\\\midrule\endhead
Nombres y unidades de SimConnect & SDK de Prepar3D y definición del puente.\\
$q=\rho V^2/2$ & Fundamento físico general; referencia NASA.\\
Resorte dependiente de presión normalizada & Adaptación conceptual documentada de TelemFFB.\\
120 kt, $r=0.15$, $g_R=0.40$, $g_P=0.80$, $G=0.90$ & Valores locales de diseño. No se halló una medición de avión que los determine.\\
$r+(1-r)Q$, filtro de 100 ms y límite $Q=4$ & Implementación local observable; no copia exacta de la ley completa de TelemFFB.\\
Paquete SC8, watchdog, calibración ADC y ganancias persistidas & Infraestructura local de MiYoke.\\
Dos PWM según signo & Antecedente verificable en la calibración del proyecto original.\\
Escala $\times255$, redondeo simétrico y techo 220 & Implementación vigente de MiYoke, no conversión a torque físico.\\
\bottomrule
\end{longtable}
La atribución es una comparación de fuentes y comentarios, no una reconstrucción de autoría mediante historial Git: la carpeta principal inspeccionada no es un repositorio Git.

\subsection{Qué ganancias y calibraciones se pueden ajustar}
\subsubsection{Ajustes disponibles en el menú}
\begin{longtable}{p{0.24\linewidth}p{0.15\linewidth}p{0.53\linewidth}}
\toprule Ajuste & Inicio & Efecto sobre el resorte\\\midrule\endhead
Ganancia total, $G$ & 0.90 & Multiplica ambos ejes y los demás efectos; no sirve para balancear roll contra pitch.\\
Centrado roll, $g_R$ & 0.40 & Modifica sólo la intensidad de centrado roll.\\
Centrado pitch, $g_P$ & 0.80 & Modifica sólo la intensidad de centrado pitch.\\
Centro reposo, $r$ & 0.15 & Cambia la base y pendiente con presión según $r+(1-r)Q$. Es común a ambos ejes.\\
\bottomrule
\end{longtable}
Estos cuatro ajustes admiten valores de 0 a 1; los botones avanzan 0.05, es decir, cinco puntos porcentuales. No es un incremento relativo del 5\% sobre el valor anterior. El menú permite apagar centrado roll y pitch por separado. Apagar ``Centro reposo'' hace efectivo $r=0$: elimina la base, pero conserva el resorte aerodinámico si $Q>0$.

Las ganancias se guardan en memoria persistente y se cargan al arrancar. Cambiar un \code{define} no garantiza que se use el nuevo valor si hay ajustes guardados válidos. \code{YOKE_SETTINGS_VERSION} permite invalidar el registro de ganancias; el registro de potenciómetros tiene formato y versión separados. Hay que registrar los valores efectivamente cargados al comparar ensayos.

\subsubsection{Ajustes en el puente y en el firmware}
\begin{itemize}
\item \code{--reference-knots}: cambia la sensibilidad aerodinámica común. A presión fija, subir de 120 a 150 kt multiplica $Q$ por $(120/150)^2=0.64$, mientras no intervenga el recorte. La base $r$ no se multiplica por 0.64. La referencia es compartida con otros cálculos de la Pico.
\item Calibración ADC: determina centro, rango y sentido de posición. Debe corregirse antes de intentar resolver asimetrías con ganancias.
\item \code{ROLL_MOTOR_DIRECTION} y \code{PITCH_MOTOR_DIRECTION}: invierten el comando del motor. No corrigen un centro o un recorrido mal calibrado.
\item \code{YOKE_POT_HYSTERESIS_ADC} y \code{YOKE_POT_SAMPLES}: reducen ruido de entrada a costa de seguimiento y resolución efectiva. No endurecen directamente el resorte.
\item El número 100 en \code{SimEffects.h}: constante del filtro de rigidez en ms. Menor valor sigue más rápido los cambios de presión; mayor valor los suaviza más. No aumenta la tasa de SimConnect.
\item \code{SIMCONNECT_EXTRA_LIMIT}, \code{PWM_MAX} y \code{MOTOR_PWM_HARD_LIMIT}: fijan límites o escala; no corrigen la pendiente perdida cuando ya hay saturación.
\end{itemize}

\subsubsection{Cómo elegir ganancias sin saturar demasiado pronto}
Para mantener el resorte aislado dentro del límite hasta una posición $x_{\max}$ y una presión $Q_{\max}$ previstas, en régimen:
\begin{equation}
g_a\le\frac{L}{G\,[r+(1-r)Q_{\max}]\,|x_{\max}|}.
\end{equation}
Con $G=0.90$, $r=0.15$, $Q_{\max}=4$ y $|x_{\max}|=1$, se necesita $g_a\lesssim0.270$. Las ganancias iniciales 0.40 y 0.80 superan ese valor; por eso se recorta el recorrido a alta presión. Esto no obliga a usar 0.270: muestra la relación entre la dureza elegida y el margen antes de saturar. Si se necesita margen para otros efectos, el límite asignado al resorte debe ser menor que $L$.

\subsection{Procedimiento práctico para pulir el centrado}
\subsubsection{Primero, referencia mecánica y posición}
\begin{enumerate}
\item Confirmar qué modo de posición se compiló y qué ajustes están guardados. En potenciómetros, comprobar ADC negativo, centro y positivo por separado en cada eje.
\item Verificar que $x$ recorra aproximadamente $-1\to0\to+1$ en los sentidos físicos elegidos. La convención propuesta por el código es roll derecha positivo y pitch tirar positivo.
\item Comprobar el centro físico con el ADC retenido y no sólo con la animación del simulador. Si queda desplazado, corregir la calibración; bajar ganancia sólo oculta parte del error.
\end{enumerate}
En un sensor de pendiente creciente, si se cambia el centro configurado en $\delta C$ cuentas y el recorrido es aproximadamente $S$, cerca del centro el error de posición es $\delta x\simeq-\delta C/S$. Antes de saturar, introduce un error de comando de magnitud aproximada $Gg_ak_f|\delta C|/S$: el descentrado se nota más cuando aumenta la presión.

\subsubsection{Después, aislar el efecto y comprobar el signo}
\begin{enumerate}
\item Con simulación pausada de forma reconocida por el puente, o comunicación desconectada, usar el menú para dejar activos solamente los centrados que se ensayarán. Apagar los demás aportes, incluidos damper y fricción. El firmware bloquea la edición mientras la telemetría indica actividad sin pausa.
\item Empezar con ganancias pequeñas y un eje por vez. Desplazar ligeramente a ambos lados: el motor debe oponerse al desplazamiento. Si empuja a alejarse, detener la prueba y corregir coherencia de signos antes de subir ganancia.
\item Mantener posiciones fijas para separar resorte de cargas que dependen del movimiento. Probar centro, cuarto, media carrera y una posición mayor en ambos sentidos.
\end{enumerate}
El argumento \code{--no-buffets}, por sí solo, no aísla centrado: no apaga todas las otras contribuciones. No se ha modificado el firmware para este análisis.

\subsubsection{Finalmente, ajustar reposo, vuelo y margen}
\begin{enumerate}
\item Con presión baja conocida, ajustar $r$ para obtener el centrado residual deseado. No interpretar automáticamente avión detenido como presión cero.
\item Cerca de $Q=1$, ajustar $g_R$ y $g_P$ por separado: en ese punto $k^*=1$ independientemente de $r$, lo que facilita comparar ejes.
\item Repetir con varias presiones observadas en el CSV y esperar al menos unos 0.5 s tras cada cambio para comparar con el régimen calculado. Si la salida llega a 220 pronto, bajar ganancia de eje o revisar la referencia, en vez de atribuirlo a una falta de posición.
\item Medir fuerza en puntos del recorrido si se pretende reproducir una aeronave. La relación de comandos 2:1 no asegura una relación de fuerzas 2:1 entre una rotación de roll y una traslación de pitch.
\end{enumerate}

\subsubsection{Qué datos registrar}
El CSV registra \code{pressure_psf} y posiciones del simulador; $q_{\rm Pa}$ y $Q$ se calculan en la Pico. En modo 2 esas posiciones no son las posiciones locales que cierran el resorte. El diagnóstico local registra ADC, posición, comando y PWM; \code{pico.log} permite guardar esa salida.

Registrar por ensayo $Q$, posición local, ganancias efectivas, PWM y fuerza medida. La columna histórica \code{Centrado_Disponible} ya no se genera en SC8. Un diagnóstico futuro de $k_f$ y del aporte aislado permitiría observar los transitorios con más detalle.

\subsection{Mejoras posibles, todavía no implementadas}
\subsubsection{Centro calibrado y trim}
La mejora más directa es capturar un centro físico independiente de los extremos en el menú. Otra ampliación es incorporar un centro de resorte desplazable $x_0$ para trim, usando $u=-Ggk_f(x-x_0)$. Eso requiere decidir cómo relacionar el trim del avión con el mecanismo; no basta con sumar sin calibrar una variable angular a una posición normalizada.

\subsubsection{Separar rigidez de reposo y crecimiento aerodinámico}
Si no se desea que subir reposo ablande a $Q>1$, puede estudiarse una ley $k=k_0+k_qQ$ con parámetros independientes, o una curva acotada elegida explícitamente. Son propuestas distintas de la fórmula actual. Cambiarlas exige recalcular ganancias, saturación y transición al activar.

\subsubsection{Medición física y comportamiento cerca del centro}
Conviene medir fuerza frente a PWM para varios sentidos y posiciones, y considerar una regulación de corriente si se requiere control de torque. La forma del resorte en comando puede ser lineal y la fuerza real no serlo. No hay compensación PWM mínima en \code{setCommand}; órdenes pequeñas pueden no vencer la fricción mecánica, aunque el cálculo sea correcto.

Por redondeo, el primer PWM distinto de cero aparece aproximadamente cuando $|u|\ge0.5/255$. Antes de saturar, esto implica $|x|\gtrsim0.5/(255Gg_ak_f)$. A $Q=0$ con los valores iniciales resulta aproximadamente 3.63\% del semirrecorrido en roll y 1.82\% en pitch. Es cuantización eléctrica, a la que pueden añadirse histéresis y fricción; no es una banda muerta programada expresamente en la ley del resorte.

\subsection{Mapa de fuentes y alcance de la comprobación}
\subsubsection{Archivos que permiten reconstruir el cálculo}
Las rutas siguientes son relativas a la carpeta MiYoke\_PD\_V2.4, salvo el original indicado explícitamente:
\begin{longtable}{p{0.39\linewidth}p{0.53\linewidth}}
\toprule Archivo & Elementos a consultar\\\midrule\endhead
\code{SimConnectBridge/main.cpp} & \code{Receive}: datos crudos; \code{Send}: SC8; \code{main}: argumentos. $Q$ se calcula en \code{PicoForceModels.h}.\\
\code{SimTelemetry.h} & \code{Receiver::feed}, \code{fresh}: decodificación, rangos y vigencia.\\
\code{SimEffects.h} & \code{Effects::update}: $r$, rigidez y filtro; \code{commandWithPositionSource}: elección de posición; \code{command}: resorte y saturación.\\
\code{mis_defines.h} & Fuente predeterminada, ganancias, pines, tiempos y límites.\\
\code{GainSettings.h} & Valores iniciales, ajustes efectivos, pasos, encendido/apagado y persistencia.\\
\code{AxisPotentiometers.h} & Muestreo y promedio ADC.\\
\code{PotentiometerPosition.h} & Histéresis y normalización con tres valores.\\
\code{PotCalibration.h} & Centro por promedio de extremos y formato persistente.\\
\code{CalibrationMenu.h} & Captura de cuatro extremos, edición y confirmación de ajustes.\\
\code{CalibrationControls.h} & Carga y guardado de configuración.\\
\code{YokeMotion.h} & Posición limitada sin filtro temporal adicional de posición.\\
\code{MiYoke_PD_V2.4.ino} & Lazo de 2 ms, fuente de posición, aplicación de ganancias, polaridad, driver y diagnóstico.\\
\code{MotorDriver.h} & Escala, redondeo, techo absoluto y dos salidas PWM.\\
\bottomrule
\end{longtable}

\subsubsection{Precaución con los documentos históricos}
\code{LEEME md/CENTRADO_Y_RODAJE.md} se identifica como antecedente de versiones anteriores. Contiene rangos, ganancias y referencias a modos que no describen íntegramente el código actual. Para este documento se priorizaron el código vigente y la comparación directa con las copias de TelemFFB y del original. En particular, ahora se observan SC8, presión hasta 4, ganancias iniciales roll/pitch 0.40/0.80, límite $220/255$ y posición por potenciómetros.

\subsubsection{Referencias externas utilizadas}
\begin{enumerate}
\item NASA Glenn: \href{https://www1.grc.nasa.gov/beginners-guide-to-aeronautics/dynamic-pressure-2/}{Dynamic Pressure}. Fundamento de presión dinámica; consultado el 7 de octubre de 2026.
\item Lockheed Martin: \href{https://www.prepar3d.com/SDKv3/LearningCenter/utilities/variables/simulation_variables.html}{Prepar3D SDK v3, Simulation Variables}. Nombres y unidades de variables; consultado el 7 de octubre de 2026.
\item VPforce/TelemFFB: \href{https://github.com/walmis/VPforce-TelemFFB/blob/refactor/telemffb/sim/msfs_xp/MsfsXpFlightControlsMixIn.py}{archivo de controles en la rama refactor}. Referencia de localización pública; comparación realizada contra la copia local incluida en este proyecto.
\item Proyecto original local: \code{../../PROEYCTO ORIGINAL/Arduino_FFB_Yoke}, especialmente \code{src/Joystick.cpp}, \code{joystick.ino}, \code{AxisCalibration.cpp} y \code{Arduino.ino}. No se atribuye una URL de origen no comprobada.
\end{enumerate}

\textbf{Alcance:} esta explicación reconstruye el software inspeccionado y sus consecuencias matemáticas. Los ejemplos suponen ganancias iniciales y centrado aislado. La validación de fuerza real, signos del mecanismo, firmware cargado y parámetros persistidos requiere comprobar el montaje; no se presenta como realizada.
\newpage
\section{TURBULENCIAS}
\label{sec:turbulencias}

\subsection{Qué hace este efecto}
El proyecto convierte las \textbf{variaciones del viento del simulador} en empujes variables de roll y pitch. Separa el viento sostenido de sus cambios, combina las tres componentes del viento, pondera el resultado por presión dinámica y lo suaviza antes de enviarlo a la placa. La placa elimina señales muy pequeñas, aplica ganancias y obtiene el PWM.

Una forma sencilla de entenderlo es: \emph{``detectar cuánto cambió el viento respecto de su valor reciente y transformar ese cambio en un empuje del mando''}. No hay una tabla de posiciones a las que deba ir el yugo. El resultado es una orden de fuerza normalizada, no una posición objetivo.

\subsubsection{Diferencia con el centrado}
El centrado necesita el desplazamiento del mando para calcular su fuerza. La turbulencia aislada no multiplica por posición ni velocidad del yugo: puede actuar con el mando centrado y quieto. Potenciómetros, encoders y posición de SimConnect no determinan su amplitud en la fórmula. En la integración general sí se validan datos y pueden coexistir cargas de otros efectos; eso no convierte la turbulencia en un resorte.

\subsubsection{Qué representa y qué no representa}
Se trata de una \textbf{aproximación háptica basada en telemetría}. No calcula fuerzas de bisagra, no utiliza una geometría aerodinámica del avión ni devuelve newtons. Tampoco genera ruido aleatorio, una senoide fija o un campo atmosférico Dryden/Von Kármán. Si las componentes de viento permanecen constantes en los ejes del avión, la perturbación termina desapareciendo.

El nombre ``turbulencia'' debe interpretarse con ese alcance: el algoritmo detecta cambios de viento en coordenadas del avión. Una maniobra puede modificar esas componentes aun con viento atmosférico constante, porque cambian los ejes con los que se mide. El algoritmo actual no separa inequívocamente ráfaga atmosférica y cambio de orientación.

\subsection{Recorrido completo de roll y pitch}
\begin{center}
\fbox{\begin{minipage}{0.91\linewidth}
\textbf{En Windows:} SimConnect entrega viento X/Y/Z y presión; el puente selecciona la muestra más reciente, añade marca y configuración y transmite por USB SC8.

\textbf{En la Pico:} seguimiento del viento lento, extracción de variaciones, mezcla, escala de 5 kt, presión, ganancia, límite vectorial y filtro de 80 ms; luego zona muerta, ganancias de cada eje y PWM.

\textbf{En la placa:} zona muerta continua del 2\% $\rightarrow$ ganancia de turbulencia del eje $\rightarrow$ ganancia general $\rightarrow$ suma de efectos y límite $220/255$ $\rightarrow$ polaridad del motor $\rightarrow$ escala 255, redondeo y techo PWM 220 $\rightarrow$ puente H.
\end{minipage}}
\end{center}
El cálculo de viento y sus filtros se ejecutan en la Pico 2, en \code{GustModel.h}, invocado por \code{PicoForceModels.h}. \code{SimEffects.h} usa directamente sus salidas locales.

\subsection{Variables de SimConnect utilizadas}
\subsubsection{Entradas de la fórmula}
\begin{longtable}{p{0.30\linewidth}p{0.13\linewidth}p{0.49\linewidth}}
\toprule Variable & Unidad & Papel en el efecto\\\midrule\endhead
\code{AIRCRAFT WIND X} & knots & Componente lateral $w_X$; sus variaciones alimentan roll.\\
\code{AIRCRAFT WIND Y} & knots & Componente vertical $w_Y$; sus variaciones son el aporte principal a pitch.\\
\code{AIRCRAFT WIND Z} & knots & Componente longitudinal $w_Z$; sus variaciones agregan un aporte regulable a pitch.\\
\code{DYNAMIC PRESSURE} & psf & Ponderación de la perturbación según presión dinámica.\\
\bottomrule
\end{longtable}
Los nombres, unidades y ejes se verifican en el \href{https://www.prepar3d.com/SDKv3/LearningCenter/utilities/variables/simulation_variables.html}{SDK de Prepar3D, Simulation Variables}. El código solicita los datos como \code{FLOAT64}. Las componentes ya se reciben en ejes del avión: MiYoke no aplica una matriz de rotación adicional.

La tabla del SDK identifica los ejes, pero no basta por sí sola para certificar el sentido de fuerza que debe sentir el piloto. En el código actual una variación X positiva produce comando roll positivo; una Y positiva produce pitch positivo. Para convertir ese signo en movimiento físico hay que considerar el cableado y la polaridad del motor.

\subsubsection{Estados necesarios}
\code{SIM ON GROUND} debe ser cero: \textbf{el modelo de turbulencias se desactiva y reinicia en tierra}. \code{IS SLEW ACTIVE} debe ser cero, el evento \code{Sim} debe indicar actividad y \code{Pause} ausencia de pausa. También se exige continuidad de muestras y la opción \code{--gusts}. Estos estados no multiplican gradualmente la fuerza: habilitan o reinician el modelo.

\subsubsection{Entradas compartidas que no forman la turbulencia}
La presión usada es la que informa el simulador. \code{AIRSPEED INDICATED} no se utiliza directamente en la ecuación de turbulencia. Tampoco se utilizan trim, actitud, aceleraciones, RPM, flaps, tren, spoilers ni posiciones del yugo para calcular los dos empujes de viento.

Hay una dependencia indirecta de la validación global del paquete: una posición \code{YOKE X/Y POSITION} no finita o fuera de $[-1,1]$, una IAS no finita o un estado de suelo inválido pueden desarmar la salida completa, aunque esos datos no formen parte de la mezcla de viento. Conviene diferenciar \emph{variable que calcula la intensidad} de \emph{variable cuya invalidez impide ejecutar el conjunto}.

\subsection{Símbolos para seguir las fórmulas}
\begin{longtable}{p{0.18\linewidth}p{0.74\linewidth}}
\toprule Símbolo & Significado y unidad\\\midrule\endhead
$w_i$ & Viento recibido en el eje $i=X,Y,Z$, en kt.\\
$m_i$ & Estimación del viento lento o sostenido, en kt.\\
$h_i=w_i-m_i$ & Variación respecto de esa estimación, en kt.\\
$\Delta t$ & Diferencia entre marcas de captura procesadas por la Pico, en segundos.\\
$\tau_H=1$ s & Constante temporal del seguimiento del viento lento.\\
$W=5$ kt & Escala de normalización de las variaciones de viento.\\
$k_Z=0.25$ & Peso longitudinal hacia pitch; regulable entre $-1$ y $+1$.\\
$Q_T$ & Presión normalizada usada por turbulencia, limitada a $[0,1]$.\\
$\gamma=1$ & Ganancia configurada en el puente y aplicada en la Pico, \code{--gust-gain}; admite $0\ldots10$.\\
$(v_R,v_P)$ & Objetivo bruto antes de saturación vectorial; adimensional.\\
$(t_R,t_P)$ & Objetivo después de saturación vectorial; norma no mayor que 1.\\
$(o_R,o_P)$ & Salida suavizada del modelo de la Pico.\\
$z_R,z_P$ & Componentes normalizadas calculadas localmente en la Pico.\\
$c_R,c_P$ & Ganancias de turbulencia de la placa: 0.35 y 0.70 inicialmente.\\
$G=0.90$ & Ganancia general de la placa; compartida con los demás efectos.\\
$u_R,u_P$ & Comandos de turbulencia tras zona muerta y ganancias.\\
\bottomrule
\end{longtable}
Se usan nombres distintos para evitar confundir la ganancia previa $\gamma$, las ganancias por eje $c_a$ y la ganancia general $G$. Actúan en puntos distintos de la cadena y no siempre producen resultados intercambiables.

\subsection{Paso 1: extraer las variaciones del viento}
\subsubsection{Promedio temporal y resta}
Para cada muestra y cada componente, \code{Model::update} calcula:
\begin{align}
a_n&=e^{-\Delta t_n/\tau_H},\qquad \tau_H=1\ {\rm s},\\
m_{i,n}&=m_{i,n-1}+(1-a_n)(w_{i,n}-m_{i,n-1}),\\
h_{i,n}&=w_{i,n}-m_{i,n}.
\label{eq:turb-hp}
\end{align}
$m$ sigue al viento con cierta lentitud. Al restarlo del viento actual, queda la parte transitoria $h$. \textbf{La resta usa la media ya actualizada} $m_{i,n}$, no la del ciclo anterior: ese detalle importa para reproducir los números.

Con $\Delta t=0.02$ s, $a=e^{-0.02}\simeq0.980199$ y $1-a\simeq0.019801$. En cada muestra la media avanza aproximadamente el 1.98\% de la diferencia con el viento recibido. Una ráfaga brusca no es absorbida inmediatamente por esa media y aparece en $h$.

\subsubsection{De dónde sale el filtro}
El seguimiento de media responde a la ecuación de primer orden
\begin{equation}
\tau_H\frac{dm}{dt}+m=w.
\end{equation}
Al resolverla en un intervalo con entrada constante se obtiene el factor exponencial anterior. También se puede escribir la salida de variaciones como
\begin{equation}
h_{i,n}=a_n(h_{i,n-1}+w_{i,n}-w_{i,n-1}).
\label{eq:turb-hp-equiv}
\end{equation}
Esta es una estructura de filtro pasaaltos: atenúa el nivel sostenido y deja responder a cambios. La fórmula equivalente permite compararla directamente con el pasaaltos de TelemFFB. No es una derivada pura, porque no divide el cambio de viento por $\Delta t$.

\subsubsection{Inicialización, viento constante y regreso del viento}
La primera muestra activa después de un reinicio hace $m_i=w_i$ y devuelve cero. Así, conectarse con viento ya existente no produce por sí solo un golpe. Si el viento queda constante en esos mismos ejes, $h$ tiende a cero.

Si se parte de viento cero estabilizado y el vertical salta a 5 kt, en la primera muestra de 20 ms se obtiene $m_Y\simeq0.099007$ kt y $h_Y\simeq4.900993$ kt. Si permanece en 5 kt, ese $h_Y$ va decayendo. Si posteriormente el viento baja, puede aparecer una perturbación de signo contrario: también se detecta la salida de la ráfaga. El filtro de salida puede demorar el cambio de signo efectivo del comando.

\subsection{Paso 2: mezclar viento hacia roll y pitch}
Antes de aplicar presión y ganancia, se calculan las siguientes señales:
\begin{align}
b_R&=0.70\,\frac{h_X}{5},\\
b_P&=\frac{h_Y+k_Zh_Z}{5},\qquad k_Z=0.25\text{ inicialmente}.
\label{eq:turb-mix}
\end{align}
El divisor tiene unidades de kt, igual que $h$, por lo que el resultado es adimensional. La denominación $b_R,b_P$ se utiliza aquí sólo para explicar este paso intermedio.

\subsubsection{Roll: componente lateral}
Un valor $h_X=5$ kt produce $b_R=0.70$; $h_X=-5$ produce $-0.70$. El 0.70 es un peso perceptual de mezcla, no una medición del torque de alabeo del avión. No se utiliza Y o Z directamente en la fórmula de roll, aunque el límite vectorial posterior puede hacer que pitch influya en la amplitud disponible para roll.

\subsubsection{Pitch: componente vertical y longitudinal}
La componente vertical tiene peso 1. La longitudinal tiene peso $k_Z$. Con valores iniciales, $h_Y=5$ kt da $b_P=1$ y $h_Z=5$ kt, por sí solo, da $b_P=0.25$. Variaciones del mismo signo se suman; de signos opuestos pueden cancelarse. Por ejemplo, $h_Y=1$ kt y $h_Z=-4$ kt dan $b_P=0$ con $k_Z=0.25$.

\code{--gust-longitudinal 0} elimina solamente el aporte Z. Un valor negativo invierte ese aporte sin invertir Y. Este ajuste no cambia el signo de todo pitch y tampoco afecta directamente el roll.

\subsubsection{Qué significa la escala de 5 kt}
$W=5$ kt establece cuánta variación se necesita para una señal de cierta amplitud. Bajar $W$ aumenta sensibilidad; subirlo la reduce, antes de saturar. \textbf{No significa que 5 kt de viento constante generen fuerza máxima}: se divide la variación filtrada $h$, no el viento absoluto, y luego siguen presión, ganancia, límite y suavizado. El valor está escrito directamente en \code{GustModel.h}; no se modifica con el menú de la placa.

\subsection{Paso 3: presión dinámica y ganancia previa}
\subsubsection{Presión utilizada por turbulencia}
La conversión es la misma base física utilizada para el centrado:
\begin{align}
q_{\rm Pa}&=47.88025898\,q_{\rm psf},\\
q_{\rm ref}&=\tfrac12(1.225)(0.514444444\,V_{\rm ref,kt})^2,\\
Q&=\sat_{[0,4]}(q_{\rm Pa}/q_{\rm ref}),\\
Q_T&=\sat_{[0,1]}(Q).
\label{eq:turb-q}
\end{align}
El último recorte es específico de la llamada al modelo de ráfagas. \textbf{Aunque la Pico calcule presión normalizada hasta 4 y el centrado la aproveche, la turbulencia usa como máximo 1.} Con referencia 120 kt, $q_{\rm ref}\simeq2334.24$ Pa. Si $Q=2$, el centrado puede seguir endureciéndose, pero esta ponderación de turbulencia ya vale $Q_T=1$.

NASA fundamenta $q=\rho V^2/2$; aplicar $Q_T$ a la perturbación de mando es una decisión del modelo local. No demuestra que la fuerza real de bisagra sea exactamente proporcional a este producto. La densidad 1.225 fija la referencia; la presión real se toma del simulador. Los 120 kt siguen siendo un ajuste, no una Vne certificada.

\subsubsection{Ganancia previa al límite}
La señal bruta es
\begin{align}
v_R&=\gamma Q_T\,0.70\frac{h_X}{5},\\
v_P&=\gamma Q_T\,\frac{h_Y+k_Zh_Z}{5}.
\label{eq:turb-raw}
\end{align}
\code{--gust-gain} fija $\gamma$, con valor inicial 1 y rango permitido $0\ldots10$. No modifica las ganancias guardadas en la placa. Aumentarla intensifica las señales pequeñas, pero puede hacer que el límite vectorial actúe con más frecuencia y que las diferencias entre ráfagas fuertes se compriman.

Si $Q_T$ o $\gamma$ se vuelven cero, el objetivo bruto es cero. \textbf{Una salida previa no desaparece necesariamente en ese instante}: decae por el filtro de 80 ms mientras el modelo continúe activo. La pausa, el suelo o un reinicio sí limpian directamente el estado del modelo.

\subsection{Paso 4: limitar conjuntamente los dos ejes}
\subsubsection{Norma del vector y recorte radial}
Se calcula la longitud del vector de fuerza bruta:
\begin{align}
\ell&=\sqrt{v_R^2+v_P^2},\\
(t_R,t_P)&=\frac{(v_R,v_P)}{\max(1,\ell)}.
\label{eq:turb-limit}
\end{align}
Si $\ell\le1$, no cambia nada. Si $\ell>1$, ambos valores se dividen por la misma longitud. Por ejemplo, $(v_R,v_P)=(0.9,1.2)$ tiene longitud 1.5 y se convierte en $(0.6,0.8)$. Se conserva la proporción y el signo del objetivo instantáneo.

\subsubsection{Consecuencia para ajustar roll y pitch}
El modelo impone $t_R^2+t_P^2\le1$, no dos máximos simultáneos independientes. No puede producir el objetivo $(1,1)$. Una ráfaga vertical fuerte puede reducir el roll disponible aunque $h_X$ no cambie, porque aumenta $\ell$ y se divide todo el vector.

Apagar ``Turbul. pitch'' en el menú no elimina ese efecto de reparto: la placa anula pitch \emph{después} de que el modelo de la Pico haya limitado ambos ejes. Para ensayar lateral sin esa interacción, la entrada vertical/longitudinal debe ser pequeña o el objetivo no debe saturarse. Esto es una consecuencia de la ubicación de los ajustes.

\subsubsection{Condición para evitar saturación}
Antes del suavizado, el objetivo no se recorta si
\begin{equation}
\frac{\gamma Q_T}{5}
\sqrt{(0.70h_X)^2+(h_Y+k_Zh_Z)^2}\le1.
\end{equation}
Con $Q_T=\gamma=1$, una variación exclusivamente lateral alcanza el límite alrededor de $|h_X|=7.143$ kt; una exclusivamente vertical, a 5 kt; una exclusivamente longitudinal con $k_Z=0.25$, a 20 kt. Son umbrales de \emph{variaciones filtradas}, no del viento meteorológico absoluto. Con viento mixto se puede saturar antes.

\subsection{Paso 5: suavizar el objetivo}
\subsubsection{Filtro de salida de 80 ms}
Después del límite vectorial:
\begin{align}
\beta_n&=1-e^{-\Delta t_n/\tau_L},\qquad \tau_L=0.08\ {\rm s},\\
o_{a,n}&=o_{a,n-1}+\beta_n(t_{a,n}-o_{a,n-1}),\quad a=R,P.
\label{eq:turb-lp}
\end{align}
$o_a$ es la salida anterior corregida una fracción de la distancia al nuevo objetivo. El código usa \code{-expm1(-dt/0.08)}, una forma numéricamente conveniente de calcular $1-e^{-dt/0.08}$. Con 20 ms, $\beta\simeq0.221199$: se recorre el 22.12\% de la diferencia pendiente por muestra.

La fórmula procede del mismo tipo de ecuación de primer orden que el seguimiento de media. El valor de 80 ms es un ajuste del proyecto. Si el objetivo fuese constante, desde cero alcanzaría el 63.2\% en 80 ms y el 95.0\% en 240 ms. Ante un escalón de viento, el objetivo no es constante porque el pasaaltos ya está decayendo; no deben confundirse ambas respuestas.

\subsubsection{Por qué la salida sigue dentro de $\pm1$}
El filtro combina la salida anterior y un objetivo dentro del disco de radio 1, con pesos positivos que suman uno. Como parte desde cero, conserva $o_R^2+o_P^2\le1$, salvo diferencias numéricas. Por lo tanto, cada componente queda entre $-1$ y $+1$ sin necesidad de usar toda esa escala en cada instante.

El límite radial conserva la dirección del objetivo en una muestra, pero el filtro combina historia de distintas direcciones. Luego, la zona muerta y las ganancias diferentes por eje también cambian la proporción final. No corresponde afirmar que toda la cadena hasta los motores conserve siempre la dirección del vector de viento.

\subsubsection{Qué sensaciones favorecen ambos filtros}
La referencia continua del pasaaltos tiene corte aproximado $f_H=1/(2\pi\tau_H)\simeq0.159$ Hz; el pasabajos, $f_L\simeq1.989$ Hz. Favorecen perturbaciones transitorias y reducen variaciones extremadamente lentas o muy rápidas. No crean una vibración de frecuencia fija ni anulan bruscamente todo lo situado fuera de esos valores.

Subir $\tau_H$ hace que el viento sostenido tarde más en ser absorbido por la media y prolonga el transitorio; baja la frecuencia de corte del pasaaltos. Subir $\tau_L$ suaviza más y retrasa la respuesta de salida. Los cortes son orientativos para la cadena lineal; muestreo, saturación y zona muerta modifican la respuesta completa.

\subsection{Paso 6: datos USB y ejecución en la placa}
La trama SC8 lleva viento X/Y/Z y presión en sus unidades originales, además de los parámetros de configuración. La Pico calcula $(o_R,o_P)$ en \code{GustModel.h}; el comando usa directamente esos flotantes como $z_R,z_P$, sin cuantización intermedia a enteros.

Los datos se solicitan por frame, el puente selecciona la última muestra y limita envíos a uno cada 20 ms, y los motores se actualizan nominalmente cada 2 ms. Los filtros de viento sólo avanzan una vez por nueva marca aceptada. Su salida se retiene entre muestras y se reinicia al deshabilitar o interrumpir el cálculo local.

\subsection{Paso 7: zona muerta y ganancias de la placa}
\subsubsection{Zona muerta continua del 2\%}
\code{SimTelemetry.h::gustCommand} aplica, con $d=0.02$:
\begin{equation}
D(z)=\begin{cases}
0,&|z|\le d,\\[2pt]
\displaystyle\operatorname{sgn}(z)\frac{|z|-d}{1-d},&|z|>d.
\end{cases}
\label{eq:turb-dead}
\end{equation}
Primero recorta $z$ a $[-1,1]$. Luego elimina las señales pequeñas; por encima del umbral resta la zona eliminada y divide por 0.98 para mantener $D(1)=1$. Ejemplo: $z=0.10$ produce $D=0.08/0.98\simeq0.081633$, no 0.10. Al cruzar 0.02 no hay salto finito: la función empieza desde cero.

El 2\% es un ajuste local de rechazo de señal pequeña, no un porcentaje de viento en kt ni una propiedad aerodinámica. Un valor de 0.02 en roll y otro en pitch se evalúan individualmente, aunque el modelo previo aplique un límite conjunto.

\subsubsection{Fórmulas por eje y valores actuales}
Después de esa transformación:
\begin{align}
u_R&=D(z_R)\,c_R\,G=0.315\,D(z_R),\\
u_P&=D(z_P)\,c_P\,G=0.630\,D(z_P).
\label{eq:turb-firmware}
\end{align}
Se usaron $c_R=0.35$, $c_P=0.70$ y $G=0.90$, definidos inicialmente por \code{SIMCONNECT_GUST_ROLL}, \code{SIMCONNECT_GUST_PITCH} y \code{SIMCONNECT_TOTAL_GAIN}. Las ganancias guardadas pueden reemplazarlos. ``Turbul. roll'' y ``Turbul. pitch'' permiten habilitar/deshabilitar cada aporte por separado.

No hay otro factor automático 0.5 en esta ruta: la diferencia de eje ya está en las ganancias 0.35/0.70 y en la mezcla lateral 0.70 del modelo. Los factores de otros efectos presentes en \code{SimEffects.h} no deben multiplicarse otra vez por la turbulencia.

\subsubsection{Por qué roll no es siempre la mitad de pitch}
Si llegan $z_R=z_P$, las ganancias iniciales dan comandos en relación 1:2. Pero para variaciones iguales $h_X=h_Y$ y $h_Z=0$, el modelo de la Pico produce roll con factor 0.70. Lejos de la zona muerta y sin otras distorsiones, la relación final es aproximadamente $u_R/u_P\simeq0.70(0.35/0.70)=0.35$. Cerca de la zona muerta la relación cambia aún más. No existe una proporción fija para todas las ráfagas.

\subsection{Paso 8: del comando de turbulencia al PWM}
\subsubsection{Cadena eléctrica para el efecto aislado}
La ruta es común con el centrado, pero se explicita aquí para poder leer esta sección de forma independiente. Si los demás efectos aportan cero:
\begin{align}
u_{{\rm out},a}&=\sat_{[-L,L]}(u_a),\qquad L=220/255,\\
z_{{\rm motor},a}&=\sat_{[-1,1]}(s_a u_{{\rm out},a}),\\
p_a&=\sat_{[-220,220]}\left(
\operatorname{round}_{\rm sim}(M z_{{\rm motor},a})\right),\quad M=255.
\end{align}
$s_a$ es la polaridad del motor, inicialmente $+1$; $M$ es \code{PWM_MAX}. El redondeo suma 0.5 a los positivos, resta 0.5 a los negativos y convierte a entero. Un comando 0.10 da 26 PWM y uno $-0.10$ da $-26$. No se multiplica por 220: se multiplica por 255 y se recorta después.

Con otros efectos activos, el límite $L$ se aplica a la suma de sus comandos. La turbulencia puede quedar parcial o totalmente ocultada por esa saturación, o reducir la salida si tiene signo contrario al resto. Amortiguación y fricción se suman como cargas independientes; no modifican aquí la fórmula de viento.

\subsubsection{Máximos con valores predeterminados}
Como $|D(z)|\le1$, los máximos individuales teóricos son
\begin{equation}
|u_R|\le0.315\Rightarrow |p_R|\le80,\qquad
|u_P|\le0.630\Rightarrow |p_P|\le161.
\end{equation}
Son límites individuales, no una afirmación de que ambos se alcancen juntos: el vector del modelo tiene norma no mayor que uno. La turbulencia aislada, con ganancias iniciales, no alcanza el techo 220. Con ganancias de menú mayores puede hacerlo; para un eje hace falta aproximadamente $G c_a\ge L$ y suficiente señal de entrada.

\subsubsection{Activación de pines}
Con $p\ge0$, LEFT queda en cero y RIGHT recibe $p$; con $p<0$, RIGHT queda en cero y LEFT recibe $|p|$. Para Pico: roll GP6/GP7, pitch GP8/GP9, enable GP10/GP11. Los signos físicos dependen del montaje y de \code{ROLL_MOTOR_DIRECTION}/\code{PITCH_MOTOR_DIRECTION}.

La frecuencia PWM configurada es 490 Hz y el rango 0--255. El duty cycle activo es $|p|/255$. Un PWM de 14 equivale a 5.49\%; uno de 161, a 63.14\%. No mide la fuerza real: la corriente, el motor y la transmisión determinan cuánto siente el piloto. El límite duro de 220 no regula corriente.

\subsection{Ejemplos numéricos desde viento hasta PWM}
\subsubsection{Ráfaga vertical de 5 kt, primera muestra}
Supuestos: avión en vuelo, modelo ya inicializado con viento cero, $\Delta t=0.02$ s, $Q_T=0.5$, $\gamma=1$, $k_Z=0.25$, ganancias de placa iniciales y demás efectos apagados. Llega $(w_X,w_Y,w_Z)=(0,5,0)$ kt.
\begin{enumerate}
\item Filtro de viento: $a=e^{-0.02}=0.9801987$; $m_Y=0.0990066$ kt; $h_Y=4.9009934$ kt. X y Z siguen en cero.
\item Mezcla: $b_P=h_Y/5=0.9801987$ y $b_R=0$.
\item Presión y ganancia: $v_P=(0.5)(1)(0.9801987)=0.4900993$; $v_R=0$.
\item Límite vectorial: $\ell=0.4900993<1$, por lo que $t_P=v_P$.
\item Suavizado: $\beta=1-e^{-0.25}=0.2211992$; desde salida cero, $o_P=0.1084096$.
\item La Pico usa directamente $z_P=o_P$; para las cuentas ilustrativas siguientes se redondea a 0.1084, sin representar un entero USB.
\item Zona muerta: $D(z_P)=(0.1084-0.02)/0.98\simeq0.0902041$.
\item Ganancias: $u_P=(0.0902041)(0.70)(0.90)=0.0568286$.
\item PWM: $\operatorname{round}(255\times0.0568286)=14$. Roll permanece en cero.
\end{enumerate}
Con polaridad pitch $+1$ se activa RIGHT con 14 y LEFT queda en cero. Si se invierte la polaridad, se invierte la salida activa. La identificación del sentido mecánico requiere comprobar el montaje.

\subsubsection{Qué ocurre si ese viento de 5 kt permanece constante}
Con los mismos supuestos y muestras cada 20 ms, la fuerza primero crece por el suavizado y después cae porque la media va absorbiendo el viento sostenido. Los siguientes números incluyen redondeo de los ejemplos y redondeo PWM:
\begin{center}
\begin{tabular}{rrrrr}
\toprule
Tiempo desde el cambio (s) & $h_Y$ (kt) & $o_P$ & $z_P$ & PWM pitch\\
\midrule
0.02 & 4.90099 & 0.108410 & 0.1084 & 14\\
0.08 & 4.61558 & 0.298876 & 0.2988 & 46\\
0.20 & 4.09365 & 0.396526 & 0.3965 & 62\\
0.50 & 3.03265 & 0.325448 & 0.3254 & 50\\
1.00 & 1.83940 & 0.198022 & 0.1980 & 29\\
2.00 & 0.67668 & 0.072849 & 0.0728 & 9\\
4.00 & 0.09158 & 0.009859 & 0.0098 & 0\\
\bottomrule
\end{tabular}
\end{center}
La tabla supone viento constante en ejes del avión, presión fija y muestreo regular. No es una medición de un vuelo real ni una promesa de amplitud para cualquier configuración meteorológica de ``5 kt''.

\subsubsection{Ráfaga mixta: el límite vectorial sí interviene}
Se parte nuevamente de viento y salida cero ya inicializados. Ahora llegan $w_X=w_Y=5$ kt, $w_Z=0$, con $Q_T=\gamma=1$ y $\Delta t=20$ ms. Las variaciones X/Y son $4.900993$ kt:
\begin{align*}
v_R&=0.6861391,&v_P&=0.9801987,\\
\ell&\simeq1.196485,&(t_R,t_P)&\simeq(0.5734623,0.8192319),\\
(o_R,o_P)&\simeq(0.1268494,0.1812135),&
(z_R,z_P)&=(0.1268,0.1812).
\end{align*}
Tras zona muerta y ganancias: $u_R\simeq0.0343286$ y $u_P\simeq0.1036286$. Se aplican aproximadamente \textbf{9 PWM roll y 26 PWM pitch}. No salen dos valores iguales ni una relación 1:2: intervinieron la mezcla lateral 0.70, el reparto vectorial y la zona muerta.

\subsubsection{Conversión rápida de señal recibida a PWM}
Esta tabla parte de $z$ ya recibido, por lo que no especifica una ráfaga única capaz de producirlo. Cada eje se evalúa por separado, con sus ganancias iniciales y polaridad positiva:
\begin{center}
\begin{tabular}{rrrrr}
\toprule
$z$ & $u_R$ & PWM roll & $u_P$ & PWM pitch\\
\midrule
0.02 & 0 & 0 & 0 & 0\\
0.03 & 0.003214 & 1 & 0.006429 & 2\\
0.05 & 0.009643 & 2 & 0.019286 & 5\\
0.10 & 0.025714 & 7 & 0.051429 & 13\\
0.25 & 0.073929 & 19 & 0.147857 & 38\\
0.50 & 0.154286 & 39 & 0.308571 & 79\\
1.00 & 0.315000 & 80 & 0.630000 & 161\\
\bottomrule
\end{tabular}
\end{center}
Para $z<0$ cambian los signos del comando y PWM; las magnitudes se conservan. Las dos columnas no implican que $z_R=z_P=1$ sea posible simultáneamente desde el modelo vectorial.

\subsection{Habilitación, cortes y reinicios}
\subsubsection{Condiciones del puente}
El modelo se calcula si \code{--gusts} está presente, la simulación corre sin pausa ni slew, las muestras son recientes y \code{SIM ON GROUND=0}. \code{--gusts} es falso por defecto en el ejecutable; el lanzador \code{Iniciar-Simulador.ps1} sí lo agrega junto con \code{--gust-gain 1}.

Para enviar telemetría armada se necesita \code{--enable} y puerto reconocido como \code{YOKE17SC8}. El modo monitor sin placa sólo registra las entradas; no calcula efectos.

\subsubsection{Validación específica del modelo}
\code{GustModel.h} reinicia y devuelve cero si hay una entrada no finita, $\Delta t<0$ o $\Delta t\ge0.25$ s, presión normalizada fuera de $[0,1]$, ganancia fuera de $[0,10]$, $|k_Z|>1$ o alguna componente de viento mayor que 1000 kt en módulo. Los 1000 kt son un filtro de sanidad de datos, no un umbral meteorológico de turbulencia.

Con $\Delta t=0$ y el modelo ya inicializado, conserva la salida anterior y no inventa tiempo transcurrido. Tras un reinicio, la primera muestra inicializa la media y devuelve cero. Si llega un salto temporal largo, no se interpreta todo ese salto como una ráfaga continua.

Un viento finito pero mayor que 1000 kt puede reiniciar sólo el modelo de ráfagas sin que necesariamente se desarme todo el puente; en cambio una entrada no finita interceptada por la validación global produce desarme. No todas las comprobaciones tienen idéntico alcance.

\subsubsection{Condiciones de la placa y latencia de corte}
La placa exige \code{SIMCONNECT_ADDITIONS_ENABLED}, máscara de turbulencias, ganancia efectiva del eje, telemetría armada de menos de 300 ms y DTR activo. El menú de calibración fuerza los comandos de motor a cero. Al guardar, también se detienen antes de escribir memoria.

El suelo o una pausa reinician el modelo en la Pico; la retirada física depende de que llegue y se procese la trama correspondiente. Ante ausencia de tramas, interviene el watchdog. No debe interpretarse ``reinicio inmediato del modelo'' como latencia eléctrica absolutamente nula.

El reinicio de \code{Effects::update} por pérdida de habilitación o intervalo local mayor que 50 ms limpia todos los modelos SC8, incluyendo turbulencia, además de centro y portadoras. La siguiente muestra dinámica inicializa referencias sin crear un impulso artificial.

\subsection{Procedencia de las fórmulas y constantes}
\subsubsection{Qué aporta TelemFFB}
Se compararon las copias locales de \code{TurbulenceModulator.py}, \code{TurbulenceMixIn.py} y la definición de \code{RelWind} en \code{SimConnectManager.py}. El comentario de \code{GustModel.h} señala expresamente a TelemFFB como referencia para filtrado, mezcla y limitación vectorial.

El pasaaltos de TelemFFB usa $\alpha=\alpha_{\rm ref}^{\Delta t/\Delta t_{\rm ref}}$ y $h_n=\alpha(h_{n-1}+w_n-w_{n-1})$. Se relaciona con MiYoke mediante
\begin{equation}
\tau_H=-\frac{\Delta t_{\rm ref}}{\ln\alpha_{\rm ref}},
\qquad \alpha=e^{-\Delta t/\tau_H}.
\end{equation}
La equivalencia es de estructura matemática; no implica que se hayan usado los mismos tiempos. Con los argumentos predeterminados del método de TelemFFB, $\alpha_{\rm ref}=0.95$ a 60 Hz equivale aproximadamente a 0.325 s, mientras MiYoke usa 1 s. Los ajustes efectivos de TelemFFB pueden reemplazar esos argumentos.

TelemFFB utiliza pesos vertical 1, longitudinal 0.25 y lateral 0.70. Para roll aplica signo negativo; MiYoke usa \textbf{positivo}, según su comentario para conservar la convención del puente existente. La comparación no valida qué polaridad física es correcta para este mecanismo.

También hay límite conjunto de pitch/roll y suavizado exponencial en la referencia. TelemFFB ajusta sensibilidad mediante un divisor entre 10 y 1 y una intensidad posterior al límite de magnitud. MiYoke usa 5 kt, presión normalizada y ganancia \emph{antes} del límite unitario. Por tanto, \code{--gust-gain} no equivale exactamente al parámetro de intensidad de TelemFFB.

\subsubsection{La entrada de viento no es la misma}
\code{SimConnectManager.py} de TelemFFB define \code{RelWind} con \code{RELATIVE WIND VELOCITY BODY X/Y/Z} en m/s. Este proyecto usa \code{AIRCRAFT WIND X/Y/Z} en kt. No son nombres intercambiables: cambiar de familia de variables modifica el significado de la entrada, y cambiar a m/s exige revisar la escala de 5 kt.

TelemFFB convierte los resultados en magnitud y dirección para un efecto de fuerza constante del dispositivo FFB; ``constante'' describe el tipo de orden enviado, no que su magnitud no varíe con el tiempo. MiYoke transmite entradas crudas por USB serie; la Pico calcula las componentes y las convierte a PWM. No utiliza la conversión angular ni el firmware Rhino para generar esta salida.

\subsubsection{Qué procede del proyecto alemán original}
En las fuentes inspeccionadas de \code{../../PROEYCTO ORIGINAL/Arduino_FFB_Yoke} no se identificó una cadena con \code{AIRCRAFT WIND}, SimConnect y estos filtros de ráfagas. El original aporta antecedentes de lectura de ejes, manejo FFB y accionamiento de motor; no se le atribuyen las ecuaciones de viento actuales.

La arquitectura de dos PWM según signo tiene un antecedente concreto en \code{AxisCalibration.cpp::MoveMotor}. Como se explicó en centrado, el original contiene también otra conversión de fuerza en \code{Arduino.ino}, por lo que no debe afirmarse identidad completa entre todos los mapas de PWM del original y los actuales.

\subsubsection{Ajustes propios, fundamentos y trazabilidad}
\begin{longtable}{p{0.43\linewidth}p{0.49\linewidth}}
\toprule Parte & Procedencia y alcance\\\midrule\endhead
Variables de viento y presión & SDK de Prepar3D; describe telemetría, no una ley de fuerza de mando.\\
$q=\rho V^2/2$ & Fundamento físico general, referencia NASA; no fija ganancias hápticas.\\
Pasaaltos, mezcla y límite vectorial & Estructura comparable con TelemFFB y antecedente declarado en el código local.\\
Pesos 1, 0.25 y 0.70 & Presentes en TelemFFB; adaptación local con signo roll positivo y $k_Z$ regulable. No son mediciones de un avión concreto.\\
1 s, 80 ms y 5 kt & Valores de ajuste de MiYoke documentados como continuidad del modelo local; no constantes universales ni defaults idénticos a TelemFFB.\\
Referencia 120 kt y producto $\gamma Q_T$ & Implementación local de ponderación y normalización.\\
Zona muerta 0.02 y ganancias 0.35/0.70 & Ajustes locales vigentes de la placa.\\
SC8, watchdog y menú persistente & Implementación local de comunicación y control.\\
Redondeo, escala 255 y techo 220 & Ruta eléctrica de MiYoke; no calibración de fuerza física.\\
\bottomrule
\end{longtable}
\code{SimConnectBridge/GustGain.h} contiene \code{scaleGust}, pero no está incluido ni llamado por el puente actual: aparece en pruebas. No se agregó su fórmula como una etapa extra, porque hacerlo duplicaría un cálculo que no está en la ruta de ejecución analizada.

\subsection{Qué puede ajustarse y cómo cambia la sensación}
\subsubsection{Ganancias disponibles en la placa}
\begin{longtable}{p{0.25\linewidth}p{0.13\linewidth}p{0.54\linewidth}}
\toprule Ajuste & Inicio & Resultado esperado\\\midrule\endhead
Turbul. roll, $c_R$ & 0.35 & Escala sólo roll después de recibir y transformar la señal.\\
Turbul. pitch, $c_P$ & 0.70 & Escala sólo pitch después de recibir y transformar la señal.\\
Ganancia total, $G$ & 0.90 & Escala ambos ejes y el resto de efectos.\\
\bottomrule
\end{longtable}
El rango es $0\ldots1$ y el paso de botones 0.05. Las ganancias se guardan en memoria; sus valores efectivos pueden ser distintos de los \code{define}. Los interruptores por eje hacen efectiva una ganancia cero sin cambiar el número guardado. Subir estas ganancias aumenta PWM, pero no recupera variaciones que el modelo ya saturó o que la zona muerta ya eliminó.

\subsubsection{Ajustes del puente}
\begin{longtable}{p{0.34\linewidth}p{0.58\linewidth}}
\toprule Parámetro & Efecto y límite de interpretación\\\midrule\endhead
\code{--gust-gain}, $0\ldots10$ & Amplifica antes del límite vectorial y del filtro. Puede hacer visibles ráfagas pequeñas y saturar las grandes.\\
\code{--gust-longitudinal}, $-1\ldots1$ & Peso de Z sobre pitch. Cero sirve para evaluar sólo aporte vertical; un valor negativo invierte sólo Z.\\
\code{--reference-knots}, 40--250 & Regula $q_{\rm ref}$. Más referencia reduce $Q_T$ a presión fija, salvo regiones recortadas. También afecta centrado.\\
\code{--gusts} & Habilita el cálculo de ráfagas; no sustituye a \code{--enable} para actuar.\\
\bottomrule
\end{longtable}
Subir referencia de 120 a 150 kt multiplica la presión normalizada \emph{sin recortar} por 0.64. Si ambos casos siguen por encima de uno, $Q_T$ permanece en uno y no cambia la ponderación de turbulencia. Es distinto del centrado SC8, que conserva valores hasta cuatro.

\subsubsection{Cambios que requieren editar y recompilar}
La escala 5 kt, el peso lateral 0.70 y los tiempos 1 s/80 ms están escritos en \code{GustModel.h}: requieren recompilar el firmware. La zona muerta 0.02 se encuentra en \code{SimTelemetry.h::gustCommand}: requiere recompilar y cargar firmware. Si se cambia $d$, se debe cambiar coherentemente el divisor $1-d$, actualmente 0.98.

Las polaridades de motor y límites PWM requieren revisar \code{mis_defines.h}. Invertir la polaridad de un motor invierte todos sus efectos, incluido el centrado. No conviene usar ese ajuste para corregir exclusivamente el signo de una turbulencia cuando el resorte ya restaura correctamente: en ese caso hay que revisar el mapeo del viento del modelo.

\subsubsection{Umbral efectivo por cuantización PWM}
Superar $d=0.02$ genera un comando matemático, pero puede seguir redondeando a cero. Para un eje aislado con $M=255$ y $Gc_a>0$, la primera cuenta aparece aproximadamente cuando
\begin{equation}
|z_a|\gtrsim d+(1-d)\frac{0.5}{255Gc_a}.
\end{equation}
Con valores iniciales resulta 0.02610 para roll y 0.02305 para pitch. Además intervienen la precisión de flotantes y las pérdidas del mecanismo. Una cuenta de PWM no garantiza movimiento o una fuerza perceptible. Por eso, observar una señal pequeña en el CSV no implica necesariamente sentirla en la mano.

\subsection{Procedimiento para calibrar y pulir el efecto}
\subsubsection{Primero comprobar que existe señal de viento}
Usar el modo monitor permite registrar las entradas sin armar motores. Por ejemplo, desde la carpeta del ejecutable:
\begin{quote}\small
\code{YokeSimConnectBridge.exe --seconds 60 --gusts --gust-gain 1}
\end{quote}
Comparar viento sostenido, cambios de viento y maniobras con viento sostenido, manteniendo avión y velocidad de simulación. Si las tres componentes no varían, subir ganancias no crea turbulencia. Si varían sólo al maniobrar, hay que considerar el efecto de las coordenadas del avión antes de atribuirlo a la meteorología.

\subsubsection{Ensayo de fuerzas aisladas}
\begin{enumerate}
\item Registrar modo de firmware, ganancias guardadas, avión, meteorología y argumentos del puente. El efecto requiere vuelo: una prueba en tierra no reproduce la ruta activa.
\item Con pausa reconocida o comunicación desconectada, dejar en el menú sólo las turbulencias que se ensayarán; apagar centrado, otras vibraciones, damper y fricción. \code{--no-buffets} no apaga por sí solo todos esos aportes.
\item Retomar el ensayo con ganancias pequeñas, un eje por vez, comprobando correspondencia entre signo de señal, comando y PWM. Para actuar se necesitan puerto real, \code{--enable} y \code{--gusts}; el comando de monitor anterior no lo hace.
\item Ajustar primero el balance con $c_R$ y $c_P$. Después evaluar $\gamma$ para sensibilidad general a ráfagas pequeñas y comprobar que no mantenga el objetivo casi siempre saturado.
\item Comparar $k_Z=0$ y $k_Z=0.25$ con las mismas condiciones, para distinguir cuánto aporta el viento longitudinal a pitch.
\end{enumerate}
Mantener una configuración de referencia entre ensayos evita confundir cambios de meteorología con cambios del algoritmo. La reproducibilidad exige capturar la serie de viento; dos vuelos diferentes no son automáticamente la misma entrada.

\subsubsection{Qué cambiar ante cada síntoma}
\begin{longtable}{p{0.29\linewidth}p{0.63\linewidth}}
\toprule Síntoma & Comprobación o ajuste relacionado\\\midrule\endhead
No aparece fuerza & Revisar vuelo, \code{--gusts}, armado, frescura, DTR y ganancias. Confirmar variaciones reales del viento y señal por encima de zona muerta.\\
Ráfagas fuertes se sienten iguales & Inspeccionar si el objetivo vectorial satura o si la suma llega a 220. Bajar $\gamma$ puede recuperar diferencias antes del límite.\\
Pitch domina demasiado & Ajustar $c_P$ y evaluar $k_Z$, diferenciando contribución vertical y longitudinal.\\
Respuesta demasiado lenta & Revisar frecuencia de muestras, $\tau_L$ y presión; cambiar $\tau_H$ afecta también la duración del transitorio.\\
Fuerza asociada a giros en viento estable & Revisar cambios de coordenadas del avión; no atribuirlo automáticamente a un problema de ganancia.\\
CSV indica señal pero no se siente & Comparar señal recuperada, zona muerta, redondeo PWM y fricción del mecanismo.\\
Turbulencia cambia al mover el yugo & Aislar centrado y cargas de movimiento; la ley de viento no utiliza directamente esa posición.\\
\bottomrule
\end{longtable}

\subsection{Diagnóstico: qué observar y qué puede inducir a error}
\subsubsection{CSV del puente}
El CSV conserva viento X/Y/Z, presión original \code{pressure_psf}, configuración de ráfagas, estados y marcas de captura/envío. SC8 no registra \code{gust_roll}, \code{gust_pitch}, \code{q_ratio} ni \code{Turbulencia} calculados en la PC. Para reconstruir una prueba fuera de línea deben ejecutarse los modelos de la Pico con las marcas de las muestras enviadas. El monitor sólo recoge telemetría.

\subsubsection{Datos de la placa y mejoras de registro}
El diagnóstico local \code{f=} muestra la suma de comandos multiplicada por 10000 y \code{p=} el PWM aplicado. Sólo con el efecto aislado se puede comparar esa suma directamente con $u_R/u_P$. El LCD y las líneas de depuración no sustituyen un registro a la frecuencia del control y pueden omitir picos breves.

Una ampliación útil sería registrar $h_X,h_Y,h_Z$, $Q_T$, $\ell$, señal antes/después del límite, $o_a$, $z_a$ y aporte aislado de turbulencia. Así se distinguiría una ráfaga eliminada por el filtro de una señal recortada, una zona muerta o un techo de la suma. Este análisis no añadió esas variables al firmware ni al puente.

\subsection{Mejoras posibles, todavía no implementadas}
\subsubsection{Separar causas atmosféricas y maniobras}
Para reducir respuesta a cambios de actitud en viento terrestre constante se puede estudiar viento en coordenadas terrestres y una transformación coherente al cuerpo del avión, con actitud y tiempos alineados. Esa mejora exige nuevas variables y validación; sustituir nombres de SimVars sin revisar unidades y referencias no la resuelve.

\subsubsection{Hacer explícitos los parámetros y los límites}
Se podrían exponer $\tau_H$, $\tau_L$, $W$ y pesos de mezcla como parámetros con rangos validados y guardarlos junto al registro del ensayo. Conviene distinguir una sensibilidad previa al límite y una intensidad posterior: actualmente $\gamma$ modifica cuánto se satura el modelo, mientras $c_R/c_P$ regulan la salida final por eje.

También sería útil diagnosticar saturación vectorial aunque un eje esté apagado en la placa. Cambiar a límites independientes o hacer que el modelo respete los interruptores por eje antes del límite modifica el comportamiento; debe hacerse como decisión de diseño, no como si ya estuviera implementado.

\subsubsection{Calibrar sensación y fuerza física}
Para calibrar magnitud real hay que medir fuerza/torque frente a PWM en el mecanismo y verificar ambos sentidos. Los coeficientes de TelemFFB no sustituyen esa identificación. Una regulación de corriente sería otra ampliación distinta del modelo de viento.

La calibración ADC del yugo sigue siendo importante para centrado y cargas de movimiento, pero no es una ganancia de turbulencia. No conviene desplazar el centro del sensor para corregir una turbulencia excesiva o asimétrica; se debe revisar la serie de viento, mezcla, signo y ganancias específicas.

\subsection{Mapa de fuentes y verificación}
\subsubsection{Implementación vigente}
\begin{longtable}{p{0.39\linewidth}p{0.53\linewidth}}
\toprule Archivo & Elementos comprobados\\\midrule\endhead
\code{SimConnectBridge/main.cpp} & Variables originales, unidades, validación de transporte, argumentos, trama y CSV.\\
\code{GustModel.h} & \code{Model::update}: inicialización, media, mezcla, presión, ganancia, límite vectorial y suavizado.\\
\code{SimTelemetry.h} & Decodificación cruda, watchdog y \code{gustCommand}: zona muerta y ganancias.\\
\code{SimEffects.h} & Aplicación por eje, habilitación, convivencia con otras cargas y saturación total.\\
\code{mis_defines.h} & Ganancias iniciales 0.35/0.70/0.90, banderas, polaridades y límites.\\
\code{GainSettings.h} & Ganancias efectivas, pasos del menú y estado por eje.\\
\code{MiYoke_PD_V2.4.ino} & Lazo de control, menú, aplicación de polaridad y diagnóstico.\\
\code{MotorDriver.h} & Escala, redondeo, límite duro y salidas PWM.\\
\code{Iniciar-Simulador.ps1} & Invocación habitual con \code{--gusts --gust-gain 1}.\\
\code{tests/simconnect_checks.cpp} & Casos existentes de viento constante, mezcla, saturación, reinicio, zona muerta, tasa de muestras y comunicación.\\
\bottomrule
\end{longtable}
Las rutas son relativas a MiYoke\_PD\_V2.4. Se verificaron numéricamente las tablas aplicando las ecuaciones en el mismo orden; no son medidas de motores. Las constantes y estructuras se contrastaron con el código, no únicamente con los documentos explicativos anteriores.

\subsubsection{Documentación histórica y referencias públicas}
\code{LEEME md/MODELO_TURBULENCIA.md} describe el fundamento y las ecuaciones del modelo, pero conserva pasajes antiguos sobre protocolo SC2, FFB HID y ganancias 0.10/0.20. Esos detalles no describen la integración actual. Esta sección usa SC8, ganancias 0.35/0.70 y suma exclusivamente generada por la ruta actual de SimConnect.

\begin{enumerate}
\item \href{https://www.prepar3d.com/SDKv3/LearningCenter/utilities/variables/simulation_variables.html}{Prepar3D SDK v3, Simulation Variables}: componentes \code{AIRCRAFT WIND X/Y/Z} y presión. Consultado el 7 de octubre de 2026.
\item \href{https://www1.grc.nasa.gov/beginners-guide-to-aeronautics/dynamic-pressure-2/}{NASA Glenn, Dynamic Pressure}: fundamento de presión dinámica, ya utilizado en la sección de centrado.
\item \href{https://github.com/walmis/VPforce-TelemFFB/blob/refactor/telemffb/utils/TurbulenceModulator.py}{TelemFFB, TurbulenceModulator.py}: referencia pública para localizar filtros, mezcla y saturación; comparación realizada contra la copia local.
\item \href{https://github.com/walmis/VPforce-TelemFFB/blob/refactor/telemffb/sim/msfs_xp/TurbulenceMixIn.py}{TelemFFB, TurbulenceMixIn.py}: uso de componentes y envío de fuerza FFB; copia local inspeccionada.
\item \href{https://github.com/walmis/VPforce-TelemFFB/blob/refactor/telemffb/telem/SimConnectManager.py}{TelemFFB, SimConnectManager.py}: definición de \code{RelWind} con viento relativo en m/s; copia local inspeccionada.
\end{enumerate}
Las ramas públicas pueden evolucionar. Los tiempos, ganancias y conclusiones de esta sección corresponden a los archivos locales inspeccionados. No se presentan como identificaciones experimentales de una aeronave ni como una validación del firmware efectivamente cargado.
\newpage
\section{VIBRACIONES EN SUELO}
\label{sec:suelo}

\subsection{Actualización: textura suave de rodaje sobre pista}
La salida de suelo combina ahora las sacudidas de telemetría con una textura sintética suave, calculada en la Pico en cada ciclo de motores. Esta textura representa rodaje aun cuando Prepar3D entrega aceleraciones prácticamente constantes sobre pista. No identifica material ni baches individuales. Las secciones de filtrado siguientes describen la componente de sacudidas; la salida total incorpora también esta etapa.
\begin{align*}
T(t)&=0.12E\,[0.65\sin(2\pi\,7t)+0.35\sin(2\pi\,11.3t)],\\
F_{a,\mathrm{total}}&=\sat_{[-E,E]}(F_a+T(t)).
\end{align*}
$E$ conserva exactamente la rampa de 0.5 a 40 kt. Las frecuencias son fijas y la velocidad modifica la intensidad. El reloj local de textura se reinicia al detenerse, salir de suelo o deshabilitar efectos; se repite cada 10 s sin discontinuidad de fase. No depende de que lleguen nuevas aceleraciones en cada ciclo. Se aplica únicamente en la ruta SC8 y se suma antes de ganancias y límites existentes.

Con suelo 0.450 y total 0.90, a $E=1$ el techo teórico de textura aislada es $255(0.12)(0.450)(0.90)=12.393$ PWM en roll, aproximadamente 12 tras redondeo, y 1 en pitch por su divisor diez. Los baches mantienen los techos conjuntos de suelo de 40/4 PWM. Las senoides aportan variedad sin introducir ruido aleatorio; su percepción en el mecanismo todavía requiere ensayo físico.

El diagnóstico \code{gr/gp} conserva la componente de aceleraciones previa a la textura; R/P incluyen la suma final. Las cifras de reproducción de la revisión sin textura son históricas: con esta textura, el mismo CSV genera 800 muestras con PWM de suelo no nulo, frente a 375 sin textura. La suite actual ejecutó 96\,764 comprobaciones, incluyendo pista con aceleración constante, reposo y salida de suelo. Los ejemplos siguientes del modelo de aceleraciones no incluyen la fase de textura salvo indicación expresa.

\subsection{Qué hace el efecto en la versión SC8}
El efecto responde a las \textbf{variaciones de aceleración lateral y vertical que entrega Prepar3D mientras el avión rueda}. La velocidad respecto del terreno conserva exactamente la progresión solicitada: hasta 0.5 kt no hay aporte; entre 0.5 y 40 kt aumenta el multiplicador de intensidad; desde 40 kt ese multiplicador vale uno.

La explicación sencilla es: \emph{``transformar las sacudidas medidas del avión en comandos de motor, y hacer que una misma sacudida se transmita con mayor intensidad cuanto más rápido se rueda, hasta el límite elegido''}. La componente de baches puede cambiar de signo o desaparecer si las aceleraciones se estabilizan; la textura base continúa mientras el avión rueda por encima de 0.5 kt. \textbf{Máxima ganancia desde 40 kt no significa fuerza máxima permanente.}

\subsubsection{Cambio respecto de la implementación anterior}
En SC3--SC6 la placa sintetizaba una triangular de 8 Hz y suavizaba su amplitud con 100 ms. Esa ruta permanece para recibir protocolos antiguos, pero \textbf{no se ejecuta para el suelo de SC8}. Ahora la Pico calcula dos señales con signo y aplica ganancias y límites directamente. La componente de aceleraciones no tiene portadora ni suavizado adicional de 100 ms; se suma la textura local descrita arriba. Los buffets aerodinámicos mantienen sus portadoras y filtros anteriores.

A igual velocidad, dos recorridos pueden producir comandos distintos si cambian las aceleraciones. El algoritmo no consulta un mapa de baches, no distingue cada rueda ni lee tipo de superficie: responde al movimiento que el simulador efectivamente publica. Puede responder también a maniobras bruscas en suelo. No hay un efecto separado de touchdown, ni reproducción mecánica calibrada de la transmisión del tren al yugo.

\subsection{Recorrido completo desde SimConnect hasta los motores}
\begin{center}
\fbox{\begin{minipage}{0.91\linewidth}
\textbf{En Windows:} aceleraciones X/Y, velocidad y estado de suelo se recogen y transmiten en sus unidades originales, con marca temporal.

\textbf{En la Pico:} validación, pasaaltos de 3 Hz, pasabajos de 20 ms, mezcla y normalización, saturación y rampa de velocidad; todo una vez por nueva muestra SC8.

\textbf{En la placa:} señales de suelo calculadas localmente $\rightarrow$ ganancia ``Vib. suelo'' y total $\rightarrow$ recorte propio por eje $\rightarrow$ división adicional de pitch por diez $\rightarrow$ suma con otros efectos $\rightarrow$ límite general, polaridad, redondeo y PWM.
\end{minipage}}
\end{center}
El signo de las perturbaciones se calcula en la Pico. Entre paquetes se conserva la componente calculada de aceleraciones; la textura continúa con el reloj local.

\subsection{Entradas y condiciones de ejecución}
\begin{longtable}{p{0.34\linewidth}p{0.15\linewidth}p{0.43\linewidth}}
\toprule Variable & Unidad & Uso\\\midrule\endhead
\code{ACCELERATION BODY X} & ft/s$^2$ & Componente lateral de entrada.\\
\code{ACCELERATION BODY Y} & ft/s$^2$ & Componente vertical de entrada.\\
\code{GROUND VELOCITY} & kt & Multiplicador progresivo 0.5--40 kt.\\
\code{SIM ON GROUND} & Bool & Debe valer exactamente 1.\\
\code{IS SLEW ACTIVE} & Bool & Debe valer cero.\\
\bottomrule
\end{longtable}
El puente solicita estas entradas como \code{FLOAT64}. Los eventos \code{Sim} y \code{Pause} exigen simulación en marcha y sin pausa. El \href{https://www.prepar3d.com/SDKv4/sdk/references/variables/simulation_variables.html}{SDK de Prepar3D v4, Simulation Variables} documenta ambas aceleraciones en pies por segundo al cuadrado. No enumera la variable \code{CONTACT POINT COMPRESSION} usada por TelemFFB en MSFS; por ello esta adaptación utiliza aceleraciones documentadas y no presupone compatibilidad de esa variable.

No se usan IAS, presión dinámica, posición o velocidad del yugo en la fórmula de suelo. Tampoco se solicitan compresiones de ruedas, RPM de neumáticos, geometría de suspensión ni tipo de superficie. Un avión detenido permanece sin este aporte aunque tenga viento relativo. Las opciones \code{--gusts}, \code{--gust-gain}, \code{--gust-longitudinal}, \code{--reference-knots} y \code{--no-buffets} no modifican este modelo.

La validación compartida sigue pudiendo desarmar el conjunto por presión o posiciones del simulador inválidas, IAS/viento no finitos o estado de suelo inválido. Las comprobaciones particulares de aceleración y velocidad, en cambio, reinician únicamente el modelo de suelo y devuelven sus fuerzas en cero.

\subsection{Símbolos y parámetros implementados}
\begin{longtable}{p{0.25\linewidth}p{0.67\linewidth}}
\toprule Símbolo & Significado\\\midrule\endhead
$V_g$, $E$ & Velocidad en kt y multiplicador de velocidad, entre 0 y 1.\\
$a_{X,n},a_{Y,n}$ & Aceleraciones de entrada, en ft/s$^2$.\\
$\Delta t_n$ & Diferencia entre marcas de captura aceptadas, en segundos.\\
$h_{i,n}$, $\ell_{i,n}$ & Salida del pasaaltos y del pasabajos, respectivamente.\\
$\tau_H=1/(2\pi\,3)$ s & Constante del pasaaltos; referencia continua de corte 3 Hz.\\
$\tau_L=0.020$ s & Constante del pasabajos.\\
$a_*=8.043512$ ft/s$^2$ & Escala aproximada de 0.25 g para normalización.\\
$s_R,s_P$ & Señales normalizadas y limitadas antes de ponderar velocidad.\\
$F_R,F_P$ & Señales con signo calculadas por la Pico: $F_a=E s_a$.\\
$\widehat F_R,\widehat F_P$ & Señales locales de la Pico, sin cuantización USB intermedia.\\
$c_g=0.18$, $G=0.90$ & Ganancia de suelo y ganancia total iniciales.\\
$L_g=40/255$, $r_g=10$ & Techo propio y divisor adicional para pitch.\\
$u_R,u_P$ & Aportes normalizados de suelo a la suma de efectos.\\
\bottomrule
\end{longtable}
Las ganancias pueden ser reemplazadas por ajustes persistidos en la placa. Los tiempos de filtro, pesos de mezcla y escala de aceleración son constantes del firmware; no son parámetros del menú ni valores identificados mediante mediciones del avión.

\subsection{Paso 1: conservar la progresión con velocidad}
\code{GroundModel.h::envelope} conserva la ley:
\begin{equation}
E=\begin{cases}
0,&\text{si no está activo, no está en suelo o }V_g\text{ no es finita},\\
0,&V_g\le0.5\ {\rm kt},\\
\displaystyle\frac{V_g-0.5}{39.5},&0.5<V_g<40\ {\rm kt},\\[3pt]
1,&V_g\ge40\ {\rm kt}.
\end{cases}
\label{eq:ground-envelope}
\end{equation}
``Activo'' identifica la habilitación del modelo; en la ruta SC8 de producción requiere telemetría armada y vigente, además de conexión USB. El modo monitor SC8 sólo registra las entradas y no ejecuta el modelo.

El denominador es $40-0.5=39.5$ kt. A 10 kt resulta $E\simeq0.240506$; a 20.25 kt, $E=0.5$; a 40 y 60 kt, $E=1$. Una velocidad negativa da cero: no se toma su valor absoluto. No se implementa histéresis alrededor de 0.5 kt.

Esta rampa es una decisión local de sensación. No se deriva de un modelo de suspensión ni de TelemFFB. Multiplica la perturbación después de filtrarla y limitarla, de modo que el máximo normalizado disponible por eje es $E$.

\subsection{Paso 2: validar e inicializar las aceleraciones}
El modelo devuelve cero y borra su memoria si $E=0$, si alguna aceleración no es finita o tiene magnitud mayor que 640 ft/s$^2$, o si $\Delta t$ no es finito, es negativo o alcanza 0.25 s. La cota 640 es una validación de sanidad, no un límite operativo del avión.

La primera muestra válida después de un reinicio establece las dos referencias previas:
\begin{equation}
a_{i,\rm prev}=a_{i,n},\qquad h_i=\ell_i=0,\qquad F_R=F_P=0.
\end{equation}
No se interpreta el nivel inicial de aceleración como un golpe. Esto evita un impulso artificial al arrancar o reanudar, pero también significa que la primera muestra válida tras un contacto con suelo no constituye un efecto de impacto separado. Sólo los cambios posteriores generan respuesta.

El tiempo se obtiene de la marca de captura de la muestra seleccionada por el puente. Los lotes se reducen a su último dato antes del envío. La Pico procesa una vez cada nueva marca; ante $\Delta t=0$ conserva estado y salida en lugar de reiniciar. No se supone un período exacto de 50 Hz ni se asigna artificialmente 20 ms a datos simultáneos.

\subsection{Paso 3: filtrar las variaciones de movimiento}
Para cada componente $i\in\{X,Y\}$, en una actualización válida posterior a la inicialización:
\begin{align}
\alpha_n&=\frac{\tau_H}{\tau_H+\Delta t_n}, & \tau_H&=\frac{1}{2\pi\,3}\simeq0.053052\ {\rm s},\\
h_{i,n}&=\alpha_n\bigl(h_{i,n-1}+a_{i,n}-a_{i,n-1}\bigr),\\
\beta_n&=\frac{\Delta t_n}{0.020+\Delta t_n},\\
\ell_{i,n}&=\ell_{i,n-1}+\beta_n(h_{i,n}-\ell_{i,n-1}).
\label{eq:ground-filter}
\end{align}
Después se guarda $a_{i,n}$ como referencia para la siguiente muestra. El pasaaltos elimina la componente constante y atenúa variaciones lentas; el pasabajos suaviza los cambios rápidos. Sus referencias continuas son 3 Hz y aproximadamente 7.958 Hz, respectivamente. \textbf{No son frecuencias de osciladores ni cortes ideales}; la respuesta digital depende también del intervalo de muestreo.

Con $\Delta t=0.020$ s, $\alpha\simeq0.726221$ y $\beta=0.5$. Si una aceleración cambia y después queda constante, la salida presenta un transitorio y tiende a cero. Si permanece constante desde la inicialización, el modelo permanece silencioso. No se resta explícitamente 1 g: cualquier nivel constante se elimina por el filtrado.

El suavizado pertenece al modelo de la Pico. Cambiar los 100 ms que aún aparecen en \code{Effects::update} modifica centrado y buffets, pero no suaviza las fuerzas de suelo de SC8.

\subsection{Paso 4: obtener dos señales y ponderarlas por velocidad}
La mezcla implementada es:
\begin{align}
s_R&=\sat_{[-1,1]}\left(\frac{0.3\ell_X+0.7\ell_Y}{8.043512}\right),\\
s_P&=\sat_{[-1,1]}\left(\frac{\ell_Y}{8.043512}\right),\\
F_R&=E\,s_R, & F_P&=E\,s_P.
\label{eq:ground-signals}
\end{align}
La componente vertical también excita roll para conservar el carácter predominante de ese eje en este equipo. Los pesos 0.3/0.7 y la escala aproximada de 0.25 g son un \textbf{mapeo háptico elegido}, no una transformación geométrica del tren al mando ni una fórmula copiada de TelemFFB.

La saturación se aplica antes de $E$. A 20.25 kt las mismas señales filtradas producen la mitad de señal local que a 40 kt; a partir de 40 kt no aumenta ese multiplicador. Esto no garantiza que dos recorridos reales a velocidades diferentes produzcan la misma aceleración de entrada. Tampoco garantiza una fuerza física proporcional después de saturación, redondeo y transmisión mecánica.

Las componentes de aceleraciones pueden diferir en magnitud y signo y no tienen fase común. La textura agregada sí comparte fase entre ejes. La suma no impone una relación instantánea obligatoria 10:1. Un transitorio puede producir fuerza de un solo signo durante un intervalo.

\subsection{Paso 5: procesar la telemetría SC8 en la Pico}
La trama transporta aceleraciones X/Y en bytes 90--93 y 94--97, velocidad de suelo en 78--81 y estado de suelo en 30--33, todos float32. La marca temporal ocupa 6--9. Envolvente y fuerzas se calculan en la placa. Antes de ganancias, \code{SimEffects.h} suma la textura y limita: $\widehat F_a=\sat_{[-E,E]}(F_a+T(t))$. No existe cuantización intermedia con escala 10000 ni corte por envolvente codificada cero.

La rampa se aplica una sola vez dentro del modelo. El modelo de aceleraciones se consume una vez por revisión del receptor y se mantiene entre tramas, sin triangular ni interpolación. La textura se actualiza en cada ciclo de motores. El paso cero no borra los filtros. Al detenerse hasta 0.5 kt o abandonar suelo, las fuerzas se anulan en la siguiente muestra válida, junto con la textura.

El receptor valida CRC, secuencia, banderas, datos y avance de marca. Exige trama armada con menos de 300 ms y DTR. El puente desarma por pausa, simulación detenida, slew o pérdida de datos; exige datos de menos de 250 ms para armar. Permanecen rutas históricas para SC1--SC7, pero el puente actual sólo emite SC8.

\subsection{Paso 6: ganancias, límites y suma con otros efectos}
La etapa de suelo en \code{Effects::command} aplica:
\begin{align}
g_R&=\sat_{[-L_g,L_g]}(\widehat F_R c_g G),\\
g_P&=\sat_{[-L_g,L_g]}(\widehat F_P c_g G),\\
u_R&=g_R, & u_P&=g_P/10, & L_g&=40/255.
\label{eq:ground-gains}
\end{align}
Hay \textbf{dos señales y dos recortes}, seguidos de la reducción adicional de pitch. El divisor 10 limita su sensibilidad, pero ya no significa que pitch sea siempre un décimo de roll. Los máximos aislados con escala PWM 255 siguen siendo 40 cuentas roll y 4 pitch.

Con ganancias iniciales $c_gG=0.18\times0.90=0.162$, un eje se recorta cuando $|\widehat F_a|$ supera aproximadamente 0.968289. No existe una velocidad de saturación universal: también depende de la perturbación filtrada. Sólo en el caso extremo $|s_a|=1$, ignorando cuantización, se obtiene la referencia $V_g\simeq38.7474$ kt para comenzar ese recorte.

El suelo se suma después del damper y la fricción; esas resistencias no multiplican su señal. Para cada eje, con $u_{{\rm otros},a}$ como suma de los demás aportes:
\begin{align}
u_{{\rm total},a}&=\sat_{[-220/255,220/255]}(u_{{\rm otros},a}+u_a),\\
p_a&=\sat_{[-220,220]}\left(\operatorname{round}_{\rm sim}(255\,d_a u_{{\rm total},a})\right).
\end{align}
$d_a$ es la polaridad de motor, inicialmente $+1$ para ambos ejes. El redondeo es simétrico, con mitades alejándose de cero. Ver más de 40 PWM durante rodaje no contradice el techo propio: pueden sumarse centrado u otras cargas.

PWM positivo activa RIGHT y anula LEFT; negativo hace lo contrario. Roll usa GP6/GP7, pitch GP8/GP9, con habilitaciones GP10/GP11. La frecuencia eléctrica configurada sigue siendo 490 Hz; \textbf{la componente de baches no tiene frecuencia impuesta; la textura utiliza 7 y 11.3 Hz}. Un PWM 40 equivale a $40/255\simeq15.69\%$ de duty, no a 40 N ni a una corriente medida. Si \code{PWM_MAX} cambia, la conversión usa esa escala en lugar de 255, mientras el techo normalizado sigue referido a 255.

\subsection{Ejemplos numéricos de la nueva cadena}
\subsubsection{Una variación vertical a 20.25 kt}
Supuestos: filtros inicializados y sin transitorio previo, entrada lateral sin cambio, incremento vertical de 8 ft/s$^2$, $\Delta t=20$ ms, avión en suelo, tramas válidas y ganancias predeterminadas. Los valores redondeados son:
\begin{enumerate}
\item La velocidad da $E=0.5$.
\item El pasaaltos vertical produce $h_Y=0.726221\times8\simeq5.809768$ ft/s$^2$.
\item El pasabajos produce $\ell_Y=0.5h_Y\simeq2.904884$ ft/s$^2$; $\ell_X=0$.
\item La mezcla da $s_R\simeq0.252802$ y $s_P\simeq0.361146$.
\item La Pico obtiene $F_R\simeq0.126401$ y $F_P\simeq0.180573$.
\item Estos valores se usan directamente en coma flotante, sin codificarlos de nuevo por USB.
\item Roll: $u_R\simeq0.126401\times0.162=0.02047696$; pitch: $u_P\simeq0.180573\times0.162/10=0.00292528$.
\item El aporte aislado redondea a 5 PWM roll y 1 PWM pitch.
\end{enumerate}
No se consultó una fase triangular. Si la aceleración queda en el nuevo nivel, los filtros evolucionan y finalmente las señales tienden a cero. Si el cambio inicial fuera negativo, las señales tendrían signo opuesto bajo estos mismos supuestos.

\subsubsection{Misma señal filtrada a distintas velocidades}
Para comparar sólo la rampa, se fijan hipotéticamente $s_R=0.6$ y $s_P=0.4$. Se incluyen valores redondeados de fuerza y redondeo final, con ganancias iniciales y sin otros efectos:
\begin{center}
\begin{tabular}{rrrrrr}
\toprule
$V_g$ (kt) & $E$ & $\widehat F_R$ & $\widehat F_P$ & PWM roll & PWM pitch\\\midrule
0.5 & 0 & 0 & 0 & 0 & 0\\
10 & 0.240506 & 0.1443 & 0.0962 & 6 & 0\\
20.25 & 0.500000 & 0.3000 & 0.2000 & 12 & 1\\
40 & 1.000000 & 0.6000 & 0.4000 & 25 & 2\\
60 & 1.000000 & 0.6000 & 0.4000 & 25 & 2\\
\bottomrule
\end{tabular}
\end{center}
La tabla no predice la aceleración de una pista real. Si $s_R=s_P=0$, toda la tabla de fuerzas sería cero aunque $E=1$. Si ambos alcanzaran magnitud uno, los máximos a 40 kt quedarían recortados a 40 y 4 PWM. Una señal lateral y una vertical opuestas pueden incluso producir signos distintos entre ejes.

\subsection{Parada, pausa, vuelo y reinicios}
Al parar hasta 0.5 kt, salir del suelo, entrar en pausa o slew, perder continuidad o recibir datos particulares inválidos, el modelo devuelve cero y reinicia sus filtros. No conserva la cola del pasabajos al volver a una condición inactiva. La primera muestra válida de la reanudación establece una nueva referencia sin golpe inicial.

En la placa, la pérdida de habilitación, el primer update activo o un intervalo local mayor que 50 ms borran \code{groundSignal}. En el siguiente update normal puede volver a copiarse la última señal fresca del receptor: no hay una rampa local de arranque de 100 ms. El watchdog de 300 ms, el desarme y DTR completan la retirada de fuerza. Los tiempos describen procesamiento y comunicación, no latencia eléctrica absolutamente nula.

Se mantienen las banderas \code{SIMCONNECT_ADDITIONS_ENABLED} y \code{SIMCONNECT_GROUND_VIBRATION_ENABLED}, la máscara \code{Ground}, el menú y las ganancias persistentes. Mientras el menú está activo se inhiben los motores; guardar configuración también los detiene. La ganancia e interruptor de suelo siguen siendo comunes a ambos ejes.

\subsection{Relación con TelemFFB y procedencia de las ecuaciones}
En la copia local de \code{MotionEffectsMixIn.py}, \code{ac_update_runway_rumble} procesa compresión de rueda central y diferencia izquierda--derecha:
\begin{align}
v_1&=\operatorname{HPF}(W_0)I, & v_2&=\operatorname{HPF}(W_1-W_2)I.
\end{align}
Los pasaaltos tienen corte configurado de 3 Hz; los resultados se limitan a $\pm0.5$ y se aplican como efectos de fuerza con dirección modulada aleatoriamente y envolvente de ataque. BMS utiliza una ruta distinta de telemetría de golpes.

SC8 adopta la idea de responder a \emph{variaciones de telemetría}, pero usa aceleraciones del avión en Prepar3D. No reproduce compresión individual, dirección aleatoria ni envolvente de ataque de TelemFFB. Tampoco puede garantizar sensación equivalente entre mecanismos. La rampa 0.5--40 kt, mezcla 0.3/0.7, escala 0.25 g, pasabajos de 20 ms y reparto de ganancias son decisiones locales.

El filtro pasaaltos y el pasabajos son discretizaciones de primer orden. La normalización y la saturación convierten la señal a un comando limitado; no constituyen identificación de rigidez, amortiguación o fuerza del tren. El accionamiento con dos PWM según signo conserva el antecedente del proyecto alemán original, pero no se atribuye a ese original este modelo de rodaje por SimConnect.

\subsection{Ajustes y diagnóstico}
``Vib. suelo'' ajusta $c_g$ de 0 a 1, inicialmente 0.18, en pasos de 0.01. ``Ganancia total'' ajusta $G$, inicialmente 0.90, y afecta todos los efectos. El menú no ofrece ganancias de suelo separadas por eje. Modificar un valor inicial del código no reemplaza necesariamente una configuración persistida válida.

Para una muestra sin recorte, $|p_R|$ antes de redondear es $255Gc_g|\widehat F_R|$. Una estimación de ganancia para un comando deseado $P_R$ es:
\begin{equation}
c_g\simeq\frac{P_R}{255G|\widehat F_R|},\qquad |\widehat F_R|>0.
\end{equation}
Ahora se necesita conocer la perturbación, no sólo la velocidad. Con $G=0.90$ y $|\widehat F_R|=0.3$, un objetivo de 10 cuentas da $c_g\simeq0.14524$. No es una calibración de fuerza física. Para evitar el recorte propio incluso con señal unitaria debe cumplirse $c_gG\le40/255$.

Los filtros, mezcla y escala se modifican en \code{GroundModel.h} y requieren recompilar el firmware. El divisor \code{SIMCONNECT_GROUND_ROLL_PITCH_RATIO} y techo \code{SIMCONNECT_GROUND_MAX_PWM} están en el firmware. El literal 8 que aún figura en el arreglo de frecuencias sólo gobierna la compatibilidad con suelo anterior a SC7; cambiarlo no altera esta nueva señal.

\subsubsection{CSV y límites de observación}
El CSV actual registra \code{ground_speed_knots}, \code{ground_accel_x_ft_s2}, \code{ground_accel_y_ft_s2}, \code{on_ground}, \code{sample_time_ms} y los estados de envío. No contiene \code{ground_amp}, \code{ground_roll}, \code{ground_pitch} ni \code{Vibracion_Tierra}: esas magnitudes ya no se calculan en Windows.

En \code{pico.log}, \code{gr}/\code{gp} son las señales de suelo de la Pico por 10000, mientras \code{R}/\code{P} son PWM totales aplicados. Se debe aislar el efecto para compararlos directamente. LCD y diagnóstico son muestras cercanas a 100 ms, pueden omitir impulsos breves y no miden fuerza física.

\subsubsection{Procedimiento práctico de verificación}
\begin{enumerate}
\item Verificar firmware \code{YOKE17SC8} y puente nuevo. El lanzador apunta a \code{SimConnectBridge/bin/Win32/RealismoSC8}.
\item En monitor, registrar aceleraciones y velocidad; las fuerzas se observan en el diagnóstico de la Pico durante una prueba habilitada. Para registrar sin ordenar motores puede usarse \code{YokeSimConnectBridge.exe --seconds 60 --no-buffets} sin puerto ni \code{--enable}.
\item Comprobar cero hasta 0.5 kt y al despegar; contrastar pista lisa y terreno irregular por sus señales, no sólo por el nombre de superficie.
\item Para evaluar el mecanismo, aislar ``Vib. suelo'' mediante el menú con pausa reconocida y usar ganancia reducida. Comparar registros de comando con la sensación.
\item Confirmar retirada en pausa y pérdida de comunicación. A partir de 40 kt la envolvente debe permanecer en uno, pero las sacudidas pueden seguir cambiando.
\end{enumerate}
Si no se siente vibración, comprobar velocidad de suelo, habilitación, ganancias y PWM. La textura debe aportar señal aunque las aceleraciones sean constantes; las sacudidas adicionales dependen de sus variaciones. El rozamiento mecánico puede ocultar comandos pequeños. Si hay fuerza parado, revisar otros efectos y la velocidad informada. No se debe corregir la calibración de posición para ajustar este modelo: la posición no interviene en él.

\subsection{Validación, archivos y límites del alcance}
Se compilaron el puente Win32 y el firmware Pico 2 ARM/TinyUSB para potenciómetros. La suite con reproducción del CSV ejecutó 96\,764 comprobaciones correctas. El firmware ocupa 99\,748 bytes de programa y 18\,408 bytes de RAM global. Con ganancia de suelo 0.450 y total 90\%, la reproducción generó 1129 paquetes, 800 muestras con PWM de suelo no nulo y picos roll/pitch de 40/4. Son comandos calculados por software, no mediciones de hardware. \textbf{No se cargó el firmware ni se comprobó físicamente el rodaje o el tiempo de ejecución en la Pico.}

Las pruebas cubren rampa, mitad de señal a 20.25 kt respecto de 40 kt para las mismas entradas, saturación del multiplicador, respuesta a cambios y rechazo de carga constante, reinicios, signo por eje, CRC, rangos, límites de comando, textura de pista con aceleración constante, watchdog, pausa, desconexión y recepción de protocolos anteriores. No demuestran una equivalencia perceptual con TelemFFB.

\begin{longtable}{p{0.40\linewidth}p{0.52\linewidth}}
\toprule Archivo & Elementos a consultar\\\midrule\endhead
\code{GroundModel.h} & Rampa, validación, filtros, mezcla y fuerzas con signo.\\
\code{SimConnectBridge/main.cpp} & 22 variables crudas, selección de muestras, trama SC8 y CSV.\\
\code{SimTelemetry.h} & Datos crudos SC8, validación de reloj, compatibilidad y watchdog.\\
\code{SimEffects.h} & Selección SC8, ganancias, recortes por eje y suma.\\
\code{mis_defines.h} & Ganancia inicial, divisor, límites y banderas.\\
\code{GainSettings.h} & Menú y valores efectivos persistentes.\\
\code{MiYoke_PD_V2.4.ino} & Identificación SC8, ciclo, inhibiciones y salida.\\
\code{MotorDriver.h} & Redondeo y accionamiento PWM.\\
\code{tests/ground_checks.h} & Pruebas del modelo y de compatibilidad SC7; SC8 en \code{tests/raw_checks.h}.\\
\code{tests/simconnect_checks.cpp} & Invocación y regresión del resto de efectos.\\
\bottomrule
\end{longtable}
Las rutas son relativas a MiYoke\_PD\_V2.4. \code{LEEME md/CALCULO_PICO_SC8.md} documenta la puesta en uso; \code{SUELO_REALISTA_SC7.md} es histórico. Las explicaciones antiguas de triangular de suelo describen protocolos anteriores, no la componente de aceleraciones SC8.

Quedan como mejoras futuras una fuente válida de compresión por rueda, un efecto separado de contacto, histéresis de velocidad, ganancias independientes y calibración de fuerza. Se incorporó textura determinista con dos senoides; no se incorporaron textura aleatoria, eventos por distancia ni clasificación de superficie. La respuesta disponible depende del modelo de vuelo, del muestreo y de la sensibilidad elegida.
\newpage
\section{BUFFET POR TREN DE ATERRIZAJE}
\label{sec:buffet-tren}

\subsection{Qué hace este efecto}
Genera una vibración periódica de roll y pitch cuando el avión está en vuelo, tiene tren retráctil y su tren izquierdo está al menos parcialmente desplegado. La intensidad crece con la extensión y con la velocidad indicada, entre 100 y 150 kt. La frecuencia actual es 10 Hz y la onda es triangular.

La idea sencilla es: \emph{``representar una vibración aerodinámica asociada al tren expuesto al flujo de aire''}. El simulador no envía esa vibración instantánea: informa posición de tren y velocidad; la Pico calcula tanto la amplitud como la oscilación.

\subsubsection{Qué no está incluido}
No representa impactos de ruedas, rodaje, ruido del motor de extensión, vibración del actuador ni golpe de enclavado. Tampoco calcula fuerza de arrastre ni torque de bisagra. El buffet puede continuar con el tren inmóvil y extendido: no necesita que el tren esté moviéndose.

La entrada observada es únicamente el tren izquierdo. No se promedian tren derecho y de nariz, ni se verifica que los tres estén trabados. Un avión de tren fijo queda excluido por la condición \code{IS GEAR RETRACTABLE=1}, aunque en la realidad pueda transmitir vibraciones. Esa exclusión es una decisión del modelo.

\subsection{Recorrido y variables de SimConnect}
\subsubsection{Cadena de cálculo}
\begin{center}
\fbox{\begin{minipage}{0.91\linewidth}
IAS y extensión de tren $\rightarrow$ verificar vuelo y tren retráctil $\rightarrow$ rampa 100--150 kt $\times$ extensión $\rightarrow$ envolvente $0\ldots1$ $\rightarrow$ filtro local de 100 ms $\rightarrow$ triangular de 10 Hz $\rightarrow$ ganancia del efecto y ganancia total $\rightarrow$ roll recibe la mitad de pitch $\rightarrow$ límite general $\rightarrow$ PWM.
\end{minipage}}
\end{center}

\subsubsection{Variables que forman o habilitan el efecto}
\begin{longtable}{p{0.32\linewidth}p{0.18\linewidth}p{0.42\linewidth}}
\toprule Variable & Unidad solicitada & Uso\\\midrule\endhead
\code{AIRSPEED INDICATED} & knots & Velocidad $V$ de la rampa de intensidad.\\
\code{GEAR LEFT POSITION} & percent over 100 & Extensión $d$ del tren izquierdo: 0 retraído, 1 extendido.\\
\code{IS GEAR RETRACTABLE} & Bool & Debe ser exactamente 1.\\
\code{SIM ON GROUND} & Bool & Debe ser 0: sólo se calcula objetivo en vuelo.\\
\code{IS SLEW ACTIVE} & Bool & Debe ser 0 para habilitar.\\
\bottomrule
\end{longtable}
Los nombres y unidades se contrastaron con el \href{https://www.prepar3d.com/SDKv3/LearningCenter/utilities/variables/simulation_variables.html}{SDK de Prepar3D}. Se usa posición real publicada del tren, no \code{GEAR HANDLE POSITION}. Los eventos \code{Sim}/\code{Pause}, la frescura y el armado completan la habilitación.

\subsubsection{Variables que no intervienen en la amplitud}
No se usa presión dinámica, velocidad sobre suelo, viento, AoA, flaps, posición de yugo ni velocidad del mando. \code{--reference-knots} y los argumentos de ráfagas no cambian esta ley. La validación compartida puede desarmar el paquete si otras variables son inválidas; esa dependencia de integridad no significa que sean factores de intensidad.

\subsection{Paso 1: fórmula de amplitud en la Pico}
\subsubsection{Rampa y multiplicación por despliegue}
Definimos $\sat_{[0,1]}(x)=\min(1,\max(0,x))$. En \code{BuffetModel.h::calculate}:
\begin{align}
R_g(V)&=\sat_{[0,1]}\left(\frac{V-100}{150-100}\right),\\
E_g&=d\,R_g(V).
\label{eq:gear-env}
\end{align}
$V$ es IAS en kt; $d$ es extensión fraccionaria; $R_g$ es una rampa adimensional. El 100 resta la velocidad inicial de vibración; $150-100=50$ kt es el ancho de la transición. El resultado $E_g$ es una envolvente no negativa, todavía sin signo alternante.

Hasta 100 kt, $R_g=0$; a 125 kt vale 0.5; desde 150 kt vale 1. Con el tren a mitad, $d=0.5$, la envolvente es la mitad. Así, a 125 kt y media extensión, $E_g=0.25$. No existe un umbral adicional de despliegue del 10\% en MiYoke: cualquier $d>0$ válido puede generar señal si $V>100$, sujeto a cuantización.

\subsubsection{Procedencia y límites de los umbrales}
Los 100 y 150 kt son ajustes locales de sensación. No se solicitan Vle ni Vlo al simulador y no se deducen estos umbrales de los límites operativos del avión. Cambiar de aeronave no adapta automáticamente la rampa. Una ausencia de buffet bajo 100 kt es una consecuencia del código, no una conclusión sobre lo que vibra en un avión real.

\subsubsection{Validación específica}
El modelo acepta IAS finita entre 0 y 1000 kt. Fuera de ese rango devuelve cero para los tres buffets. Tren requiere $d$ finito en $[0,1]$ y retráctil exactamente igual a 1. Un despliegue 1.2 no se interpreta como totalmente desplegado: deshabilita este objetivo. La ausencia o invalidez de datos de pérdida o spoiler puede anular esos efectos sin anular un tren válido.

\subsection{Paso 2: obtención local de la envolvente}
El puente transmite IAS, extensión de tren, condición retráctil y estado de suelo. La Pico calcula $E_g$ en \code{BuffetModel.h}, lo guarda en \code{PicoForceModels::buffet[1]} y lo usa directamente como $E_{g,\rm rx}$. No transmite un entero de amplitud. El objetivo se actualiza por nueva muestra, hasta unos 50 Hz; filtro y portadora avanzan con el ciclo local nominal de 2 ms.

\subsection{Paso 3: filtro y portadora}
\subsubsection{Suavizado de 100 ms}
En la placa, con $\Delta t$ en ms:
\begin{align}
\beta_n&=\frac{\Delta t_n}{100+\Delta t_n},\\
A_{g,n}&=A_{g,n-1}+\beta_n(E_{g,\rm rx,n}-A_{g,n-1}).
\end{align}
$A_g$ es la amplitud suavizada. Con 2 ms, $\beta=2/102$. Es una discretización del filtro de primer orden $\tau\dot A+A=E$, con $\tau=100$ ms. Desde cero y objetivo constante, alcanza aproximadamente 62.85\% a 100 ms, 94.87\% a 300 ms y 99.29\% a 500 ms.

El número 100 se comparte con otras envolventes y el centrado dentro de \code{Effects::update}. Modificarlo directamente cambia más de un efecto. No es un tiempo físico medido del tren.

\subsubsection{Onda triangular y fase}
El generador usa
\begin{align}
\phi_{g,n}&=(\phi_{g,n-1}+\Delta t_{\rm ms}\,10)\bmod1000,\\
T(\phi)&=\begin{cases}0.004\phi-1,&\phi<500,\\3-0.004\phi,&\phi\ge500,\end{cases}\\
w_g&=A_gT(\phi_g).
\end{align}
La fase es una cuenta de milésimas de ciclo. $10$ Hz implica un período de 100 ms. El factor 0.004 es $2/500$: pendiente de la recta que lleva la triangular de $-1$ a $+1$. En fases 0, 250, 500 y 750, $T$ vale $-1$, 0, $+1$ y 0.

Al reiniciar se usa fase 250, un cruce por cero. Durante operación habilitada la fase sigue avanzando incluso con objetivo cero. Retraer y volver a extender el tren no reinicia necesariamente la fase; perder habilitación general o sufrir un update de más de 50 ms sí reinicia amplitud y fase. No hay aleatoriedad en la onda.

\subsubsection{Cola al retirar el tren o tocar tierra}
Cuando llega $E_g=0$, el filtro mantiene inicialmente parte de $A_g$ y la reduce suavemente. Con ciclo de 2 ms y sin reinicio, $A_{g,n}=A_{g,0}(100/102)^n$. No existe para los buffets la rama de corte instantáneo que sí tiene vibración de suelo.

Al tocar tierra, la Pico pone los tres objetivos de buffet en cero pero puede seguir armado para suelo y centrado. Puede haber una breve coexistencia entre cola de buffet y comienzo de rodaje. En cambio, pausa, desarme, DTR perdido o watchdog vencido deshabilitan efectos y limpian el estado. Es importante distinguir objetivo cero de pérdida de habilitación.

\subsection{Paso 4: comandos de roll y pitch}
Con $c_g=0.16$ como ganancia ``Buffet tren'' y $G=0.90$ como ganancia total:
\begin{align}
u_{P,g}&=G c_g w_g=0.144 A_gT(\phi_g),\\
u_{R,g}&=0.5G c_g w_g=0.072 A_gT(\phi_g).
\end{align}
Pitch recibe toda la amplitud y roll la mitad, con la misma fase. La relación 1:2 corresponde a comando antes del redondeo, no a fuerza física. No son las ganancias ``Centrado roll/pitch'' ni ``Turbul. roll/pitch'': cada una pertenece a otra ruta.

No se utiliza el desplazamiento del yugo y tampoco la zona muerta de 0.02 de turbulencias. Las señales pequeñas de buffet llegan a la conversión PWM y pueden desaparecer por redondeo. Damper y fricción pueden sumarse al resultado, pero no multiplican la portadora de tren.

\subsection{Paso 5: de comando normalizado a PWM}
\subsubsection{Conversión completa}
Con el tren aislado, se aplica
\begin{align}
u_{{\rm out},a}&=\sat_{[-220/255,220/255]}(u_{a,g}),\\
p_a&=\sat_{[-220,220]}\left[\operatorname{round}_{\rm sim}\!left(255\,
\sat_{[-1,1]}(s_a u_{{\rm out},a})\right)\right].
\end{align}
$s_a$ es polaridad de motor, inicialmente $+1$. El redondeo suma 0.5 a positivos y resta 0.5 a negativos antes de convertir a entero. El driver multiplica por \code{PWM_MAX}, inicialmente 255; el 220 es techo posterior, no el multiplicador.

Con ganancias iniciales los picos ideales son $0.144\times255=36.72$ cuentas en pitch y 18.36 en roll: redondean a \textbf{37 y 18}. No hay techo propio de 40 como en suelo. El límite general pertenece a la suma de todos los efectos.

\subsubsection{Pines y significado eléctrico}
PWM positivo deja LEFT en cero y activa RIGHT; negativo hace lo contrario. En Pico, roll usa GP6/GP7 y pitch GP8/GP9; enable GP10/GP11. La escala eléctrica es 0--255 y la frecuencia PWM configurada 490 Hz. Por ejemplo, 18 cuentas son 7.06\% de duty; no son 18 N ni 18\% de torque medido. La portadora de 10 Hz modula el comando sobre ese PWM.

\subsection{Ejemplos numéricos}
\subsubsection{Tren extendido a 125 kt}
En vuelo, con tren retráctil y $d=1$: $E_g=(125-100)/50=0.5$. La Pico calcula 0.5 sin cuantización intermedia. En régimen $A_g\simeq0.5$ y, en una fase de ejemplo con $T=0.8$:
\begin{align*}
w_g&=0.4, &u_{P,g}&=0.0576,&p_P&=15,\\
&&u_{R,g}&=0.0288,&p_R&=7.
\end{align*}
Se obtiene $p_P=\operatorname{round}(14.688)=15$ y $p_R=\operatorname{round}(7.344)=7$. Si $T=-0.8$, se invierten los signos. Si el tren estuviera a mitad, la amplitud previa al redondeo se reduciría a la mitad.

\subsubsection{Tabla de picos ideales a extensión completa}
Supuestos: filtro estabilizado, $|T|=1$, ganancias iniciales y demás efectos nulos.
\begin{center}\begin{tabular}{rrrr}
\toprule IAS (kt)&$E_g$&Pico PWM roll&Pico PWM pitch\\\midrule
100&0&0&0\\110&0.2&4&7\\120&0.4&7&15\\125&0.5&9&18\\140&0.8&15&29\\150&1&18&37\\180&1&18&37\\\bottomrule
\end{tabular}\end{center}
Son picos de la triangular ideal. Con ciclos exactamente de 2 ms y fase inicial 250, el incremento 20 no visita las fases 0 y 500: $|T|$ llega a 0.96. Así, a amplitud uno los extremos muestreados equivalen a unos 35.25 PWM pitch y 17.63 roll, que redondean a 35 y 18. Un lazo con otros intervalos puede visitar fases distintas. Esto explica diferencias pequeñas entre el pico ideal y el registro real.

\subsubsection{Arranque con objetivo 0.5}
Desde $A_g=0$, fase 250 y pasos de 2 ms, con un primer update de inicialización ya completado:
\begin{center}\begin{tabular}{rrrrrr}
\toprule ms&Fase&$A_g$&$T$&PWM roll&PWM pitch\\\midrule
2&270&0.009804&0.08&0&0\\10&350&0.047135&0.40&0&1\\20&450&0.089826&0.80&1&3\\24&490&0.105753&0.96&2&4\\74&990&0.259695&$-0.96$&$-5$&$-9$\\100&250&0.314236&0&0&0\\300&250&0.474358&0&0&0\\\bottomrule
\end{tabular}\end{center}
La última fila no significa que no haya buffet: la muestra cae en cruce por cero aunque la amplitud ya sea alta.

\subsection{Habilitación y diagnóstico}
\subsubsection{Banderas y estados}
Los buffets están activados por defecto en el puente. \code{--no-buffets} anula los tres objetivos; \code{--gusts} no es necesario. El cálculo requiere vuelo, simulación en marcha, ausencia de pausa/slew e intervalo de datos menor que 250 ms. En monitor se registran entradas; no se calcula amplitud en Windows.

La placa exige habilitación general, \code{SIMCONNECT_GEAR_BUFFET_ENABLED}, bit \code{Gear}, ganancia efectiva, DTR y trama armada de menos de 300 ms. El menú fuerza comandos a cero. Un update de efectos mayor que 50 ms reinicia ondas. Cambiar el interruptor de efecto a apagado hace efectiva ganancia cero; el objetivo puede seguir calculándose en la Pico, pero su aporte se hace cero.

\subsubsection{CSV, LCD y aliasing de 10 Hz}
Registrar IAS, extensión de tren, condición retráctil, estado de suelo y armado de las filas enviadas. SC8 no calcula ni registra \code{gear_amp} o \code{Buffet_Tren} en Windows; su objetivo se obtiene en la Pico.

La pantalla se actualiza cada 100 ms, el mismo período de la portadora de 10 Hz. Si ambas quedan sincronizadas, puede observar siempre la misma fase, incluso siempre cero. Las pruebas del proyecto muestran señal de motor variable y muestras LCD nulas en ese caso. \textbf{No debe bajarse la frecuencia del efecto sólo para hacerlo visible en el LCD}: conviene mejorar el diagnóstico o registrar máximos/RMS en una ventana.

\subsection{Procedencia y comparación con las referencias}
\subsubsection{TelemFFB}
En la copia local de \code{MotionEffectsMixIn.py::ac_update_landing_gear}, TelemFFB relaciona despliegue con una rampa de velocidad y usa senoides en dos direcciones, con frecuencia predeterminada 10 Hz. También contempla efectos de movimiento y enclavado y, para X-Plane, umbrales relacionados con Vle. MiYoke conserva la idea de buffet por despliegue/velocidad, pero usa una sola triangular y reparto roll/pitch 1:2.

Los literales 100/150 de la clase de referencia se comparan con velocidad en m/s; no son una demostración de que el original use 100/150 kt. En MiYoke se eligieron explícitamente kt. Tampoco se copió el corte de despliegue mayor que 0.1 de esa ruta. La restricción local a tren retráctil y la lectura sólo del izquierdo deben evaluarse como decisiones propias, no como reproducción completa de TelemFFB.

Referencia de localización: \href{https://github.com/walmis/VPforce-TelemFFB/blob/refactor/telemffb/sim/base/MotionEffectsMixIn.py}{TelemFFB, MotionEffectsMixIn.py}. La comparación se hizo contra la copia incluida en el proyecto.

\subsubsection{Original alemán y decisiones locales}
No se identificó en las fuentes del original local una cadena SimConnect que produzca este buffet con esta rampa. El esquema de dos PWM tiene un antecedente en \code{AxisCalibration.cpp::MoveMotor}, pero la envolvente, transporte, triangular, ganancia 0.16, factor roll 0.5 y filtro de 100 ms son de la ruta actual. No se atribuyen a una ley aerodinámica ni a una calibración física medida.

\subsection{Ajustes, calibración y mejoras posibles}
\subsubsection{Qué se puede ajustar ahora}
``Buffet tren'' admite 0--1 con paso 0.01; inicialmente 0.16. ``Ganancia total'' inicialmente 0.90 escala también otros efectos. El menú no ofrece ganancias separadas de buffet tren por eje: roll es la mitad de pitch en código. Los valores persistidos pueden reemplazar los defaults.

\code{SIMCONNECT_GEAR_BUFFET_HZ} vale actualmente 10, admite enteros de 1 a 20 y requiere recompilar firmware. Los umbrales 100/150 se editan en \code{BuffetModel.h} y requieren recompilar el firmware. Cambiar frecuencia altera el carácter y la respuesta mecánica, no la envolvente de velocidad. Cambiar el centro de potenciómetro no ajusta este efecto.

\subsubsection{Cómo elegir una ganancia}
En régimen sin saturar, un pico pitch deseado $P$ conduce a $c_g\simeq P/(255GE_g)$ para $E_g>0$. Buscar 10 PWM pitch con $E_g=0.5$ y $G=0.90$ da $c_g\simeq0.08715$, aproximable a 0.09 desde el menú. Es una cota ideal de pico; muestreo y transmisión mecánica pueden reducir la sensación.

La primera cuenta ideal aparece aproximadamente si $E_g\ge0.5/(255Gc_g)$ en pitch y el doble en roll. Con ganancias iniciales son 0.01362 y 0.02723. Señal válida y bandera encendida no garantizan vibración perceptible.

\subsubsection{Ensayo y mejoras propuestas}
Comparar primero el tren aislado: confirmar tipo retráctil y posición real en CSV, deshabilitar otros efectos desde el menú con pausa reconocida, variar despliegue y comprobar la rampa en condiciones del simulador. Los 100/150 kt no son instrucciones operativas para un avión real. Observar varios ciclos y usar un registro suficientemente rápido.

Se puede mejorar la adaptación mediante perfiles por aeronave, lecturas de todos los trenes, separación entre movimiento/impacto y buffet, ganancias por eje y diagnóstico de picos. Son propuestas; no están implementadas aquí. Si se desea usar Vle como referencia, debe definirse conscientemente la relación con la sensación, no suponer que Vle es por sí sola el inicio físico del buffet.

\subsection{Fuentes y alcance de la comprobación}
La ruta verificable es \code{BuffetModel.h::calculate} $\rightarrow$ \code{main.cpp::Send} $\rightarrow$ \code{SimTelemetry.h::Receiver::feed} $\rightarrow$ \code{SimEffects.h::update/command} $\rightarrow$ sketch $\rightarrow$ \code{MotorDriver.h}. Ganancias y frecuencia están en \code{mis_defines.h}/\code{GainSettings.h}. \code{tests/simconnect_checks.cpp} comprueba rampa, habilitación, proporción por eje y submuestreo del LCD.

\code{LEEME md/EFECTOS_SIMCONNECT.md} conserva valores históricos como frecuencia de tren 2 Hz y ganancia 0.025. La fuente vigente usa 10 Hz y 0.16, que son los valores de esta sección. Las referencias de variables son el SDK enlazado; los ejemplos son cálculo de software, no mediciones del yugo ni verificación del firmware cargado.
\newpage
\section{BUFFET POR PERDIDA}
\label{sec:buffet-perdida}

\subsection{Qué representa el efecto}
Genera una vibración de roll y pitch al aumentar el ángulo de ataque dentro de una banda previa al ángulo de pérdida informado por el avión. La intensidad depende además de flaps e IAS comparada con una velocidad de referencia. La placa sintetiza una triangular de 13 Hz y pitch recibe el doble de comando que roll.

La idea sencilla es: \emph{``cuanto más se aproxima el ángulo de ataque al umbral del modelo, aumentar el temblor, ponderándolo por el flujo que estima la velocidad''}. La vibración puede comenzar antes de alcanzar \code{STALL ALPHA}; no es un interruptor que espere a que la pérdida ya se haya producido.

\subsubsection{Qué no calcula}
No calcula caída de ala, movimiento de morro, pérdida de sustentación, un empuje de recuperación ni un stick shaker certificado. No utiliza el estado de la bocina de pérdida. No reconstruye esfuerzo de bisagra o vibración estructural en unidades físicas y no demuestra que una aeronave real comience a vibrar a la fracción de AoA elegida.

El modelo produce buffet de ángulo positivo. No usa valor absoluto de AoA: no representa simétricamente una pérdida con ángulo negativo. Una velocidad baja por sí sola tampoco activa vibración si el ángulo sigue por debajo de la banda.

\subsection{Recorrido desde SimConnect}
\begin{center}
\fbox{\begin{minipage}{0.91\linewidth}
AoA, ángulo de pérdida, flaps, VS0/VS1 e IAS $\rightarrow$ validar datos y vuelo $\rightarrow$ promedio de flaps $\rightarrow$ ángulo de comienzo y velocidad de referencia $\rightarrow$ rampa de AoA $\times$ factor de velocidad al cuadrado $\rightarrow$ envolvente limitada $0\ldots1$ $\rightarrow$ filtro local 100 ms $\rightarrow$ triangular 13 Hz $\rightarrow$ ganancias y reparto roll/pitch $\rightarrow$ límite general $\rightarrow$ PWM.
\end{minipage}}
\end{center}

\subsection{Variables, unidades y validez}
\subsubsection{Entradas directas}
\begin{longtable}{p{0.42\linewidth}p{0.16\linewidth}p{0.34\linewidth}}
\toprule Variable & Unidad solicitada & Símbolo/uso\\\midrule\endhead
\code{AIRSPEED INDICATED} & knots & $V$: factor de flujo.\\
\code{INCIDENCE ALPHA} & degrees & $\alpha$: ángulo de ataque.\\
\code{STALL ALPHA} & degrees & $\alpha_s$: extremo de la banda.\\
\code{TRAILING EDGE FLAPS LEFT PERCENT} & percent over 100 & $f_L$: flap izquierdo, 0--1.\\
\code{TRAILING EDGE FLAPS RIGHT PERCENT} & percent over 100 & $f_R$: flap derecho, 0--1.\\
\code{DESIGN SPEED VS0} & knots & $V_{s0}$: referencia de diseño VS0.\\
\code{DESIGN SPEED VS1} & knots & $V_{s1}$: referencia de diseño VS1.\\
\bottomrule
\end{longtable}
En el modelo, VS0 se usa como extremo de referencia con flap completo y VS1 como extremo limpio. No se recalculan con peso, factor de carga, hielo u otros estados. Son datos de diseño publicados por el avión, no una medición online de la velocidad de pérdida real en todas las condiciones.

\subsubsection{Conversión de unidades solicitada}
El \href{https://www.prepar3d.com/SDKv3/LearningCenter/utilities/variables/simulation_variables.html}{SDK de Prepar3D} publica ángulos como radianes y velocidades de diseño como pies/s. \code{main.cpp} pide explícitamente grados y knots en \code{SimConnect_AddToDataDefinition}: SimConnect entrega los datos en las unidades solicitadas. El modelo no vuelve a convertir radianes ni pies/s.

``Percent over 100'' se interpreta aquí como fracción 0--1: 0.5 equivale a 50\% de despliegue. No se introducen valores 50 en la fórmula. Los flaps leídos son las superficies de borde de salida, no la palanca ni un número de posición discreto.

\subsubsection{Rangos de sanidad}
Para calcular pérdida se requieren todos estos datos válidos:
\begin{center}\begin{tabular}{ll}
\toprule Dato & Rango admitido\\\midrule
IAS & 0--1000 kt\\
AoA & $-90$ a $+90$ grados\\
Ángulo de pérdida & 1--45 grados\\
Cada flap & 0--1\\
VS0 y VS1 & 10--400 kt cada una\\\bottomrule
\end{tabular}\end{center}
Todos deben ser finitos. Si el avión devuelve \code{STALL ALPHA=0}, se anula el objetivo de pérdida: no se inventa un ángulo alternativo de 15 grados. No se acepta una velocidad de referencia cero. La validación no exige $V_{s0}\le V_{s1}$: si un avión publica referencias invertidas pero dentro de rango, la interpolación utiliza esos datos tal cual.

\subsubsection{Estados y variables que no participan}
Se requiere vuelo, actividad sin pausa/slew y muestras recientes. No se utiliza directamente \code{DYNAMIC PRESSURE}, viento, posición del yugo, trim, aceleración G ni bocina de pérdida. La validación compartida de la trama puede retirar la salida si otros campos esenciales son inválidos. Los parámetros de presión de referencia o de ráfagas no ajustan este buffet.

\subsection{Paso 1: promedio de flaps}
\begin{equation}
f=\frac{f_L+f_R}{2}.
\end{equation}
$f=0$ representa ambos retraídos; $f=1$, ambos completamente desplegados. El promedio se utiliza en dos lugares distintos: ángulo de inicio y velocidad de referencia.

\subsubsection{Qué se pierde al promediar}
Flaps $(1,0)$ y $(0.5,0.5)$ producen $f=0.5$, aunque las configuraciones no sean físicamente equivalentes. El modelo no crea un buffet lateral por asimetría de flaps ni calcula cada ala por separado. Esa simplificación debe tenerse presente al interpretar el reparto fijo de roll y pitch.

\subsection{Paso 2: determinar dónde empieza el buffet}
\subsubsection{Ángulo de comienzo}
\begin{equation}
\alpha_0=\alpha_s(0.5+0.2f).
\label{eq:stall-onset}
\end{equation}
$\alpha_s$ es el ángulo de pérdida recibido. El factor 0.5 fija el comienzo limpio al 50\% de ese ángulo; el término $0.2f$ lo desplaza hasta el 70\% con flap completo. Los números 0.5/0.2 se encuentran en la referencia local de TelemFFB; son decisiones de su estimador, no una ley universal de separación de flujo.

Para $\alpha_s=15^\circ$: con $f=0$, el inicio es $7.5^\circ$; con $f=0.5$, $9^\circ$; con $f=1$, $10.5^\circ$. Los 15 grados son un dato supuesto para los ejemplos, no una constante que MiYoke imponga al avión.

\subsubsection{Rampa normalizada de ángulo}
\begin{equation}
R_\alpha=\sat_{[0,1]}\left(\frac{\alpha-\alpha_0}{\alpha_s-\alpha_0}\right).
\label{eq:stall-ramp}
\end{equation}
El numerador indica cuánto se avanzó desde el comienzo; el denominador es el ancho de la banda. Debajo o justo en $\alpha_0$, resulta cero. A mitad de banda vale 0.5 y desde $\alpha_s$ vale 1. No sigue creciendo sin límite después de ese ángulo.

Los rangos de validación garantizan $\alpha_s>0$ y $0.5\le0.5+0.2f\le0.7$, de modo que el denominador es positivo. No hay división por cero con entradas aceptadas.

\subsection{Paso 3: referencia de velocidad y factor de flujo}
\subsubsection{Interpolación de VS0 y VS1}
\begin{equation}
V_{\rm ref}=V_{s1}+(V_{s0}-V_{s1})f
=(1-f)V_{s1}+fV_{s0}.
\label{eq:stall-vref}
\end{equation}
Con flaps retraídos toma VS1; con flaps completos toma VS0. Con media extensión toma su promedio. Para VS1=60 kt y VS0=45 kt, el resultado a $f=0.5$ es 52.5 kt. Es una interpolación lineal de referencias, no una solución del modelo aerodinámico del avión.

\subsubsection{Cociente al cuadrado y límite 1.5}
\begin{equation}
B=\min\left[\left(\frac{V}{V_{\rm ref}}\right)^2,1.5\right].
\label{eq:stall-flow}
\end{equation}
$V/V_{\rm ref}$ compara IAS con la referencia de diseño; elevar al cuadrado refuerza su dependencia con velocidad. Si IAS coincide con la referencia, $B=1$; si es la mitad, $B=0.25$. El máximo 1.5 se alcanza desde $V/V_{\rm ref}=\sqrt{1.5}\simeq1.224745$.

La relación cuadrática está motivada por $q=\rho V^2/2$, cuyo fundamento general explica \href{https://www1.grc.nasa.gov/beginners-guide-to-aeronautics/dynamic-pressure-2/}{NASA Glenn}. Pero esta fórmula usa IAS y velocidades de diseño sin medir densidad: no debe presentarse como el cálculo exacto de $q/q_{\rm stall}$ ni como una medición del esfuerzo de mando. El límite 1.5 proviene del estimador de referencia TelemFFB.

\subsubsection{Flaps no significa simplemente más o menos vibración}
Al aumentar $f$, el inicio $\alpha_0$ sube y puede reducir $R_\alpha$ a AoA fijo. Si VS0 es menor que VS1, al mismo tiempo baja $V_{\rm ref}$ y aumenta $B$. Son efectos contrapuestos. Por ello no existe una regla general de que más flaps siempre aumenten o siempre reduzcan este buffet.

\subsection{Paso 4: envolvente final de la Pico}
\begin{equation}
E_s=\min(1,R_\alpha B),\qquad 0\le E_s\le1.
\label{eq:stall-envelope}
\end{equation}
La rampa y el factor de flujo se multiplican; después se limita el producto a uno. Si $R_\alpha=0$, no hay objetivo aunque IAS sea alta. Si IAS es cero, tampoco hay objetivo aunque AoA supere el ángulo de pérdida.

\subsubsection{Saturación antes de alcanzar el ángulo de pérdida}
Con $B=1.5$, la amplitud llega a uno cuando $R_\alpha\ge2/3$. Por tanto, puede saturar antes de $\alpha_s$. El ángulo de saturación, si $B\ge1$, es
\begin{equation}
\alpha_{\rm sat}=\alpha_0+\frac{\alpha_s-\alpha_0}{B}.
\end{equation}
Con $\alpha_s=15^\circ$, flaps cero y flujo limitado a 1.5, da $\alpha_{\rm sat}=12.5^\circ$. Si $B<1$, ni siquiera con $R_\alpha=1$ se llega a envolvente uno: el máximo es $B$. El modelo no afirma que a $12.5^\circ$ el avión ya esté en pérdida; únicamente saturó su señal háptica.

\subsubsection{Qué ocurre después de la banda}
Con $\alpha\ge\alpha_s$, $R_\alpha$ permanece en uno. El buffet no desaparece ni se debilita por entrar más profundamente en pérdida, salvo cambios de velocidad o datos. Su frecuencia sigue fija. No representa por sí solo la evolución completa de una pérdida profunda.

\subsection{Paso 5: USB, filtro y señal firmada}
\subsubsection{Transporte}
SC8 transmite las entradas crudas de pérdida. \code{PicoForceModels.h} calcula $E_s$ mediante \code{BuffetModel.h} y guarda el resultado en \code{buffet[0]}. $E_{s,\rm rx}=E_s$ en coma flotante, sin truncado a diezmilésimas. CRC, secuencia, armado y frescura condicionan la entrada; la síntesis local continúa nominalmente cada 2 ms.

\subsubsection{Suavizado de 100 ms}
La amplitud interna cumple
\begin{equation}
A_{s,n}=A_{s,n-1}+\frac{\Delta t_{\rm ms}}{100+\Delta t_{\rm ms}}
(E_{s,\rm rx,n}-A_{s,n-1}).
\end{equation}
Es un filtro de primer orden que suaviza entrada y salida de la condición. Con 2 ms, se corrige 1.96\% de la diferencia por ciclo. Un escalón desde cero llega al 62.85\% a 100 ms y 94.87\% a 300 ms; no hay una espera fija de 100 ms antes de vibrar.

Si el objetivo pasa a cero pero la comunicación sigue armada, queda cola: $A_{s,n}=A_{s,0}(100/102)^n$ con esos intervalos. Si el avión vuelve al suelo, el objetivo es cero, pero esa cola puede coexistir brevemente con rodaje. El desarme general sí anula el efecto y reinicia el estado.

\subsubsection{Triangular de 13 Hz}
\begin{align}
\phi_{s,n}&=(\phi_{s,n-1}+13\Delta t_{\rm ms})\bmod1000,\\
T(\phi)&=\begin{cases}0.004\phi-1,&\phi<500,\\3-0.004\phi,&\phi\ge500,\end{cases}\\
w_s&=A_sT(\phi_s).
\end{align}
La fase arranca en 250 al reiniciar, correspondiente a cero. Con ciclos de 2 ms avanza 26 unidades por ciclo; el período nominal es $1/13\simeq76.923$ ms. La triangular tiene pico unitario y media ideal cero a amplitud fija; al multiplicarla por $A_s$ aparecen las oscilaciones firmadas entre $-A_s$ y $+A_s$.

No se utiliza ruido, dirección aleatoria, variación de frecuencia con profundidad ni compensación RMS de la variante Dynamic de TelemFFB. Una frecuencia fija a 13 Hz es una simplificación local del renderizado del buffet.

\subsection{Paso 6: ganancias, roll/pitch y PWM}
\subsubsection{Comando por eje}
Con ganancia ``Buffet perdida'' $c_s=0.24$ y general $G=0.90$:
\begin{align}
u_{P,s}&=Gc_sw_s=0.216 A_sT,\\
u_{R,s}&=0.5Gc_sw_s=0.108 A_sT.
\end{align}
El signo depende de la fase, no del signo de posición del yugo. Ambos ejes comparten fase; roll tiene la mitad del comando. El efecto puede vibrar con el mando centrado. No aplica la zona muerta de turbulencias y no usa el factor 10 de suelo.

\subsubsection{Conversión eléctrica}
Se suma con los demás efectos, se limita a $L=220/255$, se aplica polaridad de motor, se limita a $[-1,1]$, se multiplica por \code{PWM_MAX}=255, se redondea simétricamente y se impone el techo duro de 220:
\begin{equation}
p_a=\sat_{[-220,220]}\left\{\operatorname{round}_{\rm sim}
\left[255\,\sat_{[-1,1]}\left(s_a\sat_{[-L,L]}(u_{a,s})\right)\right]\right\}
\end{equation}
para el efecto aislado. El pico ideal inicial es 55.08 cuentas pitch y 27.54 roll: redondea a \textbf{55 y 28}. No existe un techo propio de buffet inferior al general; estas magnitudes provienen de sus ganancias iniciales.

PWM positivo activa RIGHT y deja LEFT en cero; negativo hace lo contrario. En Pico, roll GP6/GP7 y pitch GP8/GP9, con PWM configurado a 490 Hz. Los 13 Hz describen la oscilación del comando, no esa frecuencia eléctrica. Ni comando normalizado ni PWM son fuerza física calibrada.

\subsection{Ejemplos completos}
\subsubsection{Flaps a mitad, AoA de 12 grados e IAS de 60 kt}
Se supone vuelo válido, $\alpha_s=15^\circ$, VS0=45 kt, VS1=60 kt, $f_L=f_R=0.5$, $\alpha=12^\circ$ e IAS=60 kt:
\begin{enumerate}
\item $f=0.5$; $\alpha_0=15(0.5+0.2\times0.5)=9^\circ$.
\item $V_{\rm ref}=60+(45-60)0.5=52.5$ kt.
\item $R_\alpha=(12-9)/(15-9)=0.5$.
\item $B=(60/52.5)^2=1.306122$, menor que 1.5.
\item $E_s=0.5\times1.306122=0.653061$, calculado localmente sin truncado.
\item Tras estabilizarse el filtro, $A_s\simeq0.6530$.
\item Para una fase de ejemplo con $T=0.6$, $w_s=0.3918$.
\item $u_P=0.3918\times0.24\times0.90=0.0846288$ y $u_R=0.0423144$.
\item PWM: $p_P=22$ y $p_R=11$, con polaridades positivas. En fase opuesta salen $-22$ y $-11$.
\end{enumerate}
Con portadora en pico ideal, para la misma envolvente los picos son 36 PWM pitch y 18 roll. Son ejemplos matemáticos; el avión debe publicar los datos reales de su modelo y no se supone que todo avión tenga esas referencias.

\subsubsection{Comparación de ángulo, flaps y velocidad}
Se conserva $\alpha_s=15^\circ$, VS0=45 kt y VS1=60 kt. Tabla de régimen y pico ideal $|T|=1$, con envolventes redondeadas para presentación:
\begin{center}\begin{tabular}{rrrrrr}
\toprule AoA (grados)&$f$&IAS (kt)&$E_{s,\rm rx}$&PWM roll&PWM pitch\\\midrule
7.5&0&60&0&0&0\\
10&0&60&0.3333&9&18\\
12&0&60&0.6000&17&33\\
15&0&60&1.0000&28&55\\
20&0&60&1.0000&28&55\\
12&0.5&60&0.6530&18&36\\
12&1&60&0.5000&14&28\\
12&0&90&0.9000&25&50\\
12.5&0&90&1.0000&28&55\\
15&0&30&0.2500&7&14\\\bottomrule
\end{tabular}\end{center}
La tabla presenta envolventes redondeadas; no representa enteros transmitidos. SC8 conserva precisión flotante hasta el comando final. Los valores no predicen una trayectoria aerodinámica real: se evalúan combinaciones de entradas para entender la función.

\subsubsection{Arranque con objetivo recibido 0.6530}
Desde amplitud cero y fase 250, con pasos de 2 ms:
\begin{center}\begin{tabular}{rrrrrr}
\toprule ms&Fase&$A_s$&$T$&PWM roll&PWM pitch\\\midrule
2&276&0.012804&0.104&0&0\\
10&380&0.061558&0.520&1&2\\
20&510&0.117313&0.960&3&6\\
50&900&0.254976&$-0.600$&$-4$&$-8$\\
100&550&0.410392&0.800&9&18\\
200&850&0.562865&$-0.400$&$-6$&$-12$\\
300&150&0.619512&$-0.400$&$-7$&$-14$\\\bottomrule
\end{tabular}\end{center}
La intensidad de una muestra depende de amplitud y fase a la vez. Una envolvente creciente no implica que cada muestra PWM sea mayor que la anterior.

\subsection{Procedencia: TelemFFB, física y adaptación propia}
\subsubsection{Fórmulas reconocibles de TelemFFB}
La copia local de \code{BuffetingEffectMixIn.py::ac_update_buffeting} contiene el promedio de flaps, la banda de AoA que empieza en $\alpha_s(0.5+0.2f)$, la interpolación VS0/VS1 y el factor cuadrático limitado por \code{ACCEL_STALL_Q_CAP=1.5}. La frecuencia base configurada allí es 13 Hz. Son antecedentes concretos, no sólo una semejanza de nombres.

Referencia pública de localización: \href{https://github.com/walmis/VPforce-TelemFFB/blob/refactor/telemffb/sim/base/BuffetingEffectMixIn.py}{TelemFFB, BuffetingEffectMixIn.py}. Se compara la copia local; la rama pública puede evolucionar.

\subsubsection{Diferencias que afectan los números}
\begin{itemize}
\item TelemFFB utiliza TAS en el factor de velocidad de esa ruta; MiYoke utiliza IAS. No se consideran numéricamente equivalentes en todas las condiciones.
\item La referencia admite caminos alternativos si faltan umbrales o velocidades. MiYoke exige los datos y rangos descritos; no inventa reemplazos.
\item TelemFFB incluye su intensidad en el producto \emph{antes} del límite final de magnitud. MiYoke limita primero $R_\alpha B$ a uno y después aplica ganancia de placa. No es el mismo orden de saturación.
\item Su renderizado Classic y Dynamic contempla diferencias de onda/dirección; Dynamic agrega envolvente irregular, frecuencia dependiente de profundidad y compensación de amplitud. MiYoke mantiene una triangular fija, misma fase por eje y relación 1:2.
\end{itemize}
Para ver el tercer punto, si $R_\alpha B=1.5$ y una intensidad de referencia fuese 0.24, multiplicar antes de saturar produciría 0.36; MiYoke limita primero a 1 en la Pico y aplica después 0.24, antes de la ganancia general. Por eso no debe describirse como copia numérica completa de TelemFFB.

\subsubsection{Original alemán y constantes locales}
No se identificó en las fuentes originales locales la cadena con AoA, flaps, VS0/VS1 y SimConnect de esta sección. El esquema genérico de motor y efectos FFB es un antecedente de arquitectura, no el origen de este estimador. El filtro de 100 ms, protocolo, reparto por eje, ganancia local 0.24 y conversión PWM pertenecen a la implementación modificada.

La presión dinámica aporta una motivación general para dependencia cuadrática; no justifica por sí sola las fracciones de inicio o la fuerza máxima del yugo. Las fuentes no convierten esos parámetros en datos experimentales de un avión particular.

\subsection{Habilitación y diagnóstico}
\subsubsection{Qué lo habilita y qué lo anula}
La Pico calcula buffets con la opción activada por defecto; el puente la desactiva con \code{--no-buffets}. Para este objetivo debe haber vuelo, datos recientes, simulación en marcha y ausencia de pausa/slew. En firmware se requieren habilitación general, \code{SIMCONNECT_STALL_BUFFET_ENABLED}, bit \code{Stall}, ganancia efectiva, DTR y trama armada vigente.

Si una entrada específica de pérdida es inválida, su objetivo se hace cero sin necesidad de anular tren o spoiler válidos. Mientras el conjunto siga armado, el filtro puede dejar una cola. Si la validación global provoca desarme o vence el watchdog de 300 ms, la salida se retira. El menú activo inhibe motores y un update local mayor que 50 ms reinicia amplitudes y fases.

\subsubsection{Cómo comprobar las entradas en CSV}
Registrar \code{aoa_deg}, \code{stall_alpha_deg}, flaps, VS0/VS1, IAS, estados y marcas del CSV. La envolvente se calcula en la Pico; \code{stall_amp} y \code{Buffet_Perdida} son columnas históricas SC7 que ya no existen en el registro actual.

La envolvente local no confirma PWM ni fuerza física. La cola del filtro y el muestreo del LCD pueden ocultar picos o producir aliasing.

\subsection{Qué se puede ajustar y cómo hacerlo}
\subsubsection{Ganancias disponibles}
``Buffet perdida'' cambia $c_s$ de 0 a 1, con paso 0.01; inicio 0.24. ``Ganancia total'' cambia $G$ de 0 a 1, con paso 0.05; inicio 0.90. El interruptor de pérdida apaga ambos ejes. No hay ganancias independientes de este buffet para roll y pitch en el menú.

Antes de saturar, para un pico pitch deseado $P$ a una envolvente conocida, $c_s\simeq P/(255GE_{s,\rm rx})$. En el ejemplo $E=0.6530$, buscar 20 PWM pitch da $c_s\simeq0.13345$, aproximable a 0.13. La ganancia cambia la intensidad de mando, no el ángulo matemático de comienzo $\alpha_0$.

La percepción de comienzo sí puede retrasarse si la ganancia es baja: el primer PWM ideal requiere $E\gtrsim0.5/(255Gc_s)$ en pitch y el doble en roll. Con valores iniciales, aproximadamente 0.00908 y 0.01816. Filtro, fase y fricción agregan otras limitaciones.

\subsubsection{Ajustes del estimador y del renderizado}
Cambiar 0.5/0.2 modifica el inicio con flaps; cambiar el cap 1.5 altera refuerzo por velocidad; cambiar el uso de IAS, VS0/VS1 o sus validaciones modifica el estimador. Esas decisiones están en \code{BuffetModel.h} y requieren recompilar el firmware.

La frecuencia 13 está en el primer elemento del arreglo de \code{SimEffects.h}; cambiarla requiere firmware nuevo. Los 100 ms son compartidos: conviene separarlos si se pretende modificar sólo buffet de pérdida. El factor roll 0.5 también se aplica a los otros dos buffets; no es una calibración exclusiva de pérdida.

\subsubsection{Ensayo propuesto y mejoras pendientes}
En el simulador, aislar este buffet y registrar una secuencia de aproximación y salida de la banda con configuración conocida. Comparar flaps retraídos, intermedios y completos, sin asumir que la misma IAS y AoA representan el mismo estado físico. Confirmar la coherencia de unidades y la validez de los datos del add-on. Los ejemplos de este documento no son un procedimiento de operación de una aeronave real.

Para pulir el modelo pueden incorporarse perfiles por avión, umbrales perceptuales medidos, estados adicionales sólo si se justifican y un renderizado menos periódico. También podría separarse intensidad de referencia y profundidad para evitar que saturar demasiado pronto oculte diferencias dentro de la banda. Cualquier cambio debe distinguir el estimador del efecto visual o mecánico que se quiere transmitir.

No se debe calibrar pérdida cambiando el cero del potenciómetro: no entra directamente en su fórmula. La fuerza real requiere medición del motor/transmisión; subir PWM no certifica realismo aerodinámico.

\subsection{Fuentes y alcance}
El estimador completo está en \code{BuffetModel.h::calculate}. \code{main.cpp} define unidades, estados y transporte; \code{PicoForceModels.h} calcula \code{buffet[0]}; \code{SimEffects.h} filtra y sintetiza; \code{GainSettings.h}/\code{mis_defines.h} aportan ajustes; el sketch y \code{MotorDriver.h} convierten a PWM.

\code{tests/simconnect_checks.cpp} contiene casos de comienzo, flaps, ángulo inválido, velocidad cero, independencia de efectos y comunicación. \code{LEEME md/EFECTOS_SIMCONNECT.md} documenta el antecedente pero conserva ganancias antiguas; el valor vigente usado aquí es 0.24. Las fuentes externas son el SDK, NASA para presión dinámica y TelemFFB enlazados. Las tablas verifican el cálculo de software; no se midieron vibraciones reales ni se verificó la placa cargada.
\newpage
\section{BUFFET POR SPOILER}
\label{sec:buffet-spoiler}

\subsection{Qué hace el efecto}
Genera una vibración en ambos ejes cuando el avión está en vuelo, los spoilers están desplegados y la velocidad indicada supera 80 kt. La amplitud aumenta con el promedio de despliegue y una rampa de IAS que llega al máximo en 140 kt. La placa produce una triangular de 15 Hz; pitch recibe el doble de comando que roll.

La idea sencilla es: \emph{``cuanto más spoiler expuesto y mayor velocidad dentro de la banda, más vibración''}. Se trata de una señal háptica aproximada, no del cálculo del arrastre, la sustentación perdida o la fuerza real transmitida desde las superficies al yugo.

\subsubsection{Qué se entiende por spoiler en esta implementación}
Se leen las posiciones izquierda y derecha publicadas por el simulador. No se lee solamente la palanca de aerofrenos ni la bandera de armado. Si el avión utiliza esas superficies para ayudar al alabeo y lo publica en las variables estándar, ese despliegue puede alimentar el efecto.

No distingue despliegue simétrico de aerofrenado y despliegue asimétrico de asistencia lateral. No simula golpes de movimiento del actuador ni vibraciones de extensión/retracción como efectos separados. Tampoco requiere que la superficie esté cambiando: puede continuar vibrando mientras permanece desplegada a IAS constante.

\subsection{Recorrido y variables utilizadas}
\subsubsection{Cadena completa}
\begin{center}
\fbox{\begin{minipage}{0.91\linewidth}
Spoiler izquierdo/derecho e IAS $\rightarrow$ comprobar vuelo y datos $\rightarrow$ promedio de despliegue $\times$ rampa 80--140 kt $\rightarrow$ envolvente $0\ldots1$ $\rightarrow$ filtro local 100 ms $\rightarrow$ triangular 15 Hz $\rightarrow$ ganancias $\rightarrow$ roll mitad de pitch $\rightarrow$ suma y límite general $\rightarrow$ polaridad y PWM.
\end{minipage}}
\end{center}

\subsubsection{Variables de SimConnect}
\begin{longtable}{p{0.37\linewidth}p{0.18\linewidth}p{0.37\linewidth}}
\toprule Variable & Unidad solicitada & Uso\\\midrule\endhead
\code{SPOILERS LEFT POSITION} & percent over 100 & $s_L$, despliegue izquierdo 0--1.\\
\code{SPOILERS RIGHT POSITION} & percent over 100 & $s_R$, despliegue derecho 0--1.\\
\code{AIRSPEED INDICATED} & knots & $V$, rampa de velocidad.\\
\code{SIM ON GROUND} & Bool & Objetivo sólo si es 0.\\
\code{IS SLEW ACTIVE} & Bool & Debe ser 0 para habilitar.\\
\bottomrule
\end{longtable}
Los nombres se verifican en el \href{https://www.prepar3d.com/SDKv3/LearningCenter/utilities/variables/simulation_variables.html}{SDK de Prepar3D}; las unidades solicitadas se ven en \code{main.cpp}. El puente usa fracciones 0--1, no porcentajes 0--100. Un valor 0.4 representa 40\% de despliegue.

No se solicitan \code{SPOILERS HANDLE POSITION} o \code{SPOILERS ARMED} para esta fórmula. Presión dinámica, viento, posición del yugo, velocidad del mando, flaps y estado del tren no multiplican su amplitud. La integridad del paquete compartido y los eventos de actividad/pausa sí condicionan la salida.

\subsubsection{Rangos admitidos}
IAS debe ser finita entre 0 y 1000 kt; ambos spoilers, finitos entre 0 y 1. Si un spoiler es inválido, se anula este objetivo completo: no se reemplaza por el otro ni se hace promedio de un único dato. Una IAS fuera de rango anula los tres buffets. Otros campos opcionales inválidos pueden afectar sólo a sus respectivos efectos.

\subsection{Paso 1: promedio y rampa en la Pico}
\subsubsection{Promedio de despliegue}
\begin{equation}
\bar s=\frac{s_L+s_R}{2}.
\end{equation}
Se obtiene 0 con ambos retraídos y 1 con ambos completos. Un spoiler al máximo y otro retraído producen 0.5, igual que ambos a mitad. El modelo no conserva información de qué lado está desplegado después de este promedio.

Por tanto, el roll resultante no es una fuerza dirigida al lado del spoiler. Es la misma vibración periódica del otro eje con menor ganancia. El signo alternante se crea por fase, no por $s_L-s_R$.

\subsubsection{Rampa con velocidad indicada}
\begin{align}
R_{sp}(V)&=\sat_{[0,1]}\left(\frac{V-80}{140-80}\right),\\
E_{sp}&=\bar s R_{sp}(V).
\label{eq:spoiler-env}
\end{align}
El 80 es el inicio de la rampa en kt y $140-80=60$ kt es su ancho. Hasta 80 kt, la amplitud objetivo es cero; a 110 kt se aplica la mitad del promedio de despliegue; desde 140 kt se aplica todo. El producto no puede superar uno con entradas válidas.

No hay rampa cuadrática de presión ni referencia de 120 kt en este efecto. \code{--reference-knots} no lo ajusta. La frecuencia tampoco depende de IAS: siempre es 15 Hz en esta versión.

\subsubsection{Umbral de despliegue y continuidad}
MiYoke no exige $\bar s>0.1$: permite empezar continuamente desde cero. Con $V=140$ kt y ambos spoilers a 0.05, $E_{sp}=0.05$. La referencia local de TelemFFB sí utiliza una condición de despliegue mayor que 0.1 en su ruta de buffet; esa condición no fue copiada aquí.

Los umbrales 80/140 regulan sensación; no certifican los límites operativos o estructurales de spoilers del avión. Un avión que normalmente use spoilers por debajo de 80 kt no recibirá objetivo de este modelo en esa región, independientemente de que exista alguna vibración real.

\subsection{Paso 2: cálculo en el microcontrolador}
SC8 transmite posiciones de spoilers, IAS y estados, como flotantes. \code{PicoForceModels.h} calcula $E_{sp}$ y lo guarda en \code{buffet[2]}. $E_{sp,\rm rx}=E_{sp}$ salvo precisión de flotantes. La entrada se renueva por muestra aceptada; la portadora triangular de 15 Hz y el filtro se calculan localmente en ciclos nominales de 2 ms. No hay entero de envolvente ni resolución USB de 0.0001.

\subsection{Paso 3: filtro, frecuencia y cola}
\subsubsection{Amplitud suavizada}
\begin{equation}
A_{sp,n}=A_{sp,n-1}+\frac{\Delta t_{\rm ms}}{100+\Delta t_{\rm ms}}
(E_{sp,\rm rx,n}-A_{sp,n-1}).
\end{equation}
La variable $A_{sp}$ sigue al objetivo con un filtro de primer orden de constante 100 ms. Con un ciclo de 2 ms se corrige aproximadamente el 1.96\% de la diferencia pendiente; partiendo de cero, llega al 62.85\% a 100 ms y al 94.87\% a 300 ms. El origen de la ecuación es $\tau\dot A+A=E$, discretizada, no un modelo medido del actuador.

\subsubsection{Portadora triangular}
\begin{align}
\phi_{sp,n}&=(\phi_{sp,n-1}+15\Delta t_{\rm ms})\bmod1000,\\
T(\phi)&=\begin{cases}0.004\phi-1,&\phi<500,\\3-0.004\phi,&\phi\ge500,\end{cases}\\
w_{sp}&=A_{sp}T(\phi_{sp}).
\end{align}
La fase recorre 1000 unidades por ciclo; a 15 Hz el período es aproximadamente 66.667 ms. En 2 ms avanza 30 unidades. Al reiniciar se fija en 250, cruce por cero. La triangular cambia entre valores negativos y positivos y, a amplitud estable, tiene media ideal cero.

La onda firmada $w_{sp}$ ya queda entre $-1$ y $+1$, pero después se atenúa por ganancias. No hace falta alcanzar los extremos de esa escala para sentir buffet.

\subsubsection{Retirada de spoilers y contacto con suelo}
Al retraer, bajar a 80 kt o menos, o quedar en suelo, la Pico lleva $E_{sp}$ a cero. Mientras el conjunto siga habilitado, $A_{sp}$ decae con el filtro y deja una cola. No se reinicia la fase sólo por pasar el objetivo a cero. El suelo SC8, en cambio, se anula al parar o quedar en vuelo sin conservar una cola de su filtro en la Pico; no comparte esta envolvente de spoilers.

Pausa, desarme, DTR perdido, timeout de telemetría o reinicio por intervalo local mayor que 50 ms sí retiran o reinician la salida general. Por ello un touchdown puede dejar una breve cola de este buffet mientras comienza el efecto de suelo; no significa que la Pico siga calculando un objetivo de spoiler en tierra.

\subsection{Paso 4: ganancias y conversión a PWM}
\subsubsection{Roll y pitch}
Con ``Buffet spoiler'' $c_{sp}=0.16$ y ganancia total $G=0.90$:
\begin{align}
u_{P,sp}&=Gc_{sp}A_{sp}T=0.144A_{sp}T,\\
u_{R,sp}&=0.5Gc_{sp}A_{sp}T=0.072A_{sp}T.
\end{align}
Ambos ejes comparten fase. La relación de comando roll/pitch es 1:2 antes de redondear y de aplicar polaridades. No depende del desplazamiento del yugo ni usa la zona muerta 0.02 de turbulencias. La ganancia de centrado o la de suelo tampoco se aplica a este aporte.

\subsubsection{Saturación y driver}
En modo aislado se limita el comando a $\pm220/255$, se aplica signo de motor $s_a$, se limita a $[-1,1]$ y se convierte mediante
\begin{equation}
p_a=\sat_{[-220,220]}\left[\operatorname{round}_{\rm sim}
\left(255\,\sat_{[-1,1]}\left(s_a\sat_{[-220/255,220/255]}(u_{a,sp})\right)\right)\right].
\end{equation}
El 255 es \code{PWM_MAX} predeterminado, no la ganancia del efecto. El redondeo simétrico produce enteros con igual tratamiento de signos. Los máximos ideales aislados, con $A=|T|=1$, son 36.72 cuentas pitch y 18.36 roll: \textbf{37 y 18 PWM}. No hay un tope especial de 40 como en suelo.

Con PWM positivo se activa RIGHT y LEFT queda en cero; negativo activa LEFT y anula RIGHT. Roll usa GP6/GP7 y pitch GP8/GP9 en Pico. La portadora de 15 Hz modula salidas PWM configuradas a 490 Hz. Duty cycle es $|p|/255$; no es una medición de torque. Si otros efectos están activos, el límite general corresponde a la suma.

\subsection{Ejemplos numéricos completos}
\subsubsection{IAS 110 kt, spoilers a 80\% y 40\%}
Supuestos: vuelo y datos válidos, ganancias iniciales, filtro estabilizado y demás aportes nulos.
\begin{enumerate}
\item $s_L=0.8$, $s_R=0.4$, por lo que $\bar s=0.6$.
\item $R_{sp}=(110-80)/60=0.5$.
\item $E_{sp}=0.6\times0.5=0.3$; la Pico usa directamente esa envolvente en coma flotante.
\item En régimen, $A_{sp}\simeq0.3$. Para una fase de ejemplo con $T=0.8$, $w_{sp}=0.24$.
\item $u_P=0.24\times0.16\times0.90=0.03456$; $u_R=0.01728$.
\item PWM pitch: $\operatorname{round}(8.8128)=9$; roll: $\operatorname{round}(4.4064)=4$.
\end{enumerate}
En la fase opuesta salen $-9$ y $-4$ con polaridades positivas. En pico ideal, esa envolvente da 11 cuentas pitch y 6 roll. El hecho de que el spoiler izquierdo esté más desplegado no impone signo fijo al roll.

\subsubsection{Tabla de picos ideales}
Se supone $|T|=1$, régimen estabilizado y ganancias iniciales:
\begin{center}\begin{tabular}{rrrrrr}
\toprule IAS (kt)&$s_L$&$s_R$&$E_{sp}$&PWM roll&PWM pitch\\\midrule
80&1&1&0&0&0\\95&1&1&0.25&5&9\\110&1&1&0.50&9&18\\125&1&1&0.75&14&28\\140&1&1&1&18&37\\180&1&1&1&18&37\\110&0.8&0.4&0.30&6&11\\140&1&0&0.50&9&18\\140&0.05&0.05&0.05&1&2\\\bottomrule
\end{tabular}\end{center}
Las filas de velocidad mayor que 140 muestran saturación de la \emph{rampa}, no del driver: el objetivo ya no crece. Es distinto de llegar al techo general de 220 PWM, que aquí no se alcanza con el efecto aislado y sus ganancias predeterminadas.

\subsubsection{Arranque con objetivo 0.3}
Desde amplitud cero, fase 250 y ciclos exactos de 2 ms:
\begin{center}\begin{tabular}{rrrrrr}
\toprule ms&Fase&$A_{sp}$&$T$&PWM roll&PWM pitch\\\midrule
2&280&0.005882&0.12&0&0\\10&400&0.028281&0.60&0&1\\20&550&0.053896&0.80&1&2\\50&0&0.117141&$-1.00$&$-2$&$-4$\\74&360&0.155817&0.44&1&3\\100&750&0.188542&0&0&0\\300&750&0.284615&0&0&0\\\bottomrule
\end{tabular}\end{center}
La amplitud crece aunque las muestras elegidas a 100 y 300 ms den cero. La fase determina cuándo se ve el pico, no sólo la extensión del spoiler.

\subsection{Procedencia y diferencias con TelemFFB}
\subsubsection{Promedio y umbrales}
La copia local de \code{MotionEffectsMixIn.py::ac_update_spoilers} promedia una lista de posiciones, pondera por velocidad y utiliza umbrales definidos como $80\,\code{kt2ms}$ y $140\,\code{kt2ms}$. Aquí sí hay una correspondencia directa de las referencias de 80/140 kt: TelemFFB las convierte a m/s y MiYoke opera directamente en knots.

La velocidad de esa ruta de TelemFFB es IAS si está disponible y TAS como alternativa. MiYoke exige IAS válida. También se omite la condición de despliegue $>0.1$ de la referencia, por decisión local de permitir una entrada continua desde cero.

\subsubsection{Generación de la vibración}
TelemFFB combina senoides de 15 y 16 Hz en una dirección, y de 14 y 18 Hz en otra, además de una ruta separada de movimiento de superficies. MiYoke usa una única triangular de 15 Hz para los dos ejes, con factor roll 0.5 y filtro local de 100 ms. No copia las cuatro ondas ni el efecto del actuador.

El uso de una triangular a 15 Hz puede dar una sensación más regular y con contenido armónico distinto. Que una frecuencia coincida no significa que las señales ni el RMS sean iguales. No se ha calibrado una equivalencia perceptual automática entre ambos equipos.

\subsubsection{Original alemán y decisiones propias}
En las fuentes originales locales no se identificó esta cadena SimConnect de despliegue e IAS. El accionamiento por dos PWM conserva un antecedente del original; el estimador simplificado y su integración corresponden al proyecto modificado. Las ganancias 0.16, factor roll 0.5, filtro, protocolo y límites del driver son decisiones locales, no coeficientes aerodinámicos medidos.

Fuente pública para localizar el antecedente: \href{https://github.com/walmis/VPforce-TelemFFB/blob/refactor/telemffb/sim/base/MotionEffectsMixIn.py}{TelemFFB, MotionEffectsMixIn.py}. Las afirmaciones se contrastaron con la copia local inspeccionada.

\subsection{Habilitación y diagnóstico}
\subsubsection{Banderas, datos y apagado}
La Pico calcula este buffet con la opción activada por defecto; \code{--no-buffets} pone en cero los tres objetivos. No se necesita \code{--gusts}. Se requiere vuelo, simulación en marcha sin pausa/slew y datos recientes. Para actuar, además puerto reconocido y \code{--enable}.

En la placa intervienen habilitación general, \code{SIMCONNECT_SPOILER_BUFFET_ENABLED}, máscara \code{Spoiler}, ganancia efectiva, DTR y trama armada de menos de 300 ms. El menú inhibe motores. Una entrada de spoiler inválida pone su objetivo a cero, con posible cola filtrada si sigue armado; un desarme global retira la salida. No se debe confundir ese comportamiento con una vibración persistente de actuador, que no está implementada.

\subsubsection{CSV y pantalla}
Registrar IAS, spoilers izquierdo/derecho, estado de suelo y armado. Las columnas históricas \code{spoiler_amp} y \code{Buffet_Spoilers} no existen en SC8; la Pico calcula el objetivo y su cola filtrada.

La pantalla de 10 Hz no reconstruye una señal de 15 Hz. Con fase inicial 250 y muestreo exacto cada 100 ms, la fase observada alterna entre 250 y 750, ambos cruces por cero: puede mostrar siempre cero en un caso ideal sincronizado. En otras sincronizaciones puede parecer una oscilación más lenta. El diagnóstico debe usar una tasa adecuada o medidas de pico/RMS; una única lectura no evalúa la vibración.

\subsection{Ajustes, calibraciones y mejoras}
\subsubsection{Controles existentes}
``Buffet spoiler'' admite 0--1 con paso 0.01, inicio 0.16. Su interruptor apaga ambos ejes. ``Ganancia total'' admite 0--1 con paso 0.05, inicio 0.90, y afecta todo el conjunto. Los ajustes guardados prevalecen si son válidos; recompilar defaults no garantiza que los sustituya.

No existe una ganancia de spoiler independiente por eje ni un balance izquierdo/derecho en el menú. El factor 0.5 de roll se comparte con los tres buffets. Cambiar la calibración ADC del mando no modifica esta amplitud.

\subsubsection{Rampa, frecuencia y filtro}
Los 80/140 kt están en \code{BuffetModel.h}: ajustar requiere recompilar el firmware. Debe mantenerse un ancho de banda positivo, $V_{\rm alto}>V_{\rm bajo}$. La frecuencia 15 es el tercer elemento del arreglo en \code{SimEffects.h}; se cambia en firmware. Los 100 ms son compartidos con otras envolventes y rigidez, por lo que modificar el literal no afecta únicamente spoiler.

Si la rampa ya saturó por velocidad, subir IAS no aumenta la señal. Si la suma está saturada por otros efectos, subir ganancia puede no aumentar el pico final. Son dos saturaciones distintas. Antes de cambiar parámetros hay que identificar cuál está actuando.

\subsubsection{Ejemplo de ajuste de ganancia}
En régimen y antes de recorte, $P_P\simeq255Gc_{sp}E_{sp}$ es el pico pitch ideal. Despejando, $c_{sp}\simeq P_P/(255GE_{sp})$. Para $E_{sp}=0.3$, $G=0.9$ y pico deseado de 8 cuentas, da aproximadamente 0.11619; el menú permite 0.12 como aproximación.

La primera cuenta de PWM ideal requiere envolvente aproximadamente mayor que 0.01362 en pitch y 0.02723 en roll con ganancias iniciales. Redondeo PWM, fase y mecánica pueden impedir percibir señales algo mayores. No es la zona muerta explícita de turbulencia, sino cuantización y respuesta del sistema.

\subsubsection{Ensayo y ampliaciones propuestas}
En el simulador, aislar este efecto con pausa reconocida al editar el menú. Comprobar primero posiciones publicadas y rampa en CSV; comparar despliegues simétricos y asimétricos para confirmar que el promedio se comporta como se espera. Registrar varios ciclos de PWM y esperar unos 0.5 s para comparar valores de régimen.

Se puede mejorar con perfiles por aeronave, separación entre aerofrenos y asistencia al alabeo, ganancias por eje y un renderizado menos periódico. También se podría conservar información de asimetría, pero asignarle una fuerza lateral requeriría una justificación específica del avión y transmisión: no basta con convertir automáticamente $s_L-s_R$ en un empuje del yugo.

Medir fuerza o aceleración real del mecanismo permite ajustar percepción. Los números de la referencia TelemFFB no calibran por sí solos un motor distinto. No se implementó ninguna de estas ampliaciones durante el análisis.

\subsection{Cómo conviven los tres buffets}
\subsubsection{Suma implementada y frecuencias}
En \code{SimEffects.h}, antes de la ganancia general, la vibración común es
\begin{equation}
v_b=c_sA_sT(\phi_s)+c_gA_gT(\phi_g)+c_{sp}A_{sp}T(\phi_{sp}).
\end{equation}
Pitch recibe $Gv_b$ y roll $0.5Gv_b$. Las frecuencias vigentes son respectivamente 13, 10 y 15 Hz. Las fases arrancan en 250 al reiniciar pero luego avanzan a ritmos diferentes. Por eso pueden reforzarse o cancelarse en algunos instantes; no se suman siempre sus valores absolutos.

No hay un limitador vectorial conjunto como en turbulencia. Después se suman centrado, turbulencia y cargas de movimiento según la ruta del código; el total por eje se recorta a $\pm220/255$. El suelo tiene su propio límite antes de agregarse. Los tres buffets no reciben el techo de 40 de suelo.

\subsubsection{Cota de los tres juntos}
Con amplitudes no mayores que uno y ganancias iniciales, una cota del comando de los tres buffets es
\begin{align}
|u_{P,b}|&\le0.90(0.24+0.16+0.16)=0.504,\\
|u_{R,b}|&\le0.252.
\end{align}
Equivale a cotas redondeadas de 129 PWM pitch y 64 roll. Son cotas por suma de magnitudes; no se afirma que las fases y condiciones reales permitan alcanzar esos picos simultáneamente. Con otros efectos o ganancias mayores sí puede intervenir el límite general 220.

Una lectura de PWM total no permite atribuir intensidad a cada buffet. Para evaluar cada fórmula hay que aislarlo o agregar un diagnóstico de aportes independientes. Retirar el argumento \code{--no-buffets} habilita el cálculo de los tres; el menú o las banderas individuales seleccionan cuáles llegan a la fuerza.

\subsubsection{Qué tendría sentido registrar}
Para cada efecto: entradas crudas, objetivo calculado en la Pico, amplitud filtrada, fase, portadora y aporte por eje antes de la suma. Agregar máximo, mínimo y RMS por ventana evitaría depender de una lectura sincronizada con un cruce por cero. El filtro de cada buffet y la tasa de pantalla deben distinguirse del ritmo de la telemetría.

\subsection{Fuentes, pruebas y alcance de las tres secciones}
\subsubsection{Mapa de implementación}
\begin{longtable}{p{0.39\linewidth}p{0.53\linewidth}}
\toprule Archivo & Papel en los tres buffets\\\midrule\endhead
\code{BuffetModel.h} & Validaciones y las tres envolventes: pérdida, tren y spoiler.\\
\code{SimConnectBridge/main.cpp} & SimVars, unidades, estados, argumento \code{--no-buffets}, codificación y columnas CSV.\\
\code{SimTelemetry.h} & Telemetría cruda SC8, rangos, CRC, secuencia y watchdog.\\
\code{SimEffects.h} & Filtros, portadoras 13/10/15 Hz, ganancias, proporción por eje y mezcla.\\
\code{mis_defines.h} & Ganancias iniciales, frecuencia configurable de tren y banderas.\\
\code{GainSettings.h} & Menú, ganancias efectivas e interruptores persistidos.\\
\code{MiYoke_PD_V2.4.ino} & Ciclo nominal 2 ms, inhibición por menú, polaridad y diagnóstico.\\
\code{MotorDriver.h} & Conversión $[-1,1]$ a PWM, redondeo y techo 220.\\
\code{tests/simconnect_checks.cpp} & Pruebas existentes de envolventes, validez, síntesis, estados y límites.\\
\bottomrule
\end{longtable}

\subsubsection{Referencias de consulta}
\begin{enumerate}
\item \href{https://www.prepar3d.com/SDKv3/LearningCenter/utilities/variables/simulation_variables.html}{SDK de Prepar3D v3, Simulation Variables}: nombres y unidades. Las peticiones concretas de conversión se verificaron en \code{main.cpp}.
\item \href{https://github.com/walmis/VPforce-TelemFFB/blob/refactor/telemffb/sim/base/MotionEffectsMixIn.py}{TelemFFB, MotionEffectsMixIn.py}: métodos de tren y spoilers; comparación contra copia local.
\item \href{https://github.com/walmis/VPforce-TelemFFB/blob/refactor/telemffb/sim/base/BuffetingEffectMixIn.py}{TelemFFB, BuffetingEffectMixIn.py}: estimador y renderizado de pérdida; comparación contra copia local.
\item \href{https://www1.grc.nasa.gov/beginners-guide-to-aeronautics/dynamic-pressure-2/}{NASA Glenn, Dynamic Pressure}: fundamento general de dependencia cuadrática usado para explicar la motivación del factor de flujo de pérdida, no para justificar las rampas lineales de tren/spoiler.
\item Fuentes locales del proyecto alemán original y \code{LEEME md/EFECTOS_SIMCONNECT.md}: antecedentes de arquitectura y adaptación. La documentación histórica contiene ganancias/protocolo/frecuencias que se contrastaron con el código vigente.
\end{enumerate}
\textbf{Alcance de la comprobación:} se reconstruyeron fórmulas y tablas numéricas y se distinguió objetivo, filtro, señal normalizada y PWM. No se modificaron los modelos de fuerza, no se cargó firmware y no se midió vibración física. Los defaults utilizados son 0.24/0.16/0.16 con ganancia general 0.90; los ajustes reales guardados pueden ser diferentes.
\newpage
\section{DAMPER}
\label{sec:damper}

\subsection{Qué representa y qué se siente}
El damper es una resistencia al movimiento: cuanto más rápido se mueve el yugo, mayor es el comando que intenta frenarlo. Se parece a mover una palanca dentro de un fluido viscoso. Puede reducir el rebote del conjunto al soltarlo, pero el resultado depende del motor, la transmisión, los sensores y los retardos reales.

Su diferencia fundamental con el resorte es la entrada utilizada. El resorte mira \textbf{dónde está} el mando; el damper mira \textbf{qué tan rápido se mueve}. Un mando desplazado y quieto genera centrado, pero, una vez detenida la estimación de velocidad, no genera damper. Un mando que atraviesa el centro puede tener centrado cero y damper distinto de cero.

La idea viscosa básica es $F_d=-b\,\dot{x}$. El signo menos significa oposición al movimiento, y $b$ determina cuánto cuesta moverlo rápidamente. En este firmware se implementa una versión normalizada, con filtrado, zona muerta y límites. \textbf{No se calcula directamente una fuerza en newtons ni un torque en N\,m.} Tampoco se usa la velocidad de rotación del avión: se deriva la posición del yugo seleccionado.

\subsubsection{Alcance y versión analizada}
Las fórmulas siguientes corresponden a \code{SimEffects.h::motionAdjustedForce}, \code{YokeMotion.h} y el sketch actual. Roll y pitch tienen estados de movimiento independientes, pero comparten los parámetros de filtrado, la ganancia de damper y la de fricción. La intensidad de roll se reduce a la mitad mediante un factor fijo en el código. Los valores indicados son defaults; el menú puede cargar otros desde memoria.

\subsection{Recorrido completo desde la entrada hasta el motor}
\begin{center}
\fbox{\begin{minipage}{0.91\linewidth}
\textbf{Fuente elegida al compilar:} posición de SimConnect, ADC de potenciómetro o cuentas de encoder.

\textbf{Movimiento:} posición normalizada $x\in[-1,1]$ $\rightarrow$ diferencia de posiciones dividida por tiempo real $\rightarrow$ límite de velocidad $\rightarrow$ filtro de 40 ms, tratamiento de inversión y reposo $\rightarrow$ velocidad estimada $\widehat v$.

\textbf{Resistencia:} quitar 0.02 a la magnitud de velocidad $\rightarrow$ multiplicar por ganancia de damper $\rightarrow$ sumar resistencia de fricción, si está habilitada $\rightarrow$ límite común 0.40 $\rightarrow$ ganancia general y factor del eje $\rightarrow$ signo contrario al movimiento.

\textbf{Salida:} sumar centrado, ráfagas, buffets y suelo $\rightarrow$ limitar a $\pm220/255$ $\rightarrow$ aplicar polaridad de motor $\rightarrow$ multiplicar por 255, redondear y limitar a $\pm220$ $\rightarrow$ puente H.
\end{minipage}}
\end{center}
El puente del PC \textbf{no transmite una fuerza de damper ni una velocidad de yugo}: cuando corresponde, transmite posiciones. La placa calcula la velocidad y ambas resistencias. El ciclo de control tiene un mínimo configurado de 2 ms; el cálculo emplea el intervalo real, no supone que todas las iteraciones duren exactamente 2 ms.

\subsection{Variables de SimConnect y condiciones de habilitación}
\label{sec:damper_simvars}
\subsubsection{Entradas que participan directamente en el cálculo}
En \code{YOKE_USE_SIMULATION=1}, se usan las siguientes variables solicitadas por \code{SimConnectBridge/main.cpp} como \code{SIMCONNECT_DATATYPE_FLOAT64}:
\begin{center}
\begin{tabular}{lll}
\toprule
Variable & Unidad pedida & Uso\\
\midrule
\code{YOKE X POSITION} & \code{Position} & Posición de roll $x_R$\\
\code{YOKE Y POSITION} & \code{Position} & Posición de pitch $x_P$\\
\bottomrule
\end{tabular}
\end{center}
El puente exige valores finitos dentro de $[-1,1]$. Son posiciones de mando del simulador; no son el ángulo de alabeo, el cabeceo del avión, su velocidad angular ni sus aceleraciones. El sentido físico debe contrastarse con el eje realmente asignado; el proyecto propone roll a la derecha positivo y tirar en pitch positivo.

En modo potenciómetros o encoders, \textbf{ninguna SimVar proporciona la posición ni la velocidad utilizadas por damper y fricción}. Sin embargo, el puente sigue solicitando y validando las posiciones de SimConnect y el resto de la trama. Por eso una entrada inválida de esa trama puede desarmar la salida aunque el sensor local funcione.

\subsubsection{Datos que habilitan la salida, pero no multiplican el damper}
La rama común requiere \code{SIMCONNECT_ADDITIONS_ENABLED}, trama válida y armada con antigüedad menor de 300 ms, y conexión USB CDC comprobada mediante \code{Serial.dtr()}. El puente condiciona el armado a ejecución del simulador, ausencia de pausa, \code{IS SLEW ACTIVE=0}, telemetría reciente y opción de salida habilitada. Los eventos \code{Pause} y \code{Sim} participan en ese estado; no son coeficientes de resistencia.

\code{DYNAMIC PRESSURE}, \code{AIRSPEED INDICATED}, los vientos, flaps, tren y spoilers \textbf{no aparecen en la fórmula de damper}. La presión sí modifica el resorte que se suma con él. \code{SIM ON GROUND} separa otros efectos, pero no exige estar en vuelo para estas resistencias: pueden actuar en tierra y en aire. El menú de calibración anula ambos comandos de motor mientras está activo. En consecuencia, los efectos locales actuales también necesitan la habilitación de SimConnect; no son un modo autónomo de resistencia con el PC desconectado.

\subsection{Adquisición de posición: los tres modos}
\label{sec:adquisicion_movimiento}
\subsubsection{Selector común y configuración predeterminada}
\code{YOKE_USE_SIMULATION}, en \code{mis_defines.h}, selecciona una única fuente para ambos ejes y para los efectos que utilizan movimiento:
\begin{center}
\begin{tabular}{cll}
\toprule
Valor & Fuente & Característica principal\\
\midrule
0 & Encoders & Posición relativa al arranque\\
1 & SimConnect & Posición del mando virtual recibida por USB\\
2 & Potenciómetros & Posición absoluta según calibración ADC\\
\bottomrule
\end{tabular}
\end{center}
El valor predeterminado inspeccionado es \textbf{2}. No es un selector del menú en tiempo de ejecución. El script \code{Build-Pico2.ps1} permite seleccionar encoders o potenciómetros mediante sus opciones correspondientes; el modo 1 depende de la definición de compilación. No hay mezcla automática de un sensor para el resorte y otro para el damper, ni selección independiente de fuente por eje. Tampoco se comprobó qué binario está cargado físicamente.

\subsubsection{Modo 1: SimConnect, paquete y reloj de las muestras}
La solicitud al simulador se realiza por cuadro de simulación. El puente limita el envío activo a intervalos de al menos 20 ms: como máximo unos 50 envíos por segundo en esa ruta, y menos si no hay cuadros/datos suficientes. No equivale a una garantía de muestreo perfectamente periódico.

Para cada eje, el puente transmite la posición normalizada como float32, en bytes 82--85 para roll y 86--89 para pitch. El receptor valida rango $[-1,1]$, CRC, secuencia y marca temporal.

En SC8, \code{YokeMotion} recibe \code{rx.raw.timeMs*1000UL}: captura de la muestra en Windows, convertida a microsegundos. No es el tiempo de llegada de los bytes a la Pico ni el reloj de física del simulador. Las variaciones de entrega desde SimConnect aún afectan el muestreo, pero agrupar bytes en USB no introduce pasos ficticios. Los protocolos anteriores usan el instante local \code{rx.last}.

Entre paquetes el firmware reutiliza el mismo instante. Entonces $\Delta t=0$ y el estimador no se actualiza. Esto evita interpretar cada iteración local sin paquete nuevo como una muestra de velocidad nula. También implica que, si dejan de llegar paquetes, la velocidad simulada anterior puede mantenerse hasta que una nueva muestra la cambie o la habilitación expire. El contador de reposo no avanza sólo por ejecutar el bucle; necesita nuevos tiempos de muestra. Al perder la habilitación, el estimador de SimConnect se reinicia.

Si falta \code{yokeAvailable}, el modo simulado usa velocidad cero y deshabilita el centrado dependiente de esa posición. \textbf{No cambia automáticamente a un potenciómetro o encoder que pudiera estar conectado.}

La posición virtual puede incorporar curvas, zonas muertas, asignaciones o comportamiento del addon distintos del movimiento físico. No se debe asumir que el movimiento visual producido por trim o piloto automático equivale siempre a movimiento real del eje: hay que comprobar qué publica ese avión. El sketch actual tampoco transmite estos ejes a Prepar3D; la asignación de entrada del simulador es otra parte de la instalación.

\subsubsection{Modo 2: potenciómetros, ADC y promedio}
En Pico se utilizan GP26 para roll y GP27 para pitch, con ADC de 12 bits: valores enteros de 0 a 4095. Son números GPIO, no posiciones físicas del conector. La electrónica de entrada corresponde a 3.3 V. \code{AxisPotentiometers.h} descarta una primera conversión al cambiar de canal y promedia cuatro conversiones por defecto:
\begin{equation}
A=\left\lfloor\frac{a_1+a_2+a_3+a_4+2}{4}\right\rfloor.
\end{equation}
Los $a_i$ son lecturas ADC; el $+2$ permite redondear el promedio entero. Estas cuatro conversiones se toman dentro de la actualización del canal: no constituyen un promedio de las últimas cuatro iteraciones del control.

\subsubsection{Histéresis móvil y normalización de los potenciómetros}
\code{PotentiometerPosition.h} conserva un valor $H$ que sólo se mueve cuando el promedio $A$ se aleja más de $h=4$ cuentas:
\begin{equation}
H_n=\begin{cases}
A_n-h,&A_n>H_{n-1}+h,\\
A_n+h,&A_n<H_{n-1}-h,\\
H_{n-1},&\text{en otro caso}.
\end{cases}
\end{equation}
Inicialmente $H=A$. Es una banda que sigue a la señal: reduce pequeñas oscilaciones del ADC, pero introduce holgura al invertir el movimiento. Una deriva lenta acumulada termina saliendo de la banda y sí se detecta. No debe confundirse con una zona muerta exclusivamente alrededor del centro mecánico.

Sean $N$, $C$ y $P$ las lecturas calibradas del extremo negativo, centro y extremo positivo. Si el ADC crece hacia el extremo positivo:
\begin{equation}
x=\sat_{[-1,1]}\!\left\{\begin{cases}
\dfrac{H-C}{P-C},&H\ge C,\\[3pt]
\dfrac{H-C}{C-N},&H<C.
\end{cases}\right.
\end{equation}
Por ejemplo, con $N=0$, $C=2048$ y $P=4095$, una lectura filtrada $H=2457$ da $x=409/2047\simeq0.1998$. El código también admite ADC decreciente: detecta de qué lado del centro está el extremo positivo y utiliza los denominadores con su signo. No hay que invertir necesariamente el cableado para definir el sentido.

El menú de extremos calcula el centro como el punto medio entero entre los extremos capturados; \textbf{no realiza una tercera captura independiente del centro mecánico}. Por eso un mecanismo asimétrico puede requerir mejorar ese procedimiento. Los defaults son 0/2048/4095, mientras que el punto medio entero de 0 y 4095 es 2047: no hay que confundir default con valor recalculado. La calibración persistida valida un rango de al menos 64 cuentas, pero pasar esa validación no garantiza buena resolución. Un recorrido ADC pequeño amplifica tanto el movimiento como el ruido al normalizarlo.

Al guardar, el sketch detiene las salidas y reinicia los estimadores locales y los efectos. Se evita así interpretar el cambio de calibración o la pausa de escritura como un movimiento físico. La calibración de potenciómetros y las ganancias tienen registros separados.

\subsubsection{Modo 0: encoders y normalización de cuentas}
\code{AxisEncoders.h} pone ambos contadores en cero al comenzar. Por tanto, el yugo debe estar en el centro de referencia al arrancar: no hay búsqueda de origen ni uso del índice Z. Si se arranca desplazado, se desplaza la referencia del resorte. El damper sólo necesita diferencias de posición, pero también queda afectado cuando esa referencia errónea hace alcanzar el recorte de posición.

Con $c$ cuentas ya corregidas por \code{ROLL_ENCODER_DIRECTION} o \code{PITCH_ENCODER_DIRECTION}, la normalización es:
\begin{equation}
x=\sat_{[-1,1]}\!\left\{\begin{cases}
c/C_+,&c\ge0,\\
c/|C_-|,&c<0.
\end{cases}\right.
\end{equation}
Los límites actuales son $C_\pm=\pm1000$ en roll y $\pm2546$ en pitch. El comentario del proyecto supone encoder de 1000 PPR leído en cuadratura x4, roll directo de $\pm90^\circ$, y una geometría estimada para pitch. \textbf{Las 2546 cuentas de pitch deben reemplazarse por mediciones del mecanismo}; tampoco se debe confundir PPR del fabricante con cuentas efectivas después de la decodificación.

En Pico, \code{PicoEncoder.h} utiliza los cambios de ambos canales, una tabla de transiciones de cuadratura y lectura protegida del contador. Los GPIO son 0/1 para roll y 2/3 para pitch. Una transición inválida no reconstruye pasos perdidos. En la rama AVR se utiliza la biblioteca PJRC Encoder y otro cableado: roll 0/1 y pitch 7/4; no se debe trasladar sin revisión la configuración de una placa a otra.

El signo del encoder, la escala de cuentas y la polaridad del motor son tres configuraciones distintas. Invertir el motor no corrige un centro o una escala de sensor mal calibrados. En potenciómetros, el sentido se obtiene de los extremos calibrados y no de las macros de dirección del encoder.

\subsubsection{Comparación práctica para damper y fricción}
\begin{longtable}{p{0.16\linewidth}p{0.36\linewidth}p{0.39\linewidth}}
\toprule
Fuente & Ventaja para este cálculo & Limitación que conviene medir\\
\midrule
SimConnect & Permite demostrar el algoritmo con el mando virtual; no requiere sensor local en esta rama. & Más etapas y retardo; intervalos variables, posición virtual y posibles curvas; velocidad retenida entre paquetes.\\
Potenciómetro & Posición local absoluta después de calibrar; no exige centrar al encender. & Ruido ADC, desgaste, rango aprovechado, promedio e histéresis; centro real frente al punto medio.\\
Encoder & Medición local incremental sin conversión analógica de posición; resolución determinada por cuentas y transmisión. & Centro relativo al arranque, pasos discretos a baja velocidad, posibles cuentas perdidas y escala mecánica.\\
\bottomrule
\end{longtable}
Si se busca resistir el movimiento de la mano, el sensor local mide de forma más directa esa variable. Esto no demuestra automáticamente que cualquier montaje local sea mejor: el ruido, la resolución y la transmisión deben medirse. Cambiar de fuente puede exigir reajustar el filtro y las ganancias aun manteniendo las mismas ecuaciones.

\subsection{Cómo se calcula la velocidad común}
\label{sec:velocidad_comun}
\subsubsection{Derivada discreta y unidades}
Para cada eje y cada nueva muestra:
\begin{align}
\Delta t_n&=(t_n-t_{n-1})\,10^{-6},\\
v_{\rm raw,n}&=\sat_{[-10,10]}\left(\frac{x_n-x_{n-1}}{\Delta t_n}\right).
\end{align}
Los tiempos entran en microsegundos y $10^{-6}$ los convierte a segundos. La división es la definición de velocidad media entre dos posiciones. Como $x$ es adimensional, la unidad es \textbf{recorridos centro--extremo por segundo}. Ir del centro a un extremo en un segundo significa velocidad media de magnitud 1; recorrer de $-1$ a $+1$ en un segundo significa 2.

El límite 10 es una elección del firmware para acotar la derivada. No es una velocidad máxima medida del avión. Además, se recorta la posición a $[-1,1]$ \textbf{antes} de derivar. Si el mecanismo continúa fuera del rango calibrado y la posición ya está recortada, puede resultar $\Delta x=0$: el damper no constituye un detector de tope físico.

\subsubsection{Filtro de primer orden e inversión de dirección}
Con $\tau=0.040$ s se calcula:
\begin{equation}
\beta_n=\frac{\Delta t_n}{\tau+\Delta t_n},\qquad
\widehat v_n=\widehat v_{n-1}+\beta_n(v_{\rm raw,n}-\widehat v_{n-1}).
\label{eq:damper_filtro}
\end{equation}
$\beta$ dice qué parte de la diferencia se incorpora en cada muestra. Esta expresión es la actualización de un filtro de primer orden mediante discretización implícita de $\tau\dot{v}+v=v_{\rm raw}$. Suaviza la derivada ruidosa a costa de retrasar la respuesta. A 2 ms, $\beta=2/42\simeq0.04762$; a 20 ms, $\beta=20/60=1/3$. El coeficiente cambia con el tiempo real.

Antes de aplicar esa ecuación, si $v_{\rm raw,n}\widehat v_{n-1}<0$, el código pone la velocidad anterior en cero. Así, una inversión detectada no conserva una cola del signo anterior. Es un tratamiento especial: el estimador completo no es solamente un filtro lineal. Cuando $v_{\rm raw}=0$, el producto es cero y esta condición de inversión no se activa.

\subsubsection{Reposo, primera muestra y pausas}
El contador de reposo acumula tiempo cuando $|v_{\rm raw}|\le0.02$ y se pone en cero cuando se supera ese umbral. Al alcanzar 40000 microsegundos, fija $\widehat v=0$. Esto elimina una cola residual que un filtro exponencial puro tardaría en extinguir.

La primera muestra sólo establece posición y tiempo; su velocidad es cero. Una entrada no finita reinicia el estimador. Un intervalo entre muestras superior a 250 ms también vuelve a iniciar con velocidad cero. Un intervalo nulo no hace nada. En la rama simulada, además, \code{Effects.update} reinicia sus estimadores al deshabilitarse, al iniciar y si el intervalo de actualización de efectos supera 50 ms. Los estimadores locales siguen actualizándose desde el bucle aunque la salida de efectos esté desarmada; no se deben atribuirles todos los reinicios de la rama simulada.

\subsubsection{Resolución: por qué un pequeño paso puede importar}
Con muestras locales separadas 2 ms, una cuenta de roll de $1/1000$ produce $v_{\rm raw}=0.5$. Si la velocidad previa era cero, el primer valor filtrado es $0.5/21\simeq0.02381$, apenas por encima de la zona muerta. Una cuenta de pitch de $1/2546$ da $v_{\rm raw}\simeq0.19639$ y $\widehat v\simeq0.00935$, por debajo de ella.

Un paso ADC de $1/2047$ en el lado positivo, una vez superada la histéresis, produce aproximadamente $v_{\rm raw}=0.24426$ y $\widehat v=0.01163$ en esa primera actualización. En SimConnect, un paso de $0.0001$ en 20 ms equivale a $0.005$ antes del filtro. Estos ejemplos parten de reposo y de un único paso; no son umbrales universales para movimiento continuo.

En un encoder lento, los pasos pueden estar separados por intervalos sin cambios mayores de 40 ms. El reinicio de reposo puede entonces generar pequeñas respuestas intermitentes en vez de una velocidad continua. Aumentar la ganancia no recupera información que el sensor o el estimador eliminaron. Conviene resolver primero resolución y estimación.

\subsection{Fórmula exacta del damper y significado de cada parte}
\subsubsection{Zona muerta y ganancia del eje}
Sea $v=\sat_{[-10,10]}(\widehat v)$ la velocidad que recibe la función. Se define:
\begin{equation}
m=\max(|v|-v_d,0),\qquad v_d=0.02.
\label{eq:movimiento_util}
\end{equation}
$m$ es la magnitud de movimiento útil. \textbf{Se resta el umbral}: con $|v|=0.10$, se utiliza $m=0.08$, no 0.10. La resistencia empieza continuamente desde cero al superar el umbral.

Sean $G$ la ganancia general y $s_R=0.5$, $s_P=1$. El factor aplicado es $H_a=\sat_{[0,1]}(G s_a)$. Con $G=0.90$, $H_R=0.45$ y $H_P=0.90$. Esta reducción fija no procede de la relación física entre milímetros de pitch y grados de roll.

\subsubsection{Damper aislado}
Si sólo está habilitado el damper:
\begin{equation}
\boxed{u_{d,a}=-\operatorname{sgn}(v_a)\,H_a\,\min(Bm_a,L)}
\qquad B=0.120,\quad L=0.400.
\label{eq:damper_aislado}
\end{equation}
$B$ es la ganancia efectiva de damper, $L$ limita la resistencia previa a la ganancia general, y $u_{d,a}$ es el comando normalizado adicional. Si el efecto se desactiva desde el menú, su ganancia efectiva es cero. La máscara de compilación también puede excluirlo.

Para $v>0$ el resultado es negativo; para $v<0$ es positivo. En $v=0$, $m=0$ y el resultado es cero independientemente de cómo se implemente el signo. No depende de $x$ directamente: dos posiciones distintas con igual velocidad estimada generan el mismo damper mientras la medición no esté recortada.

\subsubsection{Límite compartido con fricción}
Con ambos habilitados, el código realmente ejecuta:
\begin{align}
R_d&=Bm, & R_f&=F_c\frac{m}{v_0+m},\\
R&=\min(R_d+R_f,L), &
u_{\rm mov,a}&=-\operatorname{sgn}(v_a)H_aR,
\label{eq:movimiento_conjunto}
\end{align}
con $F_c=0.030$ y $v_0=0.200$. \textbf{El límite se aplica a la suma antes de multiplicar por $H_a$.} No se limita cada término por separado para sumarlos después. Por encima del límite ya no hay una contribución independiente observable de cada uno; sólo queda la resistencia total limitada.

El máximo de esta rama con las ganancias generales actuales es 0.36 en pitch y 0.18 en roll: aproximadamente 92 y 46 cuentas PWM si actúa sola. No es el máximo de todos los efectos juntos, que sigue siendo 220. El damper aislado alcanza el límite a $|v|=0.02+0.4/0.12\simeq3.3533$. La fricción hace alcanzar antes el límite conjunto.

\subsubsection{Procedencia de los valores}
La dependencia proporcional a velocidad es el modelo viscoso. La diferencia discreta sale de la definición de derivada; el filtro es una aproximación de primer orden. Los números 0.12, 0.02, 40 ms, 10 y 0.40 son \textbf{parámetros del proyecto}, no constantes universales ni valores identificados para un avión concreto. La relación fija roll/pitch también es una elección local. Su presencia se verifica en \code{mis_defines.h}, \code{GainSettings.h} y \code{SimEffects.h}; la comparación histórica se detalla más adelante.

\subsection{Ejemplos numéricos y conversión a PWM}
\subsubsection{Velocidad ya estabilizada, damper aislado}
Con $v=+1$, $m=1-0.02=0.98$. Se calcula $R_d=0.12\cdot0.98=0.1176$. Como es menor que 0.4, no se limita:
\begin{align}
u_{d,P}&=-0.90\cdot0.1176=-0.10584,& \mathrm{PWM}_P&=-27,\\
u_{d,R}&=-0.45\cdot0.1176=-0.05292,& \mathrm{PWM}_R&=-13.
\end{align}
Para $v=-1$, los mismos valores tienen signo positivo. Los PWM de estos ejemplos suponen polaridad de motor $+1$, ausencia de otros efectos y redondeo del driver.

\begin{center}
\begin{tabular}{rrrrr}
\toprule
$|v|$ & $m$ & $\min(0.12m,0.4)$ & $|\mathrm{PWM}_P|$ & $|\mathrm{PWM}_R|$\\
\midrule
0.00 & 0.00 & 0.0000 & 0 & 0\\
0.02 & 0.00 & 0.0000 & 0 & 0\\
0.10 & 0.08 & 0.0096 & 2 & 1\\
0.22 & 0.20 & 0.0240 & 6 & 3\\
0.50 & 0.48 & 0.0576 & 13 & 7\\
1.00 & 0.98 & 0.1176 & 27 & 13\\
2.00 & 1.98 & 0.2376 & 55 & 27\\
4.00 & 3.98 & 0.4000 & 92 & 46\\
\bottomrule
\end{tabular}
\end{center}
La tabla usa velocidad \emph{estimada}, no necesariamente la velocidad instantánea real del mecanismo. La zona muerta, el filtro y el redondeo PWM son tres causas diferentes de que un movimiento pequeño no produzca una salida apreciable.

\subsubsection{Ejemplo desde dos muestras de SimConnect}
Supóngase que la muestra previa es $x=0.10$, la velocidad previa cero y, 20 ms después, llega $x=0.12$. Tras decodificar la trama:
\begin{align}
v_{\rm raw}&=(0.12-0.10)/0.020=1,\\
\widehat v&=0+\frac{0.020}{0.040+0.020}(1-0)=\frac13,\\
m&=\frac13-0.02\simeq0.313333,\\
u_{d,P}&=-0.90\cdot0.12\cdot0.313333=-0.03384.
\end{align}
El resultado aislado es PWM $-9$ en pitch y $-4$ en roll. No son los $-27/-13$ del ejemplo anterior porque el filtro todavía está comenzando a responder. Si el movimiento continúa a esa velocidad, el valor filtrado se aproxima a 1. Si se detiene, se atenúa y finalmente se anula por el criterio de reposo.

\subsubsection{Suma final, signo eléctrico y ciclo de trabajo}
El comando que llega al driver es:
\begin{equation}
u_a=\sat_{[-220/255,\,220/255]}
\left(u_{c,a}+u_{t,a}+u_{b,a}+u_{\rm mov,a}+u_{g,a}\right).
\end{equation}
Las letras representan centrado, turbulencia, suma de buffets, resistencia conjunta al movimiento y vibración en suelo. Con $d_a\in\{-1,+1\}$ como dirección de motor y \code{PWM_MAX=255}:
\begin{equation}
\mathrm{PWM}_a=\sat_{[-220,220]}\!\left(\operatorname{redondeo}_{\rm sim}(255d_au_a)\right).
\end{equation}
El redondeo es al entero más cercano, con mitades alejándose de cero. PWM positivo activa la salida derecha y deja la izquierda en cero; PWM negativo activa la izquierda con su magnitud y deja la derecha en cero. En Pico se emplean 6/7 para roll y 8/9 para pitch, con habilitaciones 10 y 11. La frecuencia configurada es 490 Hz. Por ejemplo, una magnitud 27 representa $27/255\simeq10.59\%$ de ciclo de trabajo; \textbf{no representa 27 N ni garantiza 10.59\% del torque nominal}. No hay en esta ruta un lazo de corriente que convierta automáticamente comando en fuerza medida.

\subsection{Relación con el centrado y con los otros efectos}
\subsubsection{Suma con el resorte: posición frente a velocidad}
Usando el resorte ya documentado:
\begin{equation}
u_{c,a}=-Gg_ak_fx_a,\qquad
k_{\rm objetivo}=r+(1-r)Q,
\end{equation}
con $g_R=0.4$, $g_P=0.8$, $r=0.15$ y presión normalizada $Q$ filtrada a través de la rigidez con 100 ms. Sin fricción ni saturación, el conjunto se aproxima a un resorte más amortiguador:
\begin{equation}
u_a=-Gg_ak_fx_a-H_aB\operatorname{sgn}(v_a)\max(|v_a|-0.02,0).
\end{equation}
No es un seguidor de trayectoria con consigna móvil ni un PID completo; no tiene término integral en esta cadena. El centro ideal sigue siendo cero.

\begin{longtable}{p{0.30\linewidth}p{0.29\linewidth}p{0.32\linewidth}}
\toprule
Estado del yugo & Centrado & Damper\\
\midrule
Quieto fuera del centro & Mantiene comando restaurador. & Se anula al extinguirse la velocidad estimada.\\
Se aleja del centro: $xv>0$ & Intenta devolverlo. & Se opone al alejamiento: mismo sentido que el resorte.\\
Regresa al centro: $xv<0$ & Favorece el regreso. & Frena el regreso: sentido opuesto al resorte.\\
Atraviesa $x=0$ con $v\ne0$ & Es cero en ese instante. & Sigue oponiéndose al movimiento.\\
\bottomrule
\end{longtable}
Con $Q=1$, rigidez estabilizada $k_f=1$, pitch $x=0.20$ y $v=+1$, el resorte da $-0.144$ y el damper $-0.10584$: total $-0.24984$, PWM $-64$. Si el mismo mando regresa con $v=-1$, total $-0.03816$, PWM $-10$. El damper está reduciendo la rapidez del retorno, no cambiando el centro del resorte.

\subsubsection{Presión y amortiguación: dependencia indirecta}
Cuando aumenta $Q$, aumenta la rigidez objetivo y el retorno puede ser más rápido. El damper actual no escala automáticamente con esa presión: $B$ permanece en su ajuste. Su magnitud sólo aumentará si aumenta la velocidad medida. Por eso una combinación agradable a baja rigidez puede resultar poco amortiguada a alta rigidez.

En un modelo mecánico ideal calibrado aparece el criterio $b_{\rm crítico}=2\sqrt{kJ}$ para una coordenada angular, siendo $k$ rigidez física y $J$ inercia. \textbf{No se puede sustituir sin conversión el $B=0.12$ de este firmware en esa ecuación}: faltan unidades físicas, transmisión, inercia identificada y relación comando--torque. Aquí sirve como explicación de por qué rigidez e inercia importan, no como receta numérica de calibración.

\subsubsection{Interacción con turbulencias, buffets y suelo}
El damper se suma; no multiplica las ondas por un factor de velocidad ni necesita que haya una fuerza base. Si una vibración mueve físicamente el yugo, ese movimiento genera resistencia y puede atenuar la vibración percibida. Si el mando no se mueve, la velocidad se anula y el damper deja de aportar, aunque el generador de vibraciones siga calculando una señal.

El código tampoco distingue si el movimiento lo produjo la mano, el resorte o una perturbación del motor: responde a la misma velocidad. La suma final puede saturar y recortar ondas. Aumentar mucho el damper puede hacer menos perceptibles los otros efectos, por oposición mecánica o por saturación compartida, aun cuando sus ganancias no cambien.

\subsubsection{Qué demuestra la oposición al movimiento}
Para ganancias no negativas, $u_{\rm mov}v\le0$: la resistencia calculada se opone a la \textbf{velocidad estimada}. Esto es una propiedad algebraica de esos términos, no de la suma completa con ráfagas o resorte. Tampoco demuestra estabilidad o pasividad física del mecanismo: intervienen retardo, muestreo, polaridad y correspondencia del sensor con el movimiento real. En especial, la velocidad virtual de SimConnect puede diferir de la de la mano.

\subsection{Procedencia: proyecto actual, original alemán y referencias}
\label{sec:origen_damper_friccion}
\subsubsection{Antecedente comprobable en el proyecto alemán}
Se inspeccionó la copia local en \code{PROEYCTO ORIGINAL/Arduino_FFB_Yoke}, situada fuera de la carpeta modificada. En \code{joystick.ino::UpdateEffects}, calcula diferencias de cuentas y divide por diferencias de tiempo en milisegundos, usando variables enteras. Asigna esas velocidades a \code{damperVelocity} y también a \code{frictionPositionChange}, a pesar del nombre de este último campo.

En \code{src/Joystick.cpp}, los casos \code{USB_EFFECT_DAMPER} y \code{USB_EFFECT_FRICTION} llaman a \code{ConditionForceCalculator} con una métrica normalizada, aplican ganancias y filtros específicos. Es una arquitectura de efectos condicionales HID con parámetros recibidos por ese sistema. \textbf{No es la fórmula actual completa}. No corresponde atribuirle el filtro común de 40 ms, la zona muerta 0.02, los límites actuales ni las ganancias actuales sólo porque ambos proyectos tengan efectos con el mismo nombre.

\subsubsection{TelemFFB: antecedente conceptual, no esta ecuación exacta}
En la copia local \code{VPforce-TelemFFB-refactor/telemffb/sim/base/FFBForcesMixIn.py}, el método \code{ac_update_ffb_forces} limita intensidades configuradas a $[0,1]$ e inicia efectos mediante \code{damper(force,force)} y \code{friction(force,force)}. Esa rutina de overrides no lee claves de telemetría para calcular una derivada; delega el efecto al dispositivo a través de su interfaz de efectos.

Por tanto, TelemFFB sirve como antecedente de disponibilidad y configuración de esos efectos, pero \textbf{no se ha demostrado que proporcione la ecuación local $F_cm/(0.2+m)$ ni el coeficiente $B=0.12$}. La implementación actual calcula explícitamente las resistencias en la placa y las convierte a PWM.

\subsubsection{Evolución local que explica documentación contradictoria}
La instantánea \code{tesis/auditoria_c182rg/antes_mejoras/SimEffects.h.txt} ya contiene los términos $Bm$ y $F_cm/(v_0+m)$. Sin embargo, multiplicaba su suma por:
\begin{equation}
\min\left(\frac{|u_{\rm base}|}{0.10},1\right)
\end{equation}
y evitaba que la corrección invirtiera el signo de la fuerza base. Por ello, si la fuerza base era cero, la corrección también era cero. \code{CAMBIOS_REALISMO.md} documenta el cambio hacia cargas independientes; la función actual resta la resistencia con su signo sin esa dependencia.

El comentario actual que menciona regularización conservada del ``original firmware'' debe interpretarse con cautela: \textbf{se comprobó su presencia en la versión local anterior; no se comprobó su procedencia del proyecto alemán}. No hay evidencia suficiente para asignarle un autor o una página matemática original. Las referencias históricas a \code{YOKE_MOTION_FORCE_BLEND=0.100} describen el modelo anterior: esa constante ya no participa en el cálculo. Asimismo, \code{VELOCITY_FILTER_TAU_S=0.018} y la clase \code{AxisMotion} no son el estimador activo de esta cadena: se utiliza \code{YokeMotion} con 0.040 s.

\subsubsection{Fuentes enlazadas y alcance de la atribución}
\begin{enumerate}
\item \href{https://www.prepar3d.com/SDKv3/LearningCenter/utilities/variables/simulation_variables.html}{SDK de Prepar3D: Simulation Variables}. Referencia de nombres y unidades; las solicitudes concretas se verificaron en el puente local.
\item \href{https://github.com/walmis/VPforce-TelemFFB/blob/refactor/telemffb/sim/base/FFBForcesMixIn.py}{TelemFFB: FFBForcesMixIn.py}. Comparación con la copia local de configuración de efectos damper/friction; la rama remota puede evolucionar.
\item \href{https://github.com/Ultrawipf/OpenFFBoard/blob/master/Firmware/FFBoard/Src/EffectsCalculator.cpp}{OpenFFBoard: EffectsCalculator.cpp}. Distingue posición para resorte y velocidad para damper, y suma efectos antes de recortar la salida. Es la referencia arquitectónica citada por \code{CAMBIOS_REALISMO.md}; no demuestra copia de sus coeficientes o de su ley de fricción.
\item Código local actual y copias históricas identificadas arriba. Para las fórmulas numéricas del proyecto son la fuente primaria, por encima de documentos descriptivos antiguos.
\end{enumerate}
Consulta de las referencias web del SDK y OpenFFBoard: 8 de octubre de 2026. Las comparaciones de TelemFFB y del original alemán se realizaron contra las copias disponibles en el espacio de trabajo.

\subsection{Qué ajustar y cómo interpretar el resultado}
\subsubsection{Ganancias disponibles sin recompilar}
En el menú, \code{Damper} ajusta $B$ entre 0 y 1 en pasos de 0.01 y permite desactivar el efecto. \code{Ganancia total} ajusta $G$ entre 0 y 1 en pasos de 0.05; afecta también a los demás efectos. Cambiar $B$ es la forma específica de modificar el damper. No hay dos entradas separadas de damper para roll y pitch: su relación de mando es actualmente 1:2.

\subsubsection{Parámetros que requieren cambiar código y recompilar}
\begin{longtable}{p{0.29\linewidth}p{0.16\linewidth}p{0.46\linewidth}}
\toprule
Parámetro & Default & Consecuencia del ajuste\\
\midrule
\code{YOKE_MOTION_FILTER_TAU_S} & 0.040 s & Más alto suaviza más y retrasa; más bajo responde antes y deja pasar más cuantización/ruido. Afecta también a fricción.\\
\code{YOKE_MOTION_DEAD_SPEED} & 0.02 & Más alto reduce respuesta a movimientos pequeños; también cambia el criterio de reposo.\\
\code{YOKE_MOTION_MAX_SPEED} & 10 & Acota picos de derivada; no sustituye un sensor bien calibrado.\\
\code{YOKE_MOTION_CORRECTION_LIMIT} & 0.40 & Límite conjunto de damper y fricción antes de $G$ y del factor de eje.\\
Reposo en \code{YokeMotion.h} & 40000 $\mu$s & Más tiempo conserva colas/velocidad entre pasos; menos puede cortar movimientos lentos discretos.\\
Factor roll en \code{SimEffects.h} & 0.5 & Modifica ambos efectos de movimiento en roll respecto de pitch.\\
\code{YOKE_POT_SAMPLES}, histéresis & 4 y 4 & Modifican el ruido y el movimiento pequeño que llega desde los potenciómetros.\\
\bottomrule
\end{longtable}
Ningún ajuste de \code{YOKE_MOTION_FORCE_BLEND} corrige la sensación actual porque ya no se usa. El límite de salida PWM tampoco es una ganancia específica del damper.

\subsubsection{Orden de calibración propuesto}
\begin{enumerate}
\item Medir centro, extremos, sentido y rango útil del sensor elegido. Registrar posición normalizada al mover el mecanismo a puntos conocidos; corregir escala antes de compensarla con ganancia.
\item Con el mando quieto, observar posición y velocidad. Si existe ruido, resolver ADC, histéresis, resolución o filtrado antes de subir resistencia.
\item Aislar damper desactivando fricción, centrado y perturbaciones desde el menú. Comparar movimientos lentos y rápidos en ambos sentidos. Debe oponerse al movimiento y desaparecer al detenerse.
\item Activar el resorte y ensayar retornos desde varias posiciones y presiones. Aumentar $B$ gradualmente si hay rebote; reducirlo si el retorno queda excesivamente lento. Registrar cuándo se alcanza saturación.
\item Incorporar fricción y, después, perturbaciones. Comprobar que la resistencia no oculta los efectos que se quieren percibir. Repetir por eje: el mismo comando normalizado no implica el mismo esfuerzo físico.
\end{enumerate}
Es un procedimiento propuesto, no un ensayo ya realizado. Para movimientos rápidos conviene registrar muestras a mayor frecuencia que el diagnóstico serie actual: éste emite cada 100 ms alternando ejes, por lo que cada eje se observa aproximadamente cada 200 ms. Sirve para tendencias, pero no reconstruye transitorios de 40 ms.

\subsubsection{Mejoras que requieren desarrollo e identificación}
Se podrían incorporar ganancias separadas por eje, captura real del centro del potenciómetro, estimación de velocidad de encoder basada en tiempo entre flancos, y marcas de tiempo de adquisición en la telemetría. Otra opción sería programar $B$ en función de rigidez, pero exige definir y validar esa relación; no está implementada. Para obtener valores físicos, hacen falta mediciones de fuerza/torque frente a PWM, posición y velocidad, además de la transmisión mecánica. El ajuste empírico actual no demuestra reproducción de un modelo de avión concreto.

\newpage
\section{FRICCION}
\label{sec:friccion}

\subsection{Qué representa y diferencia con damper}
La fricción intenta reproducir una resistencia de rozamiento al desplazar el mando. A diferencia del damper, cuya magnitud crece proporcionalmente a velocidad útil hasta el límite, la fricción implementada crece al comenzar a moverse y se aproxima gradualmente a una magnitud casi constante. La idea es aportar ``roce'' sin hacer que toda esa resistencia siga aumentando linealmente a altas velocidades.

Es una \textbf{aproximación suave de fricción de Coulomb}, no un modelo completo de rozamiento mecánico. En la fricción ideal de Coulomb en movimiento, la magnitud es constante y cambia de signo al cambiar la dirección. El salto brusco en cero puede resultar poco conveniente con sensores discretos. Este proyecto lo reemplaza por una fracción continua y una zona muerta.

\subsubsection{Lo que sucede cuando el mando queda quieto}
Cuando la velocidad estimada llega a cero, el comando de fricción es cero. No hay una fuerza de adherencia capaz de sostener por sí sola el mando desplazado contra el resorte. Tampoco se calcula un umbral de desprendimiento, memoria de deformación, temperatura, desgaste o una curva de Stribeck. Si el yugo real queda detenido fuera del centro, puede ser por su rozamiento físico, holgura, falta de torque o redondeo de salida; no es evidencia de adherencia estática generada por esta ecuación.

\subsection{Entradas y recorrido desde la posición hasta PWM}
\subsubsection{Variables de simulación utilizadas}
La fricción utiliza exactamente la misma velocidad estimada que el damper. En modo 1, deriva \code{YOKE X POSITION} para roll y \code{YOKE Y POSITION} para pitch, solicitadas en unidad \code{Position} y transmitidas como float32 en SC8. No utiliza una SimVar llamada ``fricción'', no deriva la actitud del avión y no emplea la velocidad respecto del suelo.

En modo 2 recibe la posición normalizada del potenciómetro; en modo 0, las cuentas normalizadas del encoder. Las ecuaciones de adquisición, calibraciones, tiempos de muestra y filtros están desarrolladas en las subsecciones \ref{sec:adquisicion_movimiento} y \ref{sec:velocidad_comun}. Se aplican íntegramente a esta sección: la fricción \textbf{no tiene un segundo sensor ni un filtro de velocidad propio}.

\subsubsection{Cadena de cálculo y habilitación}
\begin{center}
\fbox{\begin{minipage}{0.91\linewidth}
Posición de la fuente seleccionada $\rightarrow$ $x\in[-1,1]$ $\rightarrow$ derivada con tiempo real $\rightarrow$ velocidad limitada y filtrada $\widehat v$ $\rightarrow$ $m=\max(|\widehat v|-0.02,0)$ $\rightarrow$ $F_cm/(0.2+m)$ $\rightarrow$ suma con damper y límite 0.4 $\rightarrow$ ganancia general/factor de eje $\rightarrow$ signo opuesto a velocidad $\rightarrow$ suma con otros efectos $\rightarrow$ límite final, polaridad, redondeo y PWM.
\end{minipage}}
\end{center}
Se requiere la misma telemetría armada, reciente y conexión USB que para damper. Puede actuar en tierra y en vuelo; no necesita tren extendido, pérdida, spoilers ni centrado activado. Durante pausa/desarmado o menú activo, la salida común se anula. Su coeficiente no cambia directamente con presión dinámica, IAS, AoA o configuración del avión.

\subsection{Fórmula exacta explicada parte por parte}
\subsubsection{Modelo ideal y regularización elegida}
La idea ideal en movimiento es $u_f=-F_c\operatorname{sgn}(v)$. La implementación local suaviza la magnitud mediante:
\begin{equation}
R_f=F_c\frac{m}{v_0+m},\qquad
m=\max(|v|-v_d,0),\quad v_d=0.02,\quad v_0=0.20.
\label{eq:friccion_magnitud}
\end{equation}
\begin{itemize}
\item $v$ es la velocidad estimada y limitada, en recorridos centro--extremo/s. Su signo sólo se utiliza para elegir el sentido de la resistencia.
\item $m$ es la magnitud que queda después de restar la zona muerta. Dentro de ella, la resistencia es exactamente cero.
\item $F_c$ es el coeficiente efectivo de fricción: por defecto 0.030. Determina la magnitud a la que se aproxima la resistencia antes de ganancias y límites.
\item $v_0$ determina qué tan rápido aparece esa resistencia al aumentar la velocidad. Tiene la misma unidad de velocidad que $m$ y evita dividir por cero.
\item El cociente $m/(v_0+m)$ es adimensional, comienza en cero y permanece menor que uno para velocidades finitas con $v_0>0$.
\end{itemize}
Con fricción aislada:
\begin{equation}
\boxed{u_{f,a}=-\operatorname{sgn}(v_a)H_a\min\left(F_c\frac{m_a}{v_0+m_a},L\right)}.
\label{eq:friccion_aislada}
\end{equation}
Se mantienen $H_R=0.45$, $H_P=0.90$ y $L=0.40$ con los defaults. Esta fórmula sirve para aislar el efecto; cuando se habilita damper, se aplica el límite a la suma según la ecuación \ref{eq:movimiento_conjunto}.

\subsubsection{Por qué la mitad no aparece exactamente a velocidad 0.20}
La fracción vale $1/2$ cuando $m=v_0$. Como primero se resta 0.02:
\begin{equation}
|v|_{1/2}=v_d+v_0=0.22.
\end{equation}
Así, a $|v|=0.22$ la resistencia sin ganancias es $F_c/2=0.015$. El comentario que denomina a \code{YOKE_FRICTION_SPEED} ``velocidad para alcanzar media corrección'' se refiere a la velocidad útil $m$, no directamente a la velocidad total $|v|$.

\subsubsection{Comportamiento lento y rápido}
Para $m\ll v_0$:
\begin{equation}
R_f\simeq\frac{F_c}{v_0}m=0.15m.
\end{equation}
Cerca del umbral la fricción suave se parece a una resistencia viscosa adicional. Su pendiente inicial 0.15 es incluso algo mayor que la pendiente 0.12 del damper predeterminado. Por eso una ganancia de fricción aparentemente pequeña puede modificar bastante la sensación al iniciar movimientos.

Para $m\gg v_0$, $R_f\rightarrow F_c$. El crecimiento se aplana, pero no alcanza matemáticamente ese valor a velocidad finita. Como el firmware limita $|v|$ a 10, la magnitud máxima de esta expresión con defaults es:
\begin{equation}
R_{f,\max}=0.03\frac{9.98}{10.18}\simeq0.0294106.
\end{equation}
Tras las ganancias da aproximadamente 0.02647 en pitch y 0.01323 en roll, equivalentes a 7 y 3 cuentas PWM aisladas. El límite común 0.4 no actúa sobre fricción aislada con $F_c=0.03$; sí puede intervenir con ganancias más altas o al sumarse damper.

\subsubsection{Qué cambia al ajustar el parámetro de transición}
Con $F_c$ fijo, reducir $v_0$ hace que la fricción aparezca más rápidamente y se aproxime antes a su techo. Aumentarlo la vuelve más gradual y reduce la resistencia a una misma velocidad moderada. No cambia el valor asintótico $F_c$, pero sí el máximo efectivamente alcanzable dentro del límite de velocidad 10. Un $v_0$ demasiado pequeño aumenta la sensibilidad a ruido y a pasos de cuantización justo fuera de la zona muerta.

\subsection{Ejemplos completos para roll y pitch}
\subsubsection{Fricción aislada a velocidad estabilizada}
Para $v=+1$:
\begin{align}
m&=1-0.02=0.98,\\
R_f&=0.03\frac{0.98}{0.20+0.98}\simeq0.0249153,\\
u_{f,P}&=-0.90R_f\simeq-0.0224237, &\mathrm{PWM}_P&=-6,\\
u_{f,R}&=-0.45R_f\simeq-0.0112119, &\mathrm{PWM}_R&=-3.
\end{align}
Para movimiento negativo cambia el signo. La normalización utiliza el recorrido calibrado de cada eje: igual velocidad normalizada no implica igual velocidad angular o lineal física.

\begin{center}
\begin{tabular}{rrrrr}
\toprule
$|v|$ & $m$ & $R_f$ & $|\mathrm{PWM}_P|$ & $|\mathrm{PWM}_R|$\\
\midrule
0.00 & 0.00 & 0.000000 & 0 & 0\\
0.02 & 0.00 & 0.000000 & 0 & 0\\
0.05 & 0.03 & 0.003913 & 1 & 0\\
0.10 & 0.08 & 0.008571 & 2 & 1\\
0.22 & 0.20 & 0.015000 & 3 & 2\\
0.50 & 0.48 & 0.021176 & 5 & 2\\
1.00 & 0.98 & 0.024915 & 6 & 3\\
2.00 & 1.98 & 0.027248 & 6 & 3\\
10.00 & 9.98 & 0.029411 & 7 & 3\\
\bottomrule
\end{tabular}
\end{center}
El cero PWM de roll a velocidad 0.05 es consecuencia del redondeo de una orden muy pequeña, no de que la fórmula de fricción sea cero. Esto puede percibirse como una zona de baja sensibilidad adicional a la zona muerta explícita.

\subsubsection{Desde muestras nuevas hasta fricción y damper conjuntos}
Retomando el ejemplo $x:0.10\rightarrow0.12$ en 20 ms, desde velocidad previa cero, se obtiene $\widehat v=1/3$ y $m\simeq0.313333$. Entonces:
\begin{align}
R_f&=0.03\frac{0.313333}{0.513333}\simeq0.0183117,\\
u_{f,P}&\simeq-0.0164805, & u_{f,R}&\simeq-0.0082403.
\end{align}
La fricción aislada da PWM $-4$ y $-2$. Si también se activa damper, $R_d=0.0376$ y la resistencia conjunta es aproximadamente 0.0559117, menor que el límite. El comando combinado resulta $-0.0503205$ en pitch y $-0.0251603$ en roll: PWM $-13$ y $-6$.

Se suman comandos de punto flotante y luego se redondea una sola vez al convertir a PWM. Sumar los PWM ya redondeados de una tabla de efectos aislados no reproduce necesariamente la salida final.

\subsubsection{Conversión eléctrica y límites}
Después de sumar los otros efectos se limita $u$ a $\pm220/255$, se aplica la dirección del motor y se calcula el PWM con factor 255 y redondeo simétrico. Son los mismos pasos descritos en DAMPER. El signo selecciona una de las dos ramas del puente H y la magnitud determina su ciclo de trabajo. Por ejemplo, 6 cuentas equivalen a aproximadamente $2.35\%$ de PWM; no aseguran que el motor venza el rozamiento mecánico real. La resistencia física resultante depende de la instalación.

\subsection{Relación con damper, centrado y perturbaciones}
\subsubsection{Cómo se reparten damper y fricción según velocidad}
Antes de saturar, ambos tienen el mismo signo de oposición y sus magnitudes se suman:
\begin{equation}
R_{\rm total}=Bm+F_c\frac{m}{v_0+m}.
\end{equation}
Para $m>0$, la igualdad entre magnitudes ocurre cuando:
\begin{equation}
Bm=F_c\frac{m}{v_0+m}
\quad\Longrightarrow\quad m=\frac{F_c}{B}-v_0.
\end{equation}
Con defaults, $m=0.05$, es decir $|v|=0.07$. Entre el umbral 0.02 y 0.07 domina ligeramente la fricción; por encima domina el damper. Esta comparación supone ambos habilitados y sin saturación. Si se cambian ganancias, el cruce cambia o puede dejar de existir para $m>0$.

Muy cerca del umbral, la pendiente conjunta es $B+F_c/v_0=0.27$. A velocidades elevadas domina el término lineal del damper. Por ello, subir fricción para corregir un retorno rápido con rebote puede alterar demasiado los movimientos lentos sin resolver de forma adecuada la causa.

\subsubsection{Techo compartido y pérdida de sensibilidad a ajustes}
El límite conjunto se alcanza al satisfacer:
\begin{equation}
Bm+F_c\frac{m}{v_0+m}=0.4.
\end{equation}
Multiplicando por $v_0+m$:
\begin{equation}
Bm^2+(Bv_0+F_c-0.4)m-0.4v_0=0.
\end{equation}
La raíz positiva con defaults es aproximadamente $m=3.0985$, o $|v|\simeq3.1185$. A partir de allí la resistencia de movimiento queda en 0.4 antes de ganancia general/eje. Aumentar fricción puede adelantar la entrada en saturación, pero no aumentar la resistencia ya saturada. Este techo de rama es distinto del techo final de todos los efectos y del límite de PWM del driver.

\subsubsection{Relación con el centrado: no desplaza su equilibrio}
La fricción se opone al desplazamiento tanto al alejarse como al regresar al centro. No modifica la fórmula de rigidez ni la presión dinámica. Cuando $xv>0$, refuerza el sentido restaurador; cuando $xv<0$, se opone al regreso. En el centro, si el mando sigue moviéndose, puede seguir activa.

En pitch, con $Q=1$, $k_f=1$, $x=0.20$ y $|v|=1$, el centrado es $-0.144$. Damper y fricción juntos generan magnitud:
\begin{equation}
0.90(0.1176+0.0249153)\simeq0.128264.
\end{equation}
Si $v=+1$, el total es $-0.272264$, PWM $-69$. Si $v=-1$, el total es $-0.0157356$, PWM $-4$. En roll, con la misma posición y velocidad normalizadas, estos ejemplos dan aproximadamente PWM $-35$ y $-2$ respectivamente, porque tanto el resorte predeterminado como las resistencias tienen la mitad del comando de pitch.

Puede invertirse el signo total respecto del resorte sin que el algoritmo esté equivocado. Por ejemplo, en pitch $x=0.05$, $v=-1$, $Q=1$ y rigidez estabilizada, el resorte aporta $-0.036$ y la resistencia $+0.128264$: total $+0.092264$, PWM $+24$. En ese instante el mando regresa hacia cero y el comando neto intenta frenarlo. La versión histórica impedía esa inversión; la actual la permite.

Al detenerse, ambas resistencias se anulan y vuelve a quedar el resorte. Por tanto, la fricción programada no crea por sí sola un nuevo centro de equilibrio ni un rango de posiciones sostenidas por adherencia estática.

\subsubsection{Efectos periódicos, turbulencias y movimiento provocado por el motor}
La fricción puede reducir físicamente la oscilación causada por buffets, suelo o ráfagas porque se opone a la velocidad que éstos produzcan. No reduce algebraicamente su amplitud multiplicándola por una ganancia de fricción: es una suma de comandos. No cambia sus frecuencias ni consulta sus envolventes.

El cálculo no distingue una velocidad causada por la mano de una causada por el motor. Si el sensor detecta movimiento, responde. En modo SimConnect esa relación física puede ser menos directa; una posición virtual fija no produce resistencia aunque el mecanismo se mueva sin reflejarse en esa variable. Asimismo, una posición virtual que cambia puede producir resistencia calculada sin que la mano se haya movido de igual manera.

\subsection{Procedencia de la fórmula y de sus parámetros}
\subsubsection{Qué está demostrado por el código}
La fracción $m/(v_0+m)$ aparece en la instantánea local anterior a las mejoras y en el código vigente. La modificación documentada convirtió esa corrección dependiente de fuerza base en una carga independiente. El número 0.20 y la ganancia 0.03 son parámetros locales; no hay evidencia de que sean mediciones de fricción del Cessna o del mecanismo del usuario.

La comparación completa y los enlaces están en \ref{sec:origen_damper_friccion}. El original alemán utiliza \code{ConditionForceCalculator} para el efecto HID de fricción, mientras TelemFFB configura un efecto del dispositivo mediante \code{friction(force,force)}. Esas evidencias no permiten atribuirles esta fracción racional concreta. OpenFFBoard es una referencia para separación y composición de efectos, no una fuente demostrada de esta ecuación exacta.

\subsubsection{Interpretación física y límite de la simplificación}
El principio de oposición al deslizamiento motiva el signo; la aproximación a una magnitud constante motiva el parecido con Coulomb. La función racional es una elección computacional sencilla para suavizar la transición. No se identificó en el historial inspeccionado una publicación específica de la cual se hubiera copiado. Se presenta como \textbf{modelo local inspirado en rozamiento}, sin adjudicarle una procedencia que el código no demuestra.

El coeficiente se expresa en comando normalizado antes de ganancias, no en newtons. Si se quisiera representar un rozamiento físico constante en ambos ejes, habría que convertir el esfuerzo lineal o torque deseado a la escala del motor y considerar la transmisión. La relación fija 1:2 de comandos no demuestra una relación 1:2 de fuerzas reales.

\subsection{Qué se puede pulir al implementarlo}
\subsubsection{Ganancia de menú y sensibilidad del ajuste}
\code{Friccion} ajusta $F_c$ de 0 a 1 en pasos de 0.01 y permite apagar el efecto manteniendo su valor configurado. Pasar de 0.03 a 0.04 aumenta el coeficiente un $33.3\%$: el paso del menú es relativamente grande para el default. Podría convenir ofrecer pasos más finos para este efecto; sería una modificación del menú, no una opción ya disponible.

Subir $F_c$ aumenta el roce en todo el rango de velocidades útiles, hasta que interviene un límite. Ajustar la ganancia general cambia también resorte y perturbaciones, por lo que no permite identificar por separado la causa de una sensación distinta. No existen coeficientes independientes de fricción para roll y pitch en el menú actual.

\subsubsection{Calibraciones que cambian la fricción aunque su ganancia sea igual}
Un cambio en los extremos normaliza de otra manera la misma velocidad física. Si se reduce a la mitad el recorrido físico que equivale a $x=1$, la misma velocidad física se representa con aproximadamente el doble de velocidad normalizada. La fricción aparecerá antes y el damper aumentará, aunque no se cambien $F_c$ ni $B$.

Con potenciómetros, revisar rango ADC, centro y ruido. Demasiada histéresis puede esconder movimientos pequeños; muy poca puede producir pequeñas inversiones de signo. Con encoders, medir cuentas reales por recorrido y observar el tiempo entre pasos lentos. Con SimConnect, verificar curvas y zonas muertas del eje y la cadencia real de paquetes; corregir sólo la ganancia no elimina el retardo de esa ruta.

\subsubsection{Parámetros de código que afectan la sensación de roce}
\begin{longtable}{p{0.26\linewidth}p{0.23\linewidth}p{0.42\linewidth}}
\toprule
Ajuste & Efecto principal & Criterio para evaluarlo\\
\midrule
$F_c$ & Escala la resistencia. & Comparar movimientos de velocidad similar, sin damper ni resorte.\\
$v_0=0.20$ & Cambia rapidez de aparición del roce. & Comparar inicio lento y movimiento más rápido; menor valor da transición más abrupta.\\
$v_d=0.02$ & Elimina respuesta a velocidad pequeña. & Medir ruido en reposo y movimiento real mínimo que se desea sentir.\\
$\tau=40$ ms & Suaviza velocidad, retrasa resistencia. & Observar inversiones y comienzo/final del movimiento.\\
Reposo de 40 ms & Fuerza extinción de la cola. & Revisar pasos lentos de encoder y tiempo efectivo entre muestras simuladas.\\
Límite conjunto 0.4 & Acota damper más fricción. & Separar saturación de rama de la saturación de todos los efectos.\\
Factor roll 0.5 & Escala ambos efectos de movimiento. & Comparar esfuerzos físicos, no sólo números PWM.\\
\bottomrule
\end{longtable}
Excepto las ganancias del menú, esos ajustes requieren edición y recompilación. Cambiar los parámetros comunes para mejorar fricción modifica también el damper, por lo que deben volver a evaluarse juntos.

\subsubsection{Ensayo propuesto para distinguir roce de amortiguación}
\begin{enumerate}
\item Registrar la configuración actual y comprobar el sensor. Aislar fricción, dejando centrado, damper y perturbaciones desactivados, con la ruta común habilitada.
\item Mover a varias velocidades sostenidas y en ambos sentidos. La resistencia calculada debe crecer al principio y luego variar poco; debe desaparecer al quedar quieto después de actualizar el estimador.
\item Comparar cruces del centro y movimientos lejos de él. A igual velocidad estimada, la fricción debe ser la misma; una diferencia puede indicar calibración, transmisión o rozamiento real distintos.
\item Ajustar $F_c$ para el nivel de roce y, sólo si es necesario, $v_0$ para la forma de aparición. Evitar corregir ruido únicamente aumentando la zona muerta, porque también se oculta movimiento lento real.
\item Incorporar damper y resorte. Evaluar esfuerzo al alejarse, suavidad al regresar y rebote; comprobar que no se está ajustando dentro de una saturación que oculta cambios.
\item Volver a habilitar perturbaciones y verificar que siguen siendo perceptibles. Registrar posición, velocidad estimada, comando y PWM junto con los valores guardados.
\end{enumerate}
Este ensayo separa las contribuciones del software, pero no mide por sí solo la fuerza transmitida a la mano. Con los defaults, la fricción aislada utiliza pocas cuentas PWM: la zona de funcionamiento real del motor y el rozamiento mecánico pueden dominar lo percibido.

\subsubsection{Mejoras posibles y qué datos hacen falta}
Si se desea reproducir adherencia estática, hace falta otro modelo con estado y una regla de inicio de deslizamiento; aumentar $F_c$ en la ecuación actual no lo agrega. Si se desea compensar rozamiento propio del mecanismo, debe distinguirse esa compensación de la fricción virtual que se quiere añadir: una resistencia programada se suma al rozamiento mecánico existente.

Otras mejoras posibles son coeficientes separados por eje, pasos finos de menú, estimación de velocidad mejor adaptada al sensor y medición de torque/fuerza frente a velocidad. Deben definirse con datos de la instalación; no hay un valor universal correcto de fricción para todos los yugos. Tampoco se ha demostrado que aumentar fricción mejore automáticamente la estabilidad de cualquier montaje.

\subsection{Mapa de archivos para revisar o modificar ambos efectos}
\begin{longtable}{p{0.38\linewidth}p{0.53\linewidth}}
\toprule
Archivo / función & Responsabilidad comprobada\\
\midrule
\code{SimConnectBridge/main.cpp} & Solicita YOKE X/Y, valida telemetría, codifica posiciones y controla envío/armado.\\
\code{SimTelemetry.h::Receiver} & Valida trama, decodifica posición y registra tiempo de recepción y disponibilidad.\\
\code{MiYoke_PD_V2.4.ino} & Selecciona fuente, normaliza encoders, actualiza movimiento local y envía comandos a motores.\\
\code{AxisPotentiometers.h} & Lectura ADC, descarte y promedio por canal.\\
\code{PotentiometerPosition.h} & Histéresis móvil y posición normalizada según extremos y centro.\\
\code{PotCalibration.h} & Representación y validación de calibración persistida de potenciómetros.\\
\code{AxisEncoders.h}, \code{PicoEncoder.h} & Origen relativo, dirección y lectura de cuadratura.\\
\code{YokeMotion.h} & Derivada, límite, filtro de 40 ms, inversión y reposo.\\
\code{SimEffects.h::motionAdjustedForce} & Fórmulas de damper/fricción y límite conjunto.\\
\code{SimEffects.h::Effects::command} & Ganancias de eje, suma con centrado y perturbaciones, límite final.\\
\code{GainSettings.h}, \code{mis_defines.h} & Defaults, habilitaciones, ajustes de menú y parámetros de compilación.\\
\code{MotorDriver.h} & Polaridad recibida, conversión y redondeo a PWM, techo y salidas del puente.\\
\code{CAMBIOS_REALISMO.md} y copias en \code{tesis/auditoria_c182rg} & Diferencia entre el comportamiento histórico y las cargas independientes actuales.\\
\code{tests/simconnect_checks.cpp} & Comprobaciones de fuentes de movimiento, simetría/oposición de cargas, protocolo y PWM.\\
\bottomrule
\end{longtable}
\textbf{Alcance del análisis:} se documentaron las ecuaciones implementadas, su procedencia verificable, las tres fuentes de posición y su relación con los demás efectos. Las propuestas de mejora están diferenciadas del comportamiento actual. No se modificó firmware, no se cargó la placa y no se realizó una identificación experimental del esfuerzo mecánico.
\end{document}
